\IfFileExists{stacks-project.cls}{%
\documentclass{stacks-project}
}{%
\documentclass{amsart}
}
\usepackage{amssymb}

% For dealing with references we use the comment environment
\usepackage{verbatim}
\usepackage{fancyvrb}
\usepackage{pmboxdraw}
\newenvironment{reference}{\comment}{\endcomment}
%\newenvironment{reference}{}{}
\newenvironment{slogan}{\comment}{\endcomment}
\newenvironment{history}{\comment}{\endcomment}

% For source figures retained as assets
\usepackage{graphicx}

% For commutative diagrams we use Xy-pic
\usepackage[all]{xy}

% We use 2cell for 2-commutative diagrams.
\xyoption{2cell}
\UseAllTwocells

% We use multicol for the list of chapters between chapters
\usepackage{multicol}

% This is generally recommended for better output
\usepackage{lmodern}
\usepackage[T1]{fontenc}
\usepackage{textalpha}

% For cross-file-references
% We prefer the snapshot tmp/xref/#2.aux,
% which latexmkrc replaces atomically, so
% that documents compiled simultaneously
% never read a half-written .aux file.
% The last argument keeps links to #2.pdf.
\usepackage{xr-hyper}
\newcommand{\maybeexternaldocument}[2]{%
  \IfFileExists{#2.aux}{%
    \externaldocument[#1]{#2}%
  }{%
    \PackageWarning{stack}{Missing #2.aux; external references with prefix #1 unavailable}%
  }%
}
\let\externaldocumentorig\externaldocument
\renewcommand{\externaldocument}[2][]{%
  \IfFileExists{tmp/xref/#2.aux}{%
    \externaldocumentorig[#1]{tmp/xref/#2}[#2.pdf]%
  }{%
  \IfFileExists{#2.aux}{%
    \externaldocumentorig[#1]{#2}%
  }{%
    \PackageWarning{stack}{Missing #2.aux; external references with prefix #1 unavailable}%
  }%
  }%
}

% Package for hypertext links:
\usepackage{hyperref}

% For any local file, say "hello.tex" you want to link to please
% use \externaldocument[hello-]{hello}
\externaldocument[introduction-]{introduction}
\externaldocument[conventions-]{conventions}
\externaldocument[sets-]{sets}
\externaldocument[categories-]{categories}
\externaldocument[topology-]{topology}
\externaldocument[sheaves-]{sheaves}
\externaldocument[sites-]{sites}
\externaldocument[stacks-]{stacks}
\externaldocument[fields-]{fields}
\externaldocument[algebra-]{algebra}
\externaldocument[brauer-]{brauer}
\externaldocument[homology-]{homology}
\externaldocument[derived-]{derived}
\externaldocument[simplicial-]{simplicial}
\externaldocument[more-algebra-]{more-algebra}
\externaldocument[smoothing-]{smoothing}
\externaldocument[modules-]{modules}
\externaldocument[sites-modules-]{sites-modules}
\externaldocument[injectives-]{injectives}
\externaldocument[cohomology-]{cohomology}
\externaldocument[sites-cohomology-]{sites-cohomology}
\externaldocument[dga-]{dga}
\externaldocument[dpa-]{dpa}
\externaldocument[sdga-]{sdga}
\externaldocument[hypercovering-]{hypercovering}
\externaldocument[schemes-]{schemes}
\externaldocument[constructions-]{constructions}
\externaldocument[properties-]{properties}
\externaldocument[morphisms-]{morphisms}
\externaldocument[coherent-]{coherent}
\externaldocument[divisors-]{divisors}
\externaldocument[limits-]{limits}
\externaldocument[varieties-]{varieties}
\externaldocument[topologies-]{topologies}
\externaldocument[descent-]{descent}
\externaldocument[perfect-]{perfect}
\externaldocument[more-morphisms-]{more-morphisms}
\externaldocument[flat-]{flat}
\externaldocument[groupoids-]{groupoids}
\externaldocument[more-groupoids-]{more-groupoids}
\externaldocument[etale-]{etale}
\externaldocument[chow-]{chow}
\externaldocument[intersection-]{intersection}
\externaldocument[pic-]{pic}
\externaldocument[weil-]{weil}
\externaldocument[adequate-]{adequate}
\externaldocument[dualizing-]{dualizing}
\externaldocument[duality-]{duality}
\externaldocument[discriminant-]{discriminant}
\externaldocument[derham-]{derham}
\externaldocument[local-cohomology-]{local-cohomology}
\externaldocument[algebraization-]{algebraization}
\externaldocument[curves-]{curves}
\externaldocument[resolve-]{resolve}
\externaldocument[models-]{models}
\externaldocument[functors-]{functors}
\externaldocument[equiv-]{equiv}
\externaldocument[pione-]{pione}
\externaldocument[etale-cohomology-]{etale-cohomology}
\externaldocument[proetale-]{proetale}
\externaldocument[relative-cycles-]{relative-cycles}
\externaldocument[more-etale-]{more-etale}
\externaldocument[trace-]{trace}
\externaldocument[crystalline-]{crystalline}
\externaldocument[spaces-]{spaces}
\externaldocument[spaces-properties-]{spaces-properties}
\externaldocument[spaces-morphisms-]{spaces-morphisms}
\externaldocument[decent-spaces-]{decent-spaces}
\externaldocument[spaces-cohomology-]{spaces-cohomology}
\externaldocument[spaces-limits-]{spaces-limits}
\externaldocument[spaces-divisors-]{spaces-divisors}
\externaldocument[spaces-over-fields-]{spaces-over-fields}
\externaldocument[spaces-topologies-]{spaces-topologies}
\externaldocument[spaces-descent-]{spaces-descent}
\externaldocument[spaces-perfect-]{spaces-perfect}
\externaldocument[spaces-more-morphisms-]{spaces-more-morphisms}
\externaldocument[spaces-flat-]{spaces-flat}
\externaldocument[spaces-groupoids-]{spaces-groupoids}
\externaldocument[spaces-more-groupoids-]{spaces-more-groupoids}
\externaldocument[bootstrap-]{bootstrap}
\externaldocument[spaces-pushouts-]{spaces-pushouts}
\externaldocument[spaces-chow-]{spaces-chow}
\externaldocument[groupoids-quotients-]{groupoids-quotients}
\externaldocument[spaces-more-cohomology-]{spaces-more-cohomology}
\externaldocument[spaces-simplicial-]{spaces-simplicial}
\externaldocument[spaces-duality-]{spaces-duality}
\externaldocument[formal-spaces-]{formal-spaces}
\externaldocument[restricted-]{restricted}
\externaldocument[spaces-resolve-]{spaces-resolve}
\externaldocument[formal-defos-]{formal-defos}
\externaldocument[defos-]{defos}
\externaldocument[cotangent-]{cotangent}
\externaldocument[examples-defos-]{examples-defos}
\externaldocument[algebraic-]{algebraic}
\externaldocument[examples-stacks-]{examples-stacks}
\externaldocument[stacks-sheaves-]{stacks-sheaves}
\externaldocument[criteria-]{criteria}
\externaldocument[artin-]{artin}
\externaldocument[quot-]{quot}
\externaldocument[stacks-properties-]{stacks-properties}
\externaldocument[stacks-morphisms-]{stacks-morphisms}
\externaldocument[stacks-limits-]{stacks-limits}
\externaldocument[stacks-cohomology-]{stacks-cohomology}
\externaldocument[stacks-perfect-]{stacks-perfect}
\externaldocument[stacks-introduction-]{stacks-introduction}
\externaldocument[stacks-more-morphisms-]{stacks-more-morphisms}
\externaldocument[stacks-geometry-]{stacks-geometry}
\externaldocument[moduli-]{moduli}
\externaldocument[moduli-curves-]{moduli-curves}
\externaldocument[examples-]{examples}
\externaldocument[exercises-]{exercises}
\externaldocument[guide-]{guide}
\externaldocument[desirables-]{desirables}
\externaldocument[coding-]{coding}
\externaldocument[obsolete-]{obsolete}
\externaldocument[fdl-]{fdl}
\externaldocument[index-]{index}

% Heap Project documents
\externaldocument[k3-]{k3}
\externaldocument[torelli-]{torelli}
\externaldocument[linear-algebra-]{linear-algebra}
\externaldocument[groups-]{groups}
\externaldocument[complex-analysis-]{complex-analysis}
\externaldocument[functional-analysis-]{functional-analysis}
\externaldocument[categories-]{categories}
\externaldocument[commutative-algebra-]{commutative-algebra}
\externaldocument[ringed-spaces-]{ringed-spaces}
\externaldocument[homological-algebra-]{homological-algebra}
\externaldocument[lie-algebras-]{lie-algebras}
\externaldocument[representation-theory-]{representation-theory}
\externaldocument[smooth-manifolds-]{smooth-manifolds}
\externaldocument[differential-topology-]{differential-topology}
\externaldocument[algebraic-topology-]{algebraic-topology}
\externaldocument[differential-geometry-]{differential-geometry}
\externaldocument[algebraic-geometry-]{algebraic-geometry}
\externaldocument[symplectic-geometry-]{symplectic-geometry}
\externaldocument[complex-geometry-]{complex-geometry}
\externaldocument[hodge-theory-]{hodge-theory}
\externaldocument[%
variations-of-hodge-structures-]{%
variations-of-hodge-structures}
\externaldocument[%
homological-mirror-symmetry-]{%
homological-mirror-symmetry}
\maybeexternaldocument{deformations-}{deformations}
\externaldocument[vertex-algebras-]{vertex-algebras}
\externaldocument[springer-]{springer}
\externaldocument[infty-categories-]{infty-categories}
\externaldocument[seminars-differential-geometry-]{%
seminars-differential-geometry}
\externaldocument[seminars-complex-geometry-]{%
seminars-complex-geometry}
\externaldocument[seminars-algebraic-geometry-]{%
seminars-algebraic-geometry}
\externaldocument[seminars-analysis-pde-]{%
seminars-analysis-pde}
\externaldocument[surfaces-]{surfaces}
\externaldocument[birational-maps-commutative-algebra-]{birational-maps-commutative-algebra}
\externaldocument[cremona-transformations-]{cremona-transformations}
\externaldocument[derived-categories-of-sheaves-]{derived-categories-of-sheaves}
\externaldocument[geometric-invariant-theory-]{geometric-invariant-theory}
\externaldocument[geometric-stability-conditions-and-group-actions-]{geometric-stability-conditions-and-group-actions}
\externaldocument[moduli-spaces-of-sheaves-]{moduli-spaces-of-sheaves}
\externaldocument[bridgeland-stability-general-theory-]{bridgeland-stability-general-theory}
\externaldocument[hilbert-schemes-]{hilbert-schemes}
\externaldocument[noncommutative-abelian-surfaces-and-kummer-type-hyperkahler-varieties-]{noncommutative-abelian-surfaces-and-kummer-type-hyperkahler-varieties}
\externaldocument[mumford-tate-groups-in-hodge-theory-]{mumford-tate-groups-in-hodge-theory}
\externaldocument[hodge-birational-atoms-2026-]{hodge-birational-atoms-2026}
\externaldocument[xxii-egd-2026-]{%
xxii-egd-2026}
\externaldocument[pde-]{pde}
\externaldocument[probability-]{probability}

\externaldocument[linguistics-]{linguistics}

\externaldocument[%
16th-alga-meeting-2026-]{%
16th-alga-meeting-2026}

\externaldocument[%
iv-jornadas-lefschetz-2026-]{%
iv-jornadas-lefschetz-2026}

\externaldocument[reading-the-masters-]{%
reading-the-masters}

\externaldocument[real-analysis-]{%
real-analysis}

\externaldocument[optimal-transport-]{%
optimal-transport}

\externaldocument[history-]{history}

\externaldocument[seiberg-witten-]{seiberg-witten}

\externaldocument[history-philosophy-biology-]{history-philosophy-biology}

\externaldocument[classical-mechanics-]{classical-mechanics}
\externaldocument[electromagnetism-]{electromagnetism}
\externaldocument[waves-and-optics-]{waves-and-optics}
\externaldocument[wif-]{waves-in-fluids-iv}
\externaldocument[thermal-physics-]{thermal-physics}
\externaldocument[quantum-mechanics-]{quantum-mechanics}
\externaldocument[special-relativity-]{special-relativity}
\externaldocument[physics-]{physics}

\externaldocument[chemistry-]{chemistry}

\externaldocument[molecular-cell-biology-]{molecular-cell-biology}

\externaldocument[microbiology-]{microbiology}

\externaldocument[protistology-]{protistology}

\externaldocument[mycology-]{mycology}

\externaldocument[botany-]{botany}

\externaldocument[zoology-]{zoology}

\externaldocument[animal-histology-]{animal-histology}

\externaldocument[animal-physiology-]{animal-physiology}

\externaldocument[systematics-biogeography-]{systematics-biogeography}

\externaldocument[ecology-natural-resources-]{ecology-natural-resources}

\externaldocument[earth-sciences-]{earth-sciences}

\externaldocument[biostatistics-]{biostatistics}

\externaldocument[mathematical-systems-biology-]{mathematical-systems-biology}

\externaldocument[cimpa-el-salvador-2024-]{cimpa-el-salvador-2024}

\externaldocument[anthropology-]{anthropology}

\externaldocument[yang-mills-theory-]{yang-mills-theory}

\externaldocument[seminars-algebra-]{seminars-algebra}

\externaldocument[seminars-physics-]{seminars-physics}

\externaldocument[seminars-computing-]{seminars-computing}

\externaldocument[oktobermat-2026-]{oktobermat-2026}

\externaldocument[saet-2026-]{saet-2026}

\externaldocument[seminars-symplectic-geometry-]{seminars-symplectic-geometry}

\externaldocument[seminars-dynamics-]{seminars-dynamics}

\externaldocument[oktobermat-2024-]{oktobermat-2024}

\externaldocument[quantum-groups-]{quantum-groups}

\externaldocument[legal-2027-]{legal-2027}

\externaldocument[derived-categories-]{derived-categories}

% Theorem environments.
%
\theoremstyle{plain}
\newtheorem{theorem}[subsection]{Theorem}
\newtheorem{proposition}[subsection]{Proposition}
\newtheorem{lemma}[subsection]{Lemma}
\newtheorem{corollary}[subsection]{Corollary}

\theoremstyle{definition}
\newtheorem{definition}[subsection]{Definition}
\newtheorem{example}[subsection]{Example}
\newtheorem{exercise}[subsection]{Exercise}
\newtheorem{situation}[subsection]{Situation}

\theoremstyle{remark}
\newtheorem{remark}[subsection]{Remark}
\newtheorem{remarks}[subsection]{Remarks}

\numberwithin{equation}{subsection}

% Macros
%
\def\lim{\mathop{\mathrm{lim}}\nolimits}
\def\colim{\mathop{\mathrm{colim}}\nolimits}
\def\Spec{\mathop{\mathrm{Spec}}}
\def\Hom{\mathop{\mathrm{Hom}}\nolimits}
\def\Ext{\mathop{\mathrm{Ext}}\nolimits}
\def\SheafHom{\mathop{\mathcal{H}\!\mathit{om}}\nolimits}
\def\SheafExt{\mathop{\mathcal{E}\!\mathit{xt}}\nolimits}
\def\Sch{\mathit{Sch}}
\def\Mor{\mathop{\mathrm{Mor}}\nolimits}
\def\Ob{\mathop{\mathrm{Ob}}\nolimits}
\def\Sh{\mathop{\mathit{Sh}}\nolimits}
\def\NL{\mathop{N\!L}\nolimits}
\def\CH{\mathop{\mathrm{CH}}\nolimits}
\def\proetale{{pro\text{-}\acute{e}tale}}
\def\etale{{\acute{e}tale}}
\def\QCoh{\mathit{QCoh}}
\def\Ker{\mathop{\mathrm{Ker}}}
\def\Im{\mathop{\mathrm{Im}}}
\def\Coker{\mathop{\mathrm{Coker}}}
\def\Coim{\mathop{\mathrm{Coim}}}

% Boxtimes
%
\DeclareMathSymbol{\boxtimes}{\mathbin}{AMSa}{"02}

%
% Macros for moduli stacks/spaces
%
\def\QCohstack{\mathcal{QC}\!\mathit{oh}}
\def\Cohstack{\mathcal{C}\!\mathit{oh}}
\def\Spacesstack{\mathcal{S}\!\mathit{paces}}
\def\Quotfunctor{\mathrm{Quot}}
\def\Hilbfunctor{\mathrm{Hilb}}
\def\Curvesstack{\mathcal{C}\!\mathit{urves}}
\def\Polarizedstack{\mathcal{P}\!\mathit{olarized}}
\def\Complexesstack{\mathcal{C}\!\mathit{omplexes}}
% \Pic is the operator that assigns to X its picard group, usage \Pic(X)
% \Picardstack_{X/B} denotes the Picard stack of X over B
% \Picardfunctor_{X/B} denotes the Picard functor of X over B
\def\Pic{\mathop{\mathrm{Pic}}\nolimits}
\def\Picardstack{\mathcal{P}\!\mathit{ic}}
\def\Picardfunctor{\mathrm{Pic}}
\def\Deformationcategory{\mathcal{D}\!\mathit{ef}}


\begin{document}

\title{Seminars in Algebraic Geometry}

\maketitle

\phantomsection

\label{section-phantom}

\tableofcontents



\section{Classification of log Calabi--Yau
pairs}
\label{section-alves-log-calabi-yau}
%% speaker: Eduardo Alves da Silva, Institut de
%% speaker: Mathématiques d'Orsay
%% place: Seminário das Sextas, PUC-Rio
%% date: September 6, 2024

\noindent
Eduardo Alves da Silva, Institut de
Mathématiques d'Orsay.
Seminário das Sextas, PUC-Rio.
September 6, 2024.

\medskip\noindent
No abstract available.

\medskip\noindent
{\bf Sources of these notes.} Daniel's
LaTeX notes typed during the talk,
transcribed here in October 2026. There are
no slides or recording. The date is the
file date of the original notes. Gaps in
the original notes are marked as such.

\medskip\noindent
{\bf Introduction.}

\begin{definition}
\label{definition-alves-log-cy-pair}
A \emph{log Calabi--Yau pair} is a log
canonical pair $(X,D)$ consisting of a
normal projective variety $X$ and a reduced
Weil divisor $D$ such that
$K_X+D\sim_{\mathbb{Z}}0$.
\end{definition}

\begin{remark}
\label{remark-alves-volume-form}
Let $n=\dim X$. If $(X,D)$ is a log
Calabi--Yau pair, there is a rational
$n$-form $\omega=\omega_{X,D}$ on $X$,
unique up to a nonzero scalar, such that
$$
\operatorname{div}(\omega)+D=0.
$$
It is called the \emph{volume form} of the
pair.
\end{remark}

A log Calabi--Yau pair sits between two
outputs of the minimal model program. On one
side, $(X,D)$ is already a minimal model for
the log MMP, since $K_X+D$ is nef. On the
other side, running the classical MMP on
$X$ ends with a Mori fibre space, by the
chain of implications
$$
K_X+D\sim 0
\implies
-K_X=D\geq 0,
$$
so $K_X$ is not pseudo-effective; hence $X$
is uniruled by the theorem of
Boucksom--Demailly--Păun--Peternell (BDPP),
so its Kodaira dimension is $-\infty$ and
the MMP on $X$ ends with a Mori fibre space.
The original notes leave the example that
followed empty.

\begin{definition}
\label{definition-alves-volume-preserving}
Let
$f\colon (X,D_X)\dashrightarrow (Y,D_Y)$ be
a birational map of log Calabi--Yau pairs.
Then $f$ is \emph{volume preserving} if
$$
f^*\omega_{Y,D_Y}=\lambda\,\omega_{X,D_X}
$$
for some $\lambda\in\mathbb{C}^*$.
\end{definition}

The original notes drop the right-hand side
of this equation; the form above is the
standard one.

\begin{remark}
\label{remark-alves-bir-vp}
The volume preserving birational
self-maps form a subgroup
$$
\operatorname{Bir}^{\operatorname{vp}}(X,D_X)
\subset \operatorname{Bir}(X)
$$
of the group of all birational
self-maps of $X$.
\end{remark}

Equivalent definitions of volume preserving
maps:
\begin{enumerate}
\item $f$ \emph{preserves discrepancies}:
for every divisor $E$ over $X$ and $Y$,
$$
a(E,X,D_X)=a(E,Y,D_Y).
$$
\item $f$ admits a log resolution
$$
\xymatrix{
& (Z,D_Z) \ar[ld]_p \ar[rd]^q & \\
(X,D_X) \ar@{-->}[rr]^f & & (Y,D_Y)
}
$$
with
$$
p^*(K_X+D_X)=q^*(K_Y+D_Y).
$$
Warning: $D_Z$ need not be effective; the
speaker referred to Section 5.4 of his PhD
thesis for an example.
\end{enumerate}
Volume preserving maps are also called
\emph{crepant} birational maps, and one
writes $(X,D_X)\cong_{\operatorname{vp}}
(Y,D_Y)$ or
$(X,D_X)\cong_{\operatorname{cbir}}(Y,D_Y)$.

\medskip\noindent
{\bf Problem} (very hard). Classify log
Calabi--Yau pairs up to volume preserving
equivalence.

The most important invariant for this
problem is the following.

\begin{definition}
\label{definition-alves-coregularity}
The \emph{coregularity}
$\operatorname{coreg}(X,D_X)$ of a log
Calabi--Yau pair is the dimension of a
minimal log canonical centre of a dlt
modification
$$
(X^{\operatorname{dlt}},
D_{X^{\operatorname{dlt}}})
\to (X,D_X).
$$
\end{definition}

\begin{remark}
\label{remark-alves-coregularity-range}
Write $c=\operatorname{coreg}(X,D_X)$. Then
$0\leq c\leq \dim X$, and $c=\dim X$ if and
only if $X$ is Calabi--Yau and $D_X=0$.
\end{remark}

\medskip\noindent
{\bf Dimension 2.} After a minimal
resolution of singularities, a surface log
Calabi--Yau pair $(X,D_X)$ is one of the
following.
\begin{itemize}
\item $c=2$: $X$ is an abelian surface or a
K3 surface, and $D_X=0$.
\item $c=1$: either $X$ is rational and
$D_X\in |-K_X|$ is a smooth elliptic curve,
or $\pi\colon X\to E$ is a ruled surface
(not necessarily minimal) over a smooth
elliptic curve $E$ and $D_X=D_1+D_2\in
|-K_X|$ is the sum of two disjoint sections
of $\pi$.
\item $c=0$: $X$ is rational and $D_X$ is a
(possibly reducible) nodal curve of
arithmetic genus $1$.
\end{itemize}

\begin{example}
\label{example-alves-plane}
Take $X=\mathbb{P}^2$ and $D$ three lines, a
conic plus a line, a nodal cubic, or a
smooth cubic. The coregularity is $0$ in
the first three cases and $1$ for the
smooth cubic.
\end{example}

\begin{definition}
\label{definition-alves-toric-model}
A log Calabi--Yau pair $(X,D_X)$ has a
\emph{toric model} if
$(X,D_X)\cong_{\operatorname{vp}}(T,D_T)$
for a toric variety $T$, where $D_T$ is the
reduced sum of the torus invariant
divisors.
\end{definition}

\begin{theorem}[Gross--Hacking--Keel]
\label{theorem-alves-ghk}
Every surface log Calabi--Yau pair of
coregularity $0$ has a toric model.
\end{theorem}

This is false in dimension at least $3$.
Not checked.

\medskip\noindent
{\bf Dimension 3.}

\begin{theorem}[Ducat, 2023]
\label{theorem-alves-ducat}
Let $(\mathbb{P}^3,D)$ be a log Calabi--Yau
pair of coregularity $c\leq 1$. Then there
is a volume preserving birational map
$$
\varphi\colon (\mathbb{P}^3,D)
\dashrightarrow
(\mathbb{P}^1\times\mathbb{P}^2,D'),
$$
where
$$
D'=\{0\}\times\mathbb{P}^2
+\mathbb{P}^1\times E
+\{\infty\}\times\mathbb{P}^2
$$
lies in
$|-K_{\mathbb{P}^1\times\mathbb{P}^2}|$, for
a plane cubic $E\subset\mathbb{P}^2$ with
coordinates $(x:y:z)$, such that:
\begin{enumerate}
\item $c=1$ if and only if $E$ is smooth;
\item if $c=0$ then $E=\{xyz=0\}$; in
particular $D'$ is the toric boundary of
$\mathbb{P}^1\times\mathbb{P}^2$, so
$(\mathbb{P}^3,D)$ has a toric model.
\end{enumerate}
\end{theorem}

The original notes write
$\mathbb{P}^1\times\mathbb{P}^1$ and
$\mathbb{P}^4\times\mathbb{P}^2$ in this
statement; the formula for $D'$ only makes
sense on $\mathbb{P}^1\times\mathbb{P}^2$,
which is what is written here. Not checked
against Ducat's paper.

The missing case is $c=2$. Fact: $c=2$ if
and only if $D$ is an irreducible normal
quartic surface with canonical
singularities, that is, $D$ is smooth or
has ADE singularities, if and only if the
pair is canonical.

\begin{example}[Oguiso]
\label{example-alves-oguiso}
Oguiso constructed two smooth quartic
surfaces $D,D'\subset\mathbb{P}^3$,
isomorphic as abstract varieties, with a
$\varphi\in\operatorname{Bir}(\mathbb{P}^3)$
mapping $D$ birationally onto $D'$, and
such that
$(\mathbb{P}^3,D)
\not\cong_{\operatorname{vp}}
(\mathbb{P}^3,D')$.
\end{example}

In terms of coarse moduli spaces there is a
natural map
$$
\mathcal{M}^{c=2}_{(\mathbb{P}^3,D)}
\longrightarrow
\mathcal{M}^{\operatorname{can}}_{K3},
\quad
[(\mathbb{P}^3,D)]\longmapsto [D],
$$
and Oguiso's example shows that it is not
injective.

\medskip\noindent
{\bf Conjecture (trichotomy).} The speaker
compared the three values of the
coregularity of $(\mathbb{P}^3,D)$ with
curves of genus $0$, $1$ and at least $2$.
\begin{itemize}
\item Coregularity $0$, like genus $0$: the
moduli space has dimension $0$ and the group
$\operatorname{Bir}^{\operatorname{vp}}$ is
huge (``monstrous''), like
$\operatorname{PGL}(2,\mathbb{C})$.
\item Coregularity $1$, like genus $1$: the
moduli space has dimension $1$; the
original notes leave
$\operatorname{Bir}^{\operatorname{vp}}$
open, to be compared with the automorphism
group $C\rtimes\mathbb{Z}/d$ of an elliptic
curve $C$.
\item Coregularity $2$, like genus at least
$2$: the original notes leave the dimension
open, and
$\operatorname{Bir}^{\operatorname{vp}}$
should be $\operatorname{Dec}(D)$, to be
compared with the finite group
$\operatorname{Aut}(C)$, of order at most
$84(g-1)$, of a curve of genus $g$, whose
moduli space has dimension $3g-3$.
\end{itemize}
The original notes put this in a table whose
first row reads $0$, $0$, ?; the row should
read $0$, $1$, $2$, as written here.

For very general $D$, the surface $D$ is
smooth, and a volume preserving birational
map
$\varphi\colon (\mathbb{P}^3,D)
\dashrightarrow (\mathbb{P}^3,D)$
restricts to an isomorphism
$\varphi|_D\colon D\to D$. For a projective
variety $X$ and an irreducible subvariety
$Y\subset X$, the \emph{decomposition group}
$\operatorname{Bir}(Y,X)$ is the group of
those $f\in\operatorname{Bir}(X)$ whose
restriction
$f|_Y\colon Y\dashrightarrow Y$ is
birational. The notation
$\operatorname{Dec}(D)$ above presumably
means this group for
$Y=D\subset X=\mathbb{P}^3$.

\medskip\noindent
{\bf Conjecture (Shokurov).} Every rational
log Calabi--Yau threefold pair of
coregularity $0$ has a toric model.

By Theorem~\ref{theorem-alves-ducat} this
holds for $X=\mathbb{P}^3$.

\begin{definition}
\label{definition-alves-complexity}
Let $(X,D_X)$ be a log Calabi--Yau pair with
$D_X=D_1+\cdots+D_r$. Its \emph{complexity}
is the nonnegative number
$$
c(X,D_X)=\dim X
+\operatorname{rk}
\operatorname{Cl}(X)_{\mathbb{Q}}
-r.
$$
\end{definition}

Fact (Brown--McKernan--Svaldi--Zong, 2018):
if $c(X,D_X)=0$ then $(X,D_X)$ has a toric
model.

\begin{definition}
\label{definition-alves-bir-complexity}
The \emph{birational complexity} of a log
Calabi--Yau pair is
$$
c_{\operatorname{bir}}(X,D_X)
=\min c(Y,D_Y),
$$
the minimum over all pairs $(Y,D_Y)$ with
$(Y,D_Y)\cong_{\operatorname{vp}}(X,D_X)$.
\end{definition}

\begin{theorem}[Mauri--Moraga, 2023]
\label{theorem-alves-mauri-moraga}
A log Calabi--Yau pair $(X,D_X)$ satisfies
$c_{\operatorname{bir}}(X,D_X)=0$ if and only
if it has a toric model.
\end{theorem}

\begin{definition}
\label{definition-alves-cluster-type}
A log Calabi--Yau pair $(X,D_X)$ of
dimension $n$ is of \emph{cluster type} if
there is a volume preserving birational map
$$
\varphi\colon
(\mathbb{P}^n,H_0+\cdots+H_n)
\dashrightarrow (X,D_X),
$$
where the $H_i$ are the coordinate
hyperplanes, such that the exceptional
locus of $\varphi$ meets the torus
$\mathbb{G}_m^n$ in codimension at least
$2$; equivalently, the torus embeds into
$X\setminus D_X$.
\end{definition}

The original notes write $H_0+\cdots+H_{n+1}$
and are garbled at the end of this
definition; the version above is my reading,
not checked.

\begin{theorem}[Alves da
Silva--Figueroa--Moraga, 2024]
\label{theorem-alves-cluster-p3}
Let $(\mathbb{P}^3,D)$ be a log Calabi--Yau
pair of coregularity $0$, with $D$ general
in its deformation class. Then
$(\mathbb{P}^3,D)$ is of cluster type unless
one of the following holds:
\begin{enumerate}
\item $D$ is reducible, $D=H+C$ with $H$ a
plane and $C$ a cubic surface, such that
$H\cap C$ is a nodal plane cubic;
\item $D$ is irreducible and has double
points along a line.
\end{enumerate}
\end{theorem}

\medskip\noindent
{\bf Sketch.} The pair $(\mathbb{P}^3,D)$ is
of cluster type if and only if it is of
cluster type over $\mathbb{P}^1$, if and only
if some dlt modification $(X,D_X)$ of it
admits a crepant contraction onto
$\mathbb{P}^1$ with boundary
$\{0\}+\{\infty\}$. The original notes end
here, and write $\mathbb{P}^4$ for the
target; this reading is not checked.



\section{Blown-up toric surfaces with
non-polyhedral effective cone}
\label{section-tevelev-blown-up-toric}
%% speaker: Jenia Tevelev
%% place: Geometric Structures Seminar, IMPA
%% date: September 19, 2024

\noindent
Jenia Tevelev.
Geometric Structures Seminar, IMPA.
September 19, 2024.

\medskip\noindent
No abstract available. The talk is based on
joint work with Castravet, Laface and
Ugaglia: A.-M. Castravet, A. Laface,
J. Tevelev, L. Ugaglia, \emph{Blown-up toric
surfaces with non-polyhedral effective
cone}, J. reine angew. Math. 800 (2023).

\medskip\noindent
{\bf Sources of these notes.} Daniel's
LaTeX notes typed during the talk,
transcribed here in October 2026. There are
no slides or recording; the drawings made
during the talk are described in prose. The
original notes give the speaker only as
``Jenia'' and the journal as ``Crelle 23'';
the full names and reference above are mine,
not checked. Gaps in the original notes are
marked as such.

\medskip\noindent
{\bf Motivation.} Understand the birational
geometry of $\overline{M}_{0,n}$, the moduli
space of stable rational curves with $n$
marked points. The open part
$$
M_{0,n}=
\{(p_1,\ldots,p_n)\in(\mathbb{P}^1)^n :
p_i\neq p_j\}/\operatorname{PGL}(2)
$$
is a smooth affine variety of dimension
$n-3$, and $\overline{M}_{0,n}\supset
M_{0,n}$ is a smooth projective variety. The
original notes write ``dimension $n=3$'' and
draw the configuration of points on
$\mathbb{P}^1$ instead of writing it.

\medskip\noindent
{\bf Universal families.} There are
universal families
$U_n\to\overline{M}_{0,n}$ and, over the
open part, the trivial family of
$\mathbb{P}^1$'s with $n$ sections. The
boundary $\partial\overline{M}_{0,n}$ is a
simple normal crossings divisor. A stable
rational curve is a tree of $\mathbb{P}^1$'s
with $n$ distinct smooth marked points and
nodal singularities, locally
$\{xy=0\}\subset\mathbb{C}^2$, such that
every component has at least $3$ special
points (marked points or nodes).

\begin{lemma}
\label{lemma-tevelev-universal-family}
$U_n\cong\overline{M}_{0,n+1}$, and the map
$\overline{M}_{0,n+1}\to\overline{M}_{0,n}$
forgets the last point and stabilises.
\end{lemma}

\begin{example}
\label{example-tevelev-m04}
$M_{0,4}=\{(0,1,\infty,x)\}=
\mathbb{P}^1\setminus\{0,1,\infty\}$, with
the unique compactification
$\overline{M}_{0,4}=\mathbb{P}^1$; the
fibres of the universal family over
$0,1,\infty$ were drawn. This
$\mathbb{P}^1$ is the pencil of conics
through four general points
$p_1,\ldots,p_4\in\mathbb{P}^2$, so the
universal family is
$$
U_4=\operatorname{Bl}_4\mathbb{P}^2
=\overline{M}_{0,5}
\longrightarrow\overline{M}_{0,4},
$$
a del Pezzo surface of degree $5$.
\end{example}

The original notes write $U_q$, map it to
$\overline{M}_{0,5}$, and leave the degree
of the del Pezzo surface as a question mark;
the corrected version follows from
Lemma~\ref{lemma-tevelev-universal-family}.

\medskip\noindent
{\bf Review: divisors and cones.} Let $X$
be a normal $\mathbb{Q}$-factorial
projective variety. Then
$\operatorname{Pic}(X)\subset
\operatorname{Cl}(X)$, Cartier divisors
inside Weil divisors, and
\emph{$\mathbb{Q}$-factorial} means that for
every $D\in\operatorname{Cl}(X)$ some
multiple $mD$, $m>0$, lies in
$\operatorname{Pic}(X)$. So for an integral
curve $C\subset X$ the intersection number
$D\cdot C\in\mathbb{Q}$ makes sense.

The Néron--Severi spaces are
$N^1(X)$, real divisors modulo numerical
equivalence ($D\equiv 0$ if $D\cdot C=0$ for
all curves $C$), and $N_1(X)$, real
$1$-cycles modulo numerical equivalence
($C\equiv 0$ if $D\cdot C=0$ for all $D$).
They are finite dimensional and dual to each
other under the intersection pairing.
\begin{itemize}
\item The \emph{pseudo-effective cone}
$\operatorname{Eff}(X)\subset N^1(X)$ is the
closure of the cone spanned by classes of
effective divisors.
\item The \emph{Mori cone}
$\operatorname{NE}(X)\subset N_1(X)$ is the
closure of the cone spanned by classes of
effective curves.
\item The \emph{nef cone} is the dual
$$
\operatorname{Nef}(X)
=\operatorname{NE}(X)^\vee
=\{D : D\cdot C\geq 0
\ \forall C\in\operatorname{NE}(X)\}.
$$
\end{itemize}

\medskip\noindent
{\bf Linear systems.} For
$D\in\operatorname{Cl}(X)$,
$$
|D|=\mathbb{P}H^0(X,\mathcal{O}_X(D))
=\{D'\geq 0 : D'\sim D\},
$$
where $\mathcal{O}_X(D)$ is the divisorial
sheaf; it is nonempty if and only if $D$ is
linearly equivalent to an effective
divisor. There is a rational map
$$
\varphi_D\colon X\dashrightarrow |D|^\vee,
\quad
x\longmapsto\{D'\in|D| : x\in D'\},
$$
with base locus
$\operatorname{Bs}|D|=\bigcap_{D'\in|D|}D'$.
Write $D=\operatorname{Fix}(D)+M$, the
divisorial part of the base locus plus the
mobile part; then $\varphi_D=\varphi_M$. If
$\operatorname{Bs}|D|=\emptyset$, $D$ is
\emph{free} (globally generated), and then
$D$ is the pullback of a hyperplane under
$\varphi_D\colon X\to\mathbb{P}^r=|D|^\vee$:
$\mathcal{O}(D)=\varphi_D^*\mathcal{O}(1)$,
so $D$ is Cartier.

The \emph{stable base locus} is
$$
\mathbb{B}(D)=\bigcap_{m>0}
\operatorname{Bs}|mD|,
$$
and $D$ is \emph{semiample} if
$\mathbb{B}(D)=\emptyset$, that is, if $mD$
is free for some $m>0$. Semiample divisors
are nef.

\begin{theorem}[semiample fibration]
\label{theorem-tevelev-semiample-fibration}
If $D$ is semiample, the morphism
$\varphi_{|mD|}\colon X\to Y$ does not
depend on $m$ for $m\gg 0$ and divisible; it
has connected fibres and $Y$ is normal (not
necessarily $\mathbb{Q}$-factorial).
\end{theorem}

\begin{theorem}[Zariski]
\label{theorem-tevelev-zariski}
If $D\in\operatorname{Pic}(X)$ and
$\operatorname{Bs}|D|$ is finite, then $D$
is semiample.
\end{theorem}

\begin{corollary}
\label{corollary-tevelev-surfaces}
If $\dim X=2$ and $D\in\operatorname{Pic}(X)$
is effective with $\operatorname{Fix}(D)=0$,
then $D$ is semiample.
\end{corollary}

\begin{itemize}
\item $D$ is \emph{very ample} if
$\varphi_{|D|}$ is a closed embedding.
\item $D\in\operatorname{Cl}(X)$ is
\emph{ample} if $mD$ is very ample for some
$m>0$.
\item $D\in\operatorname{Cl}(X)$ is
\emph{big} if
$\dim\varphi_{|mD|}(X)=\dim X$ for some
$m>0$; equivalently $h^0(X,mD)$ grows like
$m^{\dim X}$ for $m\gg 0$ and divisible;
equivalently, by Iitaka, $\varphi_{|mD|}$ is
birational.
\end{itemize}

\medskip\noindent
{\bf Kleiman's criterion.} Although
ampleness and bigness are defined with
linear systems, they are numerical
properties.

\begin{theorem}[Kleiman]
\label{theorem-tevelev-kleiman}
$D\in\operatorname{Cl}(X)$ is ample if and
only if $D$ lies in the interior of
$\operatorname{Nef}(X)$. In particular
$\operatorname{Nef}(X)\subset
\operatorname{Eff}(X)$.
\end{theorem}

Similarly, $D$ lies in the interior of
$\operatorname{Eff}(X)$ if and only if $D$
is big, if and only if $mD$ is effective
plus ample for some $m>0$. For this one uses
the following lemma.

\begin{lemma}[Kodaira's lemma]
\label{lemma-tevelev-kodaira}
If $D$ is big and $A$ is effective, then
$mD-A$ is effective for some $m>0$.
\end{lemma}

The argument sketched was the exact
sequence
$$
0\to\mathcal{O}(mD-A)\to\mathcal{O}(mD)
\to\mathcal{O}(mD)|_A\to 0:
$$
$h^0(\mathcal{O}(mD))$ grows like
$m^{\dim X}$ while
$h^0(\mathcal{O}(mD)|_A)$ grows like
$m^{\dim X-1}$, so
$H^0(\mathcal{O}(mD-A))\neq 0$ for $m$
large. (The growth comparison is my
completion of the sketch.)

This was the end of the review.

\medskip\noindent
{\bf Questions asked by Fulton.} The
boundary $\partial\overline{M}_{0,n}$ gives
a stratification of $\overline{M}_{0,n}$, and
its one-dimensional strata are the
\emph{F-curves}.

\begin{example}
\label{example-tevelev-m05}
For $\overline{M}_{0,5}=\operatorname{Bl}_4
\mathbb{P}^2$ there are
$\binom{5}{2}=10$ boundary curves, the
$(-1)$-curves, and they generate the
effective cone.
\end{example}

\medskip\noindent
{\bf Conjecture (Fulton's F-conjecture).}
$\operatorname{NE}(\overline{M}_{0,n})$ is
generated by F-curves. Still open.

\medskip\noindent
{\bf Fulton's question (false).}
$\operatorname{Eff}(\overline{M}_{0,n})$ is
generated by boundary divisors. There are
many other extremal rays
(Castravet--Tevelev, 2013), and the cone is
not even rational polyhedral (the 2023
paper above).

The original notes write ``Fulton-Fuber''
and ``Gutrvet-Jenia''; the names above are my
reading, not checked.

\medskip\noindent
{\bf Birational contractions.} Let $X$ be a
normal $\mathbb{Q}$-factorial projective
variety and $f\colon X\dashrightarrow Y$ a
birational map to a $\mathbb{Q}$-factorial
variety, regular in codimension $1$. A
divisor $D\subset X$ is
\emph{$f$-exceptional} if
$\dim f(D)<\dim D$, equivalently if $f$ is
not an isomorphism at the generic point of
$D$. The map $f$ (or $Y$) is a
\emph{birational contraction} if there are
no $f^{-1}$-exceptional divisors.

Question from Misha (Verbitsky) in the
audience: is a contraction always regular?
Answer: no.

The map $f$ is a \emph{small
$\mathbb{Q}$-factorial modification} if
there are no $f$- or $f^{-1}$-exceptional
divisors, that is, $f$ is an isomorphism in
codimension $1$. A rational map
$f\colon X\dashrightarrow Y$ of
$\mathbb{Q}$-factorial varieties is a
\emph{rational contraction} if it is a
composition of birational contractions and
morphisms with connected fibres.

\medskip\noindent
{\bf Principle.} The birational geometry of
$X$ is the study of the rational
contractions of $X$.

Application to effective cones: if
$f\colon X\dashrightarrow Y$ is a birational
contraction, its exceptional divisors span
extremal rays of $\operatorname{Eff}(X)$.

\medskip\noindent
{\bf Hassett spaces.} These are the simplest
examples of birational contractions. (The
speaker restricted to genus $0$, but most of
this extends.) Choose rational weights
$a=(a_1,\ldots,a_n)$, $a_i>0$, with
$\sum_i a_i>2$. There is a birational
contraction
$$
\overline{M}_{0,n}\longrightarrow
\overline{M}_{0,n;a},
$$
the moduli space of $a$-stable rational
curves, restricting to the identity on
$M_{0,n}$. A curve $(C,p_1,\ldots,p_n)$ is
$a$-stable if
\begin{itemize}
\item $C$ is a nodal (semi log canonical)
curve of genus $0$ and the $p_i$ are not at
the nodes;
\item $\omega_C(\sum_i a_ip_i)$ is ample,
that is, on every component $R$, the number
of nodes on $R$ plus $\sum_{p_i\in R}a_i$ is
greater than $2$;
\item points may coincide: if $p_i=p_j$ for
all $i,j\in I\subset\{1,\ldots,n\}$, then
$\sum_{i\in I}a_i\leq 1$.
\end{itemize}

\begin{example}[Kapranov's map]
\label{example-tevelev-kapranov}
Take the weights $(1,x,\ldots,x)$ with
$1+(n-1)x=2+\varepsilon$, $\varepsilon\ll 1$,
so $p_1$ is a heavy point. A stable curve is
then $\mathbb{P}^1$ with $p_i\neq p_1$ for
$i>1$, and not all $p_i$, $i>1$, equal.
Sending $p_1\to\infty$ and $p_2\to 0$, the
remaining points $p_3,\ldots,p_n\in
\mathbb{C}$ are not all zero, so
$\overline{M}_{0,n;a}=\mathbb{P}^{n-3}$. The
map $\overline{M}_{0,n}\to\mathbb{P}^{n-3}$
is the \emph{Kapranov map}. A drawing showed
a curve with two components: the image keeps
the component carrying $p_1$ and contracts
the other one to a point. So the boundary
divisor $\Delta_I$ is contracted as soon as
$|I^c|\geq 3$.
\end{example}

\begin{exercise}
\label{exercise-tevelev-hypertree}
Take $n-2$ triangles $\Gamma_i$, each a set
of three of $n$ vertices (as for the black
triangles of a black-and-white triangulation
of a sphere with $n$ vertices).
\begin{enumerate}
\item Show the \emph{hypertree condition}
$$
\Big|\bigcup_{i\in I}\Gamma_i\Big|
\geq |I|+2
$$
for every set $I$ of triangles, with
equality if $|I|=1$ or $|I|=n-2$.
\item The inequality is strict for
$1<|I|<n-2$ unless the triangulation is a
connected sum. This is the
\emph{irreducible hypertree condition}.
\end{enumerate}
\end{exercise}

An \emph{irreducible hypertree} is a
collection $\Gamma_1,\ldots,\Gamma_{n-2}$
satisfying these conditions. It gives a
\emph{hypertree curve} $\mathcal{C}_\Gamma$:
one line for each triangle, two lines
meeting when the triangles share a vertex,
as for $\Gamma_i=\{a,b,c\}$ and
$\Gamma_j=\{a,x,y\}$. Unless the
triangulation is an octahedron, some vertex
lies on more than two black triangles.

Daniel's question in the original notes: is
this curve the union of the two lines
through $\{a,b,c\}$ and $\{a,x,y\}$? The
notes continue: the curve is locally the
union of coordinate axes.

Then
$$
M_{0,n}\subset
\operatorname{Mor}(\mathcal{C}_\Gamma,
\mathbb{P}^1),
$$
the maps that are linear on every component
of $\mathcal{C}_\Gamma$. A drawing showed
four lines with numbered intersection
points mapping to $\mathbb{P}^1$; the
intersection points go to six points of
$\mathbb{P}^1$. Composing with
$\varphi\mapsto\varphi^*\mathcal{O}(1)$
gives
$$
M_{0,n}\subset
\operatorname{Mor}(\mathcal{C}_\Gamma,
\mathbb{P}^1)\to
\operatorname{Pic}(\mathcal{C}_\Gamma),
$$
that is, a map
$M_{0,n}\to(\mathbb{C}^*)^{n-3}$, which is
birational by Riemann--Roch. Its exceptional
locus is a divisor $D_\Gamma$. A general
point of $D_\Gamma$ is obtained by mapping
$\mathcal{C}_\Gamma\to\mathbb{P}^2$,
choosing a point $x\in\mathbb{P}^2$ and
projecting from $x$, which gives $n$ points
on $\mathbb{P}^1$.

\begin{theorem}[Castravet--Tevelev]
\label{theorem-tevelev-hypertree-divisor}
If $\Gamma$ is an irreducible hypertree,
then $D_\Gamma\subset M_{0,n}$ is an
irreducible divisor, and its closure
$\overline{D}_\Gamma\subset
\overline{M}_{0,n}$ is exceptional for the
birational contraction from
$\overline{M}_{0,n}$ to a compactified
Jacobian of $\mathcal{C}_\Gamma$ extending
$M_{0,n}\to(\mathbb{C}^*)^{n-3}$. Hence
$\overline{D}_\Gamma$ spans an extremal ray
of $\operatorname{Eff}(\overline{M}_{0,n})$.
\end{theorem}

\begin{theorem}[Castravet--Laface--
Tevelev--Ugaglia]
\label{theorem-tevelev-not-polyhedral}
$\operatorname{Eff}(\overline{M}_{0,n})$ is
not finitely generated for $n\geq 10$.
\end{theorem}

\begin{lemma}
\label{lemma-tevelev-rational-contraction}
Let $f\colon X\dashrightarrow Y$ be a
rational contraction. If
$\operatorname{Eff}(X)$ is finitely
generated, so is $\operatorname{Eff}(Y)$.
Other properties preserved by rational
contractions: being a Mori dream space,
having a small $\mathbb{Q}$-factorial
modification with a nef but not semiample
divisor, and so on.
\end{lemma}

\begin{proof}
If $f$ is a birational contraction,
$\operatorname{Eff}(Y)=
f_*\operatorname{Eff}(X)$ for
$f_*\colon N^1(X)\to N^1(Y)$. If $f$ is a
morphism there is no pushforward of
divisors, but there is one of curves,
$f_*\colon N_1(X)\to N_1(Y)$. Let
$\operatorname{Mov}_1(X)$ be the cone
spanned by \emph{movable} curves, those
moving in a family that covers $X$. By the
theorem of
Boucksom--Demailly--Păun--Peternell,
$\operatorname{Mov}_1(X)$ is the dual of
$\operatorname{Eff}(X)$. If
$\operatorname{Eff}(X)$ is polyhedral, so is
$\operatorname{Mov}_1(X)$; the map
$f_*\colon\operatorname{Mov}_1(X)\to
\operatorname{Mov}_1(Y)$ is surjective, so
$\operatorname{Mov}_1(Y)$ is polyhedral, and
by duality so is $\operatorname{Eff}(Y)$.
\end{proof}

The original notes state BDPP as
``$\operatorname{Eff}(X)=
\operatorname{Mov}_1(X)$'' and attribute it
to ``Bookson Demaria PP''; the correct
statement is the duality used above.

\medskip\noindent
{\bf Goal.} Find a rational contraction of
$\overline{M}_{0,n}$ with a non-polyhedral
effective cone.

\medskip\noindent
{\bf Step 1.} Take the weights
$(1,1,\varepsilon,\ldots,\varepsilon)$, with
the heavy points at $0$ and $\infty$ and
$\varepsilon\ll 1$. The stable curves are
chains of $\mathbb{P}^1$'s with $0$ on the
leftmost component and $\infty$ on the
rightmost. The Hassett space is the
\emph{Losev--Manin space} $LM$, the toric
variety of the permutohedron, the convex
hull of
$\{(\sigma(1),\ldots,\sigma(k))\}_{\sigma\in
S_k}\subset\mathbb{R}^k$; for $k=3$ it is a
hexagon.

Universality (a lemma with Ana-Maria
Castravet): every projective toric variety
$\mathbb{P}(\Delta)$ is a rational
contraction of the permutohedral toric
variety, and therefore of
$\overline{M}_{0,n}$ for $n\gg 0$. This does
not help directly, since
$\operatorname{Eff}(\mathbb{P}(\Delta))$ is
polyhedral, generated by the toric boundary
divisors.

\medskip\noindent
{\bf Step 2.} Take the weights
$(1,1,x,\ldots,x)$, with $k$ copies of $x$
and $kx=1+\varepsilon$, $\varepsilon\ll 1$.
Then $\overline{M}_{0,n;a}=
\operatorname{Bl}_e LM$, the blow-up of the
Losev--Manin space at the identity $e$ of
the torus; a drawing showed the exceptional
divisor.

\begin{theorem}[universality, second form]
\label{theorem-tevelev-universality}
Every $\operatorname{Bl}_e\mathbb{P}(\Delta)$
is a rational contraction of
$\operatorname{Bl}_e LM$, and hence of
$\overline{M}_{0,n}$.
\end{theorem}

\begin{remark}
\label{remark-tevelev-wild}
For a lattice polygon $\Delta\subset
\mathbb{R}^2$, the surface
$X=\operatorname{Bl}_e\mathbb{P}(\Delta)$ can
be wild. A drawing showed a polygon with
vertices on a $6\times 6$ grid for which $X$
has a non-polyhedral effective cone. (Misha:
it is singular, since it has interior
lattice points.) There is an elliptic curve
$$
C\colon y^2+y=x^3-x^2-24x+54
$$
in $X$, away from the singularities, with
$C^2=0$. Then
$\mathcal{O}(C)|_C\in\operatorname{Pic}^0(C)$
is the point $(1,5)$, of infinite order, so
$h^0(mC)=1$ for all $m>0$: $C$ is not a
fibre of an elliptic fibration, and every
effective divisor in $|mC|$ is $mC$.
\end{remark}

\begin{theorem}[Nikulin]
\label{theorem-tevelev-nikulin}
If the effective cone of a surface is
polyhedral, it is generated by curves of
negative self-intersection.
\end{theorem}

The original notes give the last words as
``hesf? curves''; ``negative curves'' is my
guess. The argument then reads: if
$\operatorname{Eff}(X)$ were polyhedral,
the class of $C$, which has $C^2=0$, would
lie on a facet, so
$mC\sim xA+yB$ for some curves $A$, $B$
and some $x$, $y$, which would force
$h^0(mC)>1$, a contradiction. The original
notes also mention the cone
$\{D : D^2\geq 0\}$ between
$\operatorname{Nef}(X)$ and
$\operatorname{Eff}(X)$, garbled. Not
checked.

\medskip\noindent
{\bf Two more anomalies.}

\begin{theorem}[Goto--Nishida--Watanabe]
\label{theorem-tevelev-gnw}
There is a weighted projective plane
$\mathbb{P}(a,b,c)$ such that
$X=\operatorname{Bl}_e\mathbb{P}(a,b,c)$, in
characteristic $0$, has a nef and big
divisor that is not semiample.
\end{theorem}

\medskip\noindent
{\bf Conjectural anomaly.} For
$X=\operatorname{Bl}_e\mathbb{P}(9,10,13)$,
the effective cone should be the acute angle
between two rays: $E$, the exceptional
curve, and the ray with $D^2=0$. One must
prove that nothing outside this angle is
effective. This is equivalent to
$\mathbb{P}(9,10,13)$ having an irrational
Seshadri constant at $e$. The conjecture is
that almost every surface has a point with
irrational Seshadri constant, but no example
is known. This case would imply the Nagata
conjecture for $9\cdot 10\cdot 13$ points
on $\mathbb{P}^2$.

The original notes end mid-sentence: ``There
is this world of blown up toric surfaces
whose geometry is very''.



\section{Jacobian elliptic fibrations on K3
surfaces with a non-symplectic automorphism
of prime order}
\label{section-zingali-elliptic-k3}
%% speaker: Felipe Zingali Meira, UFRJ
%% place: Seminário das Sextas, PUC-Rio
%% date: September 20, 2024

\noindent
Felipe Zingali Meira, UFRJ.
Seminário das Sextas, PUC-Rio.
September 20, 2024.

\medskip\noindent
{\bf Abstract.} K3 surfaces are one of the
few classes of algebraic surfaces which
allow more than one distinct relatively
minimal elliptic fibration. There are
distinct ways of classifying such
fibrations. In this talk, we will describe
how to classify these fibrations with
relation to the action of a non-symplectic
automorphism $\sigma$ of prime order on a
K3 surface $X$, and how each type of
fibration relates to a linear system in the
quotient surface $X/\sigma$.

\medskip\noindent
{\bf Sources of these notes.} Daniel's
LaTeX notes typed during the talk,
transcribed here in October 2026. There are
no slides or recording; the drawings made
during the talk are described in prose.
Gaps in the original notes are marked as
such. The original notes started with a
picture of Kodaira's table of singular
fibres, taken from
\href{https://www.math.columbia.edu/~chaoli/docs/EllipticSurfaces.html}{Chao
Li's notes on elliptic surfaces}.

\medskip\noindent
{\bf Definitions.}

\begin{definition}
\label{definition-zingali-elliptic-fibration}
Let $S$ be a smooth projective surface and
$C$ a smooth projective curve. An
\emph{elliptic fibration} is a surjective
morphism $\mathcal{E}\colon S\to C$ such
that all but finitely many fibres
$F_v=\mathcal{E}^{-1}(v)$ are irreducible
curves of genus $1$. It is \emph{Jacobian}
if it has a section $s_0\colon C\to S$, and
\emph{relatively minimal} if no fibre
contains a $(-1)$-curve.
\end{definition}

The generic fibre of such a fibration is a
genus $1$ curve $E$ over the function field
$K(C)$, and there is a bijection between the
points of $E(K(C))$ and the sections
$s\colon C\to S$. Picture: the base is a
curve, the total space is a family of
elliptic curves, and a section is a curve
that meets every fibre in one point.

Facts about a Jacobian relatively minimal
elliptic fibration:
\begin{enumerate}
\item $K_S=(\chi(\mathcal{O}_S)-2)F$, where
$F$ is the class of a fibre in the
Néron--Severi group $\operatorname{NS}(S)$.
Recall that $\operatorname{NS}(S)$ is the
group of divisors modulo numerical
equivalence; here one may think of the
Picard group.
\item $\operatorname{NS}(S)$ is a lattice
with the intersection pairing.
\item The \emph{trivial lattice}
$\operatorname{Triv}(S)$ is spanned by the
zero section $\Sigma_0$, the fibre class
$F$, and the components of reducible fibres
not meeting $\Sigma_0$, and
$$
\operatorname{NS}(S)/\operatorname{Triv}(S)
\cong E(K(C)).
$$
Hence the Shioda--Tate formula
$$
\rho(S)=2+r+\sum_v (m_v-1),
$$
where $r$ is the rank of $E(K(C))$ and $m_v$
is the number of components of $F_v$.
\item The Euler number is
$e(S)=\sum_{v\in C}e(F_v)$, and Noether's
formula reads
$\chi(\mathcal{O}_S)=(e(S)+K_S^2)/12$.
\end{enumerate}
The original notes write the last formula
as $\chi(S)=e(S)/12+K_S^2$ and the sum in
the Shioda--Tate formula without
parentheses; both are corrected here to the
standard forms. Not checked against the
talk.

\medskip\noindent
{\bf Rational elliptic surfaces.} A
\emph{rational elliptic surface} is a
smooth projective rational surface $R$ with
an elliptic fibration
$\mathcal{E}\colon R\to\mathbb{P}^1$.

\begin{example}
\label{example-zingali-cubic-pencil}
Let $f,g$ be cubics on $\mathbb{P}^2$ and
consider the rational map
$$
\varphi\colon \mathbb{P}^2
\dashrightarrow \mathbb{P}^1,
\quad
p\longmapsto [f(p):g(p)].
$$
Blowing up its base points gives a
commutative triangle
$$
\xymatrix{
& R \ar[ld]_{\eta} \ar[rd]^{\mathcal{E}} & \\
\mathbb{P}^2 \ar@{-->}[rr]^{\varphi}
& & \mathbb{P}^1
}
$$
\end{example}

Over an algebraically closed field, every
relatively minimal Jacobian rational
elliptic surface arises by blowing up the
nine base points of a pencil of cubics. In
particular $\rho(R)=10$.

\medskip\noindent
{\bf K3 surfaces.}

\begin{definition}
\label{definition-zingali-k3}
A \emph{K3 surface} is a smooth projective
surface $X$ with
$q(X)=h^1(X,\mathcal{O}_X)=0$ and
$K_X=0$, that is, there is a nowhere
vanishing $\omega_X\in H^0(X,\Omega^2_X)$.
\end{definition}

\begin{proposition}
\label{proposition-zingali-euler-24}
Let $\mathcal{E}\colon X\to\mathbb{P}^1$ be
a Jacobian relatively minimal elliptic
fibration. Then $X$ is a K3 surface if and
only if $e(X)=24$.
\end{proposition}

The computation in the notes is
$$
q(X)=h^2(X,\mathcal{O}_X)
+h^0(X,\mathcal{O}_X)-\chi(\mathcal{O}_X)
=1+1-2=0.
$$

\begin{definition}
\label{definition-zingali-non-symplectic}
Let $X$ be a K3 surface and
$\sigma\in\operatorname{Aut}(X)$, so that
$\sigma^*\omega_X=\xi\omega_X$ for a root of
unity $\xi$. Then $\sigma$ is
\emph{symplectic} if $\xi=1$ and
\emph{non-symplectic} otherwise.
\end{definition}

If $\sigma$ has finite order $n$, the
quotient $X/\sigma$ is
\begin{itemize}
\item a K3 surface, if $\sigma$ is
symplectic;
\item a rational surface, if $\sigma$ is
non-symplectic and either $n\geq 3$, or
$n=2$ and $\operatorname{Fix}(\sigma)\neq
\emptyset$;
\item an Enriques surface, if $n=2$ and
$\operatorname{Fix}(\sigma)=\emptyset$.
\end{itemize}
Here the quotient is understood after
resolving its singularities.

\medskip\noindent
{\bf Base change.} How to build K3 surfaces
out of rational elliptic surfaces. Let
$\mathcal{E}\colon R\to\mathbb{P}^1$ be a
rational elliptic surface and
$\tau\colon\mathbb{P}^1\to\mathbb{P}^1$ a
cyclic cover of degree $n$. Take the fibre
product and resolve it:
$$
\xymatrix{
R \ar[d]_{\mathcal{E}}
& R\times_{\mathbb{P}^1}\mathbb{P}^1
\ar[l] \ar[d] & Y \ar[l] \ar[d] \\
\mathbb{P}^1 & \mathbb{P}^1 \ar[l]^{\tau}
& X \ar[l]^{\mathcal{E}_X}
}
$$
The original notes write ``$n-1$ cover'' and
do not label the arrows into $Y$; the
diagram above is my reading.

\begin{remark}
\label{remark-zingali-sergey}
(Sergey, from the audience.) The fibre
product may be singular at points lying over
critical points of both $\mathcal{E}$ and
$\tau$.
\end{remark}

The Kodaira type of the fibre
$F^X_v=\mathcal{E}_X^{-1}(v)$ is determined
by the Kodaira type of
$F_u=\mathcal{E}^{-1}(u)$, where
$u=\tau(v)$, and the ramification index
$r(v|u)$; there are tables for this. Several
drawings followed.

For $n=3$: $\mathcal{E}_X\colon X\to
\mathbb{P}^1$ is a K3 surface if
$F^X_a$ is of type $\mathrm{IV}$ or
$\mathrm{I}^*_n$ and $F^X_b$ is of type
$\mathrm{I}_n$, $\mathrm{II}$ or
$\mathrm{III}$, as recorded in the original
notes, which do not say what $a$ and $b$
are. Not checked.

\medskip\noindent
{\bf Towards the classification.}

\begin{definition}
\label{definition-zingali-types}
Let $X$ be a K3 surface,
$\sigma\in\operatorname{Aut}(X)$
non-symplectic of prime order $p$, and
$\mathcal{E}\colon X\to\mathbb{P}^1$ a
Jacobian elliptic fibration (on a K3
surface it is automatically relatively
minimal, by adjunction). Then
$\mathcal{E}$ is of
\begin{enumerate}
\item \emph{type 1} if $\sigma^*F=F$ and
$\sigma$ acts trivially on $\mathbb{P}^1$;
\item \emph{type 2} if $\sigma^*F=F$ and
$\sigma$ acts on $\mathbb{P}^1$ with order
$p$;
\item \emph{type 3} if $\sigma^*F\neq F$.
\end{enumerate}
\end{definition}

Three drawings showed the types. In type 1,
$\sigma$ moves points within each fibre and
fixes the base. In type 2, $\sigma$ moves
fibres to fibres and acts on the base. In
type 3 it is messier: $\sigma$ moves fibres
to horizontal curves.

Assumption: $\sigma$ acts trivially on
$\operatorname{NS}(X)$. The original notes
attribute this setting to ``Sarti,
Arteloni, Taki'', presumably Artebani,
Sarti and Taki. Not checked.

\medskip\noindent
{\bf Type 1.} Then $\sigma$ acts on each
smooth fibre $F_v$, and the order of an
automorphism of an elliptic curve fixing a
point is $2$, $3$, $4$ or $6$.
\begin{itemize}
\item $p=2$ (Garbagnati--Salgado, 2017):
the fibres are of types $\mathrm{I}_n$,
$\mathrm{I}^*_n$, $\mathrm{III}$.
\item $p=3$ (2024): the fibres are of types
$\mathrm{I}$, $\mathrm{II}^*$, $\mathrm{IV}$,
$\mathrm{IV}^*$, $\mathrm{I}_0$,
$\mathrm{I}^*_0$.
\end{itemize}
The original notes write ``Garbognati'' and
list $\mathrm{I}$ without an index for
$p=3$; not checked.

For $\operatorname{ord}(\sigma)\geq 3$ the
fixed locus may contain isolated points:
$$
\operatorname{Fix}(\sigma)=
C_g\cup C_1\cup\cdots\cup C_k
\cup\{p_1,\ldots,p_l\},
$$
where $C_g$ is a curve of genus $g$ and
$C_1,\ldots,C_k$ are rational curves. For
$p=2$, $\operatorname{Fix}(\sigma)\neq
\emptyset$.

The quotient $X/\sigma$ has a smooth
rational resolution $\tilde{R}$:
$$
\xymatrix{
\tilde{R} \ar[d]
& \tilde{X} \ar[l]_{\hat{\pi}} \ar[d] \\
X/\sigma & X \ar[l]^{\pi}
}
$$
and the fibration $\mathcal{E}$ induces a
linear system of curves $\Lambda$ on
$\tilde{R}$:
\begin{enumerate}
\item if $\mathcal{E}$ is of type 1,
$\Lambda$ is a pencil of conics;
\item if $\mathcal{E}$ is of type 2,
$\Lambda$ is a pencil of genus $1$ curves;
\item if $\mathcal{E}$ is of type 3,
$\Lambda$ is a non-complete linear system.
\end{enumerate}
The original notes end here.



\section{Smoothable compactified Jacobians of
nodal curves}
\label{section-smoot-compa-jacob-nodal}
%% speaker: Nicola Pagani, University of
%% speaker: Liverpool and Bologna
%% place: Seminar of Algebraic Geometry, UFF
%% date: August 20, 2025

\noindent
Nicola Pagani, University of Liverpool and
Bologna.
Seminar of Algebraic Geometry UFF.
August 20, 2025.

\medskip\noindent
{\bf Abstract.} Building from examples, we
introduce an abstract notion of a
'compactified Jacobian' of a nodal curve. We
then define a compactified Jacobian
to be 'smoothable' whenever it arises as the
limit of Jacobians of smooth
curves. We give a complete combinatorial
characterization of smoothable
compactified Jacobians in terms of some 'vine
stability conditions', which we
will also introduce. This is a joint work
with Fava and Viviani.

\medskip\noindent



Let $C$ be a smooth curve and $d \in
\mathbb{Z}$. Define
$$
J^d_C=\{L:\text{$L$ is a line bundle of
degree $d$}\}/\sim
$$
which is a smooth projective variety of
dimension $g(C)$.

If $C$ is nodal we still can consider
$J^d_C$.
\begin{enumerate}
\item One connected component. Then the
Jacobian is $\mathbb{P}^1$ minus two
points. This is not universally closed, so it
is not proper.
\item Two components intersecting at one
point.
The pullback of the normalization splits the
degree in
intintely many ways, givieng that $J^{-1}_C$
is an infinite set of points. This
is not of finite type, so it is not proper.
\item The curve has two components
intersecting at two points.
This gives $J^{-2}_C$,
 which is a mixture of the two former items.
 (Probably not proper too.)
\end{enumerate}

\medskip\noindent
Now consider
$$
\text{TF}_C^d=\{\mathcal{F}:
\text{coherent on $C$, torsion-free, rank-1
on $C$}\}/\sim
$$
This satisfies the existence oart of the valu
point of properness.

Now we consider the moduli. Now we consider
the ideal sheaf of the (singular?)
point(s?):
\begin{enumerate}
\item One component. The stack is proper!
\item Two components intersecting once. Now
we get stacky points,
$x=[\bullet/\mathbb{G}_m]$. These points have
generic stabilizer. The resulting
stack is not separated because a morphism of
a curve, say $\mathbb{P}^1$ minus a
point … there are infinitely many ways to
extend a morphism from this thing to a
line bundle. So you cannot include any of
these stacky points. Recall that a
sheaf is {\it simple} if its automorphism
group is $\mathbb{G}_m$.
\item  The ideal sheaf of both nodes
$\mathcal{I}(N_1,N_2)$ has a positive
dimensional automorphism
group. The stack is not proper.
\end{enumerate}

\begin{definition}
\label{definition-fine-compactified-Jacobian}
A {\it fined compactified Jacobian} of $C$ is
an open connected substack of
$\text{TF}^d(C)$ that is also proper.
\end{definition}

\begin{remark}
\label{remark-algebraic-space}
This thing is automatically an algebraic
space.
\end{remark}

\begin{definition}
\label{definition-compactified-Jacobian}
A {\it compactified Jacobian} is an open
connected of $\text{TF}^d(C)$ that
admits a proper, good moduli space.
\end{definition}

Consider the Artin stack $\mathfrak{X}
\xrightarrow{\Gamma}X$ […]
 is a {\it good
moduli space} if
\begin{enumerate}
\item Every moduli factors
$$
\xymatrix{
\mathfrak{X} \ar[r]\ar[rd]&\mathcal{I}\text{
(ACC. space)}\\
& X
}
$$
\item
$\pi_*\mathcal{O}_{\mathfrak{X}}=\mathcal{O}
_X$.
\end{enumerate}

\medskip\noindent
We expect to find a notion of stability
condition to produce these things
[…] $[\bullet/ \mathbb{G}_m]$ would be the
polystable representative.

\begin{definition}
\label{definition-smoot-compa-jacob}
A compactified Jacobian $\overline{J_C}$ is
{\it smoothable} if all smoothings
$\mathcal{C}\to \Delta=\{0,\eta\}$ (with
$\mathcal{C}_0=C$),
$$
J^d_{\mathcal{C}_\eta}\cup C \to
\overline{J_C}
$$
is proper.
\end{definition}

\medskip\noindent

\begin{definition}
\label{definition-connected}
Let $X$ be a curve.
$$
\text{BCON}(X)=\{Y \subseteq X\text{s.t.
}Y,Y^c \text{ are connected}\}
$$
\end{definition}

\begin{definition}
\label{definition-v-satbility}
A $v$-curve is a generalization of items (2)
and (3) in the lists above
[it looks like two long snakes $\sim$ that
intersect several times, and $t$ is
the number of nodes].
A {\it $v$-condition} is a pair $n=(n_1,n_2)$
such that
$$
n_1+n_2=\begin{cases}
d+1-t\qquad &\text{we say the s.c. is
nondegenerate} \\
d-t\qquad &\text{degenerate}
\end{cases}
$$
$\mathcal{F}$ on $X$ is  {\it
$n$-(semi)stable} if
$\text{deg}\mathcal{F}_{X_i}>n_i$
($\text{deg}\mathcal{F}_{X_i}\geq n_i$)
for $i=1,2$.
\end{definition}
$\mathcal{F}_{X_i}=\mathcal{F}|_{X_i}$
torsion.
$$
\text{deg}(\mathcal{F}_{X_i})+\text{deg}
(\mathcal{F}_{X_2})
=d-|\text{sing}(F)|.
$$
Then
$$
\overline{J_C}(n)=\{\mathcal{F}\text{ is
semistable}\},
$$
a smooth compact Jacobian.

\begin{definition}
\label{definition-n-degeneracy}
A {\it degeneratation} of $v$-stab. on $X$ is
$n:\text{BCON}(X) \to \mathbb{Z}$
such that
\begin{enumerate}
\item
$$
n_Y+n_{Y^c}+|Y \cap Y^c|=
\begin{cases}
d+1\qquad &\text{ we say $Y$ is
$n$-nondegenerate} \\
d\qquad &\text{$Y$ is $n$-degenerate}
\end{cases}
$$
\item $Y_i$ no pa. common component
$n_{Y_1}+n_{Y_2}+\ldots$
$$
\xymatrix{
&Y_1\ar[dr]\ar[dl]\\
Y_2\ar[ur]\ar[rr]&&Y_3\ar[ll]\ar[ul]
}
$$
\end{enumerate}
\end{definition}

\begin{theorem}[-, et al]
\label{theorem-bijec-stabi-condi-nodal}
(bijection between stability conditions and
nodal curves) The map
$$
\left\{ \substack{\text{sm. comp.} \\
\text{Jac of $X$}} \right\} \to
\left\{ \substack{\text{$v$-stab.} \\
\text{cond. of $X$}} \right\}
$$
$$
n\mapsto
\overline{J_X}(n)=\{n\text{-semistable
sheaves}\}
$$
is a bijection. (The arrow should be from
right to left!)
\end{theorem}

F. Viviani had proved it for fine compact
Jacobians.

\section{Equivariant spaces of matrices of
constant rank}
\label{section-equiv-space-matri-const}
%% speaker: Ada Boralevi, France
%% place: Algebraic Geometry Seminar, IMPA
%% date: August 27, 2025

\noindent
Ada Boralevi, France.
Algebraic Geometry Seminar, IMPA.
August 27, 2025.

\medskip\noindent
{\bf Abstract.} A space of matrices of
constant rank is a vector subspace V, say
of dimension n+1, ofthe set of matrices of
size axb over a field k, such that
any nonzero element of V has fixed rank r. It
is a classical problem to look for
different ways to construct such spaces of
matrices. In this talk I will give an
introduction up to the state of the art of
the topic, and report on my latest
joint project with D. Faenzi and D. Fratila,
where we give a classification of
all spaces of matrices of constant corank one
associated to irreducible
representation of a reductive group.

\medskip\noindent

We are interested in vector spaces $U \subset
\text{Mat}_{m,n}(\mathbb{C})$,
with $m \leq n$, of {\it constant rank}, i.e.
such that for all $f \in U$,
$r:=\text{rank}f$ is the same.

Let $\ell(r,m,n):=\text{max}\dim U:U$ is of
rank $r$.

\medskip\noindent
{\bf Questions.}
\begin{enumerate}
\item $\ell(r,m,n)=$? In general not known.
\item Find relations among $\ell, r,m$ and
$n$.
\item Construction of examples and
classification.
\end{enumerate}

\begin{example}
\label{example-what-we-know}
\begin{enumerate}
\item $\ell(1,m,n)=n$,
$$
\begin{pmatrix}
x_1&x_2&\cdots &x_n\\
\\
\\
\\
\end{pmatrix}
$$
\item $\text{rank}=2$? There are two cases
([Atkinson '83], [Eisenbud-Haus '88])
\begin{itemize}
\item Compression spaces,
$$
\begin{pmatrix}
*&*&\cdots &*\\
*&0&\cdots & 0\\
\vdots &\\
*&0&\cdots&0
\end{pmatrix}
$$
\item Skew-symmetric matrices of $3 \times
3$.
\end{itemize}
\item $\ell(r,m,n) \geq n-r+1$. Because you
can put a matrix of
$m\times n$ (with the first $r$ rows that can
have nonzero entries):
$$
\begin{pmatrix}
x_1 &  x_2 &  \cdots & & x_{n-r+1}\\
& x_1 & x_2 & \cdots &\\
\\
\\
\\
\end{pmatrix}
$$
\end{enumerate}
\end{example}

\begin{theorem}[Westwick '86]
\label{theorem-Westwick}
\begin{enumerate}
\item $n-r+1 \leq \ell(r,m,n) \leq 2n-2r+1$.
\item If $n-r+1 \not|
\frac{(m-1)!}{(r-1)!}\implies \ell$
\end{enumerate}
\end{theorem}

We can see these spaces as (subvarieties?) of
determinantal varieties
$M_n=\{f \in
\text{Mat}_{m,n}(\mathbb{C}):\text{rank}(f)
\leq  r\}$.
[Interpretation via secant varieties inside
Segre embedding.]

\medskip\noindent
Consider a map $\varphi: U \to \Hom(V,W)$.
Then $\varphi \in U^* \otimes  V^* \otimes
W$.
We get that
$\varphi$ is of constant rank if and only if
some kernel, image and cokernel are vector
bundles.

\medskip\noindent
{\bf Focus of today.} What happens when
$U,V,W$ are irreducible
representations of a complex reductive group
$G$?

\medskip\noindent
{\bf Question.} What is the natural
equivariant morphism
$$
U \to \Hom(V,W)=V^* \otimes  W
$$
of constant rank?

\medskip\noindent
Consider the case of $G=\text{SL}_2$. All the
irreducible representations (which are
self-dual) are given by
$V(m-1)\cong
\mathbb{C}[x,y]_{\text{deg}=m-1}$.

Recall the Clebsh-Gordon decomposition ($m
\leq  n$)
$$
V(m-1)\otimes
V(n-1)=\bigoplus_{i=0}^{m-1}V(n-m+2i)
$$

\begin{theorem}[B-Faenzi-Lella '22]
\label{theorem-BFL22}
$$
V(n+m-2)\hookrightarrow \Hom(V(m-1),V(n-1))
$$
is of constant rank (corank 1) if and only if
$$
n-m+2i|m-1
$$
\end{theorem}

\begin{theorem}[B. Faenzi, Fratila '25]
\label{theorem-BFF25}
Let $V(\nu)$, $V(\mu)$ and $V(\lambda)$ be
irreducible representations
of a comple reductive group $G$, with
$$
\dim(V(\mu)) \leq \dim(V(\lambda))_-
$$
If there exists a morphism of representations
$$
\varphi: V (\nu) to \Hom (V (\lambda),V(\mu))
$$
then $\varphi$ is of constant corank 1 if and
only if
there exists a simple root $\alpha_i$ such
that
\begin{enumerate}
\item $\lambda=\mu+\nu-\alpha_i$,
\item $\nu$ is a multiple of $\nu$,
\item  $\nu$ is a multple of $\omega_i$
\end{enumerate}
\end{theorem}

\section{Irreducibility of the Hilbert scheme
of points and the class of 2-step
ideals}
\label{section-graffeo-two-step-ideals}
%% speaker: Michele Graffeo, SISSA
%% place: Algebraic Geometry Seminar, UFF
%% date: September 29, 2025

\noindent
Michele Graffeo, SISSA.
Algebraic Geometry Seminar, UFF.
September 29, 2025.

\medskip\noindent
These are rough notes from the talk.
Their incomplete formulas and theorem
statements still need review.
They are preserved as a reading record,
not as verified statements.

\medskip\noindent
{\bf Abstract.}
Hilbert schemes of points on a
quasi-projective variety X are classical
objects
in algebraic geometry. Roughly speaking, they
parametrise ideals of a polynomial
ring with complex coefficients having finite
colength. Although Hilbert schemes
always have a distinguished component called
the smoothable component, their
geometry is quite pathological and one of the
main open problems around them
concerns their irreducibility. In a joint
work with Giovenzana, Giovenzana,
Lella we introduce the notion of 2-step
ideals. We show that the loci
parametrising these ideals are in general not
contained in the smoothable
component of the Hilbert scheme, thus
providing new examples of
extra-components.  In my seminar I will
discuss our class of ideals and relate
them to the compressed algebras considered by
Iarrobbino in the eighties.
Finally, I will show how to extend our result
to the nested setting.


\medskip\noindent

The talk uses the notation
$\underline{\operatorname{Hilb}}^d(X)$
for the functor of points in Definition
\ref{hilbert-schemes-definition-%
hilbert-scheme-of-points}.

\begin{remark}
\label{definition-hilbert-scheme}
The functor also makes sense for a smooth
quasi-projective variety $X$. For a family
$Z\subset B\times X$, the projection
$\pi:Z\to B$ is finite locally free of
rank $d$. Over an affine open in $B$, this
says that the finite algebra of the family
is locally free of rank $d$ as a module.
Every fibre then has length $d$.
\end{remark}

\begin{theorem}[Grothendieck]
\label{theorem-grothendieck}
$\underline{Hilb}^d(X)$ is represented by a
quasi-projective
scheme $Hilb^d(X)$
\end{theorem}

In particular, a complex point corresponds
to a length-$d$ subscheme of $X$.

A fat point $Z= \text{Spec}(A)$ $(A,
\mathfrak{m})$ is a
local artinian $\mathbb{C}$-algebra of finite
type.

If $X$ smooth can consider $X=\mathbb{A}^n$,
and it's the same to consider
\begin{itemize}
\item $Z \hookrightarrow \mathbb{A}^n$ of
length $d$.
\item $I \subset \mathbb{C}[x_1,\ldots,x_n]$
of
colength $d$.
\item $A= \mathbb{C}[x_1,\ldots,x_n]/I$ with
$\dim_\mathbb{C}=d$.
\end{itemize}

\begin{definition}
\label{definition-nested-hilbert-scheme}
$X$, $0 \leq d_1 \leq \ldots
d_r=\underline{d}$.
$$
Hilb^{\underline{d}}(X)=
\{Z_1\not\hookrightarrow
\ldots
\not\hookrightarrow Z_r  \not\hookrightarrow
X:
\text{len}(Z_i)=d_i\}
$$
\end{definition}

\begin{theorem}[Hartshorne (r=0), Fogarty,
Kalpan (r>1)]
\label{theorem-hk}
$Hilb^d(X)$ is connected.
\end{theorem}

For $r=1$, $Hilb^d(\mathbb{A}^n)$ is smooth
iff
$n \leq 2$, $d \leq 3$.
If $r>1$, $n=2$, it is smooth iff $r=2$ and
$d_1=d_2=-1$.
If $r>1$ and $n>2$, it is smooth iff $r=2$,
$(d_1,d_2)=\{(1,2),(2,3)\}$.

Iror components.
\begin{itemize}
\item $(r=1)$ $n\geq 4$,
$Hilb^d(\mathbb{A}^3)$
is irreducible iff $d\leq 7$ ($\impliedby$
[E-I], $\implies $ [Mozzola].
$n=3$ irredicuble if $d \leq  11$ (8,9 Sivic;
10 Jardim et al; 11, Jelisievv et
al.)
\item $(r=2)$, $n=2$, irreducible
[G-Rosul-Sebastian].
\item $(r \geq 3)$ $n=2$, there exist $d_1
\leq \ldots \leq d_5$
such that
$Hilb^{\underline{d}}(\mathbb{A}^i)$ is
irreducible [S-R].
\end{itemize}

There is an analogy between classical in
dimension 3
and nested in dimension 2.

Schematic structure.
\begin{itemize}
\item [I] $Hilb^{21}(\mathbb{A}^4)$ has at
least a generically
non-reduced component.
\item Problem: what about $n=3$?
\item Theorem. If $Hilb^{d_1,\ldots
d_r}(\mathbb{A}^i)$ is
irreducible $\implies $
$Hilb^{(1,d_1,\ldots,d_r)}(\mathbb{A}^n$
has a generically non reduced component.
\end{itemize}

Other results.
\begin{itemize}
\item …
\end{itemize}
Hilbert scheme function.
For a local artinian $\mathbb{C}$-algebra of
finite type there is an associated graded
$\text{Gr}_\mathfrak{m}(A)
=\bigoplus_{i \geq
0}\mathfrak{m}^1/\mathfrak{m}^{i+1}$
which allows us to define the Hilbert
function
$$
\begin{aligned}
h_A: \mathbb{Z} &\longrightarrow \mathbb{N}
\\
i &\longmapsto \dim_\mathbb{C}
\mathfrak{m}_i/\mathfrak{m}_{i+1}
\end{aligned}
$$
sistability for graded $A$-modules.
The socol is the annihilator of
the ideal which gives the algebra
e.g. $\mathbb{C}[x,y](x^2,xy,y^2)$ then
$\text{Soc}(A)=\left\langle{}x,y^2\right\rangle{}_\mathbb{C}
$.

\medskip\noindent
\begin{definition}
\label{definition-smoothable-component}
The {\it smoothable component} is
 $$
V_{sm}=\overline{\{[(z_i)_{i=1}^n \in
Hilb^{\underline{d}}(\mathbb{A}^n):
Z_r\text{ is reduced}\}}
$$
\end{definition}

\begin{definition}
\label{definition-elementary-componenet}
$V \subset
Hilb^{\underline{d}}(\mathbb{A}^n)$ is
an {\it elementary component} if
$\forall [Z_1,\ldots,Z_r] \in V$
then $Z_r$ is a fat point.
\end{definition}

Instead of looking for all components of the
Hilbert scheme
we look for elementary components. This may
not be some simple;
sometimes there's infinitely many. There
is no method to find them.

\begin{theorem}[Irrabaro]
\label{theorem-i}
Every irreducible component $V \subset
Hilb^d(\mathbb{A}^n)$
is generically étale locally product of
elementary components.
\end{theorem}

\begin{definition}
\label{definition-}
Let $h:\mathbb{Z} \to \mathbb{N}$ be a
function
with finite support, $|h|=\sum_{ i \in
\mathbb{Z}}h(i)=d$. Define
\begin{itemize}
\item $H_n=\{[A] \in Hilb^d(\mathbb{A}^n):
h_A \equiv h\}$.
This is closed by semicontinuity of the
function; but can be expressed
a by a representable functor.
\item The following also has a canonical
schematic structure
$$
\begin{aligned}
\pi_n:H_n  &\longrightarrow
\mathcal{H}_n=\{[A] \in H_n: A\text{ is
graded}\} \\
A &\longmapsto A=\text{Gr}_{\mathfrak{m}}(A)
\end{aligned}
$$
Where the first is $\not\hookrightarrow $ in
the second.
\end{itemize}
\end{definition}

Recall that
$$
T_{[I]}Hilb^d(\mathbb{A}^n)=\Hom_R(I,R/I)
$$
\begin{theorem}[B-B,J,G G G L]
\label{theorem-b-b-j-g-g-g-l}
There is a graded $[I] \in \mathcal{H}_n$.
Then
$$
\begin{aligned}
T_{[I]}\mathcal{H}_n&=\Hom_R(I,R/I)_{=0}\\
T_{[I]}\pi_n([I])&='>0\\
T_{[I]}H_n&=' \geq 0.
\end{aligned}
$$
\end{theorem}

\begin{definition}
\label{definition-compressed}
$(A,\mathfrak{m}_A)$ is {\it compressed} is
for all  $(A',\mathfrak{m}')$
with $e(A)=e(A')$ and $len(A) \geq  len(A')
$.
\end{definition}

\begin{theorem}
\label{theorem-e}
$E_h \subset H_n$ is $H$-compressed
$E_h \neq  \emptyset$
\end{theorem}

\section{Normal and Restricted Tangent
Bundles of Rational
Curves in Hypersurfaces}
\label{section-restricted-tangent-bundles}
%% speaker: Lucas Mioranci, IMPA
%% place: Algebra Seminar, IMPA
%% date: October 8, 2025

\noindent
Lucas Mioranci, IMPA.
Algebra Seminar, IMPA.
October 8, 2025.

\medskip\noindent
{\bf Abstract.}
Let $X\subset \mathbb{P}^n$ be a degree $d$
hypersurface containing a smooth
rational curve $C$ of degree $e$. The normal
bundle $N_{C/X}$ and the restricted
tangent bundle $T_X|_C$ split as direct sums
of line bundles of the form
$\bigoplus_i
\mathcal{O}_{\mathbb{P}^1}(a_i)\qquad$ called
their splitting type.
The splitting type of $N_{C/X}$ and
${T_X}|_C\;\;$ controls the local structure
of the space of rational curves and
determines how many general points of $X$ we
can interpolate by deforming the curve $C$.
By combining explicit computations
of $T_X|_C\;\;$ and an induction argument, I
classify all triples $(e,d,n)$ such
that a general degree $d$ hypersurface
$X\subset \mathbb{P}^n$ contains a
rational curve $C$ of degree $e$ whose
restricted tangent bundle ${T_X}|_C\;\;$
is balanced. The case of quadrics is
particular, in which we show that
odd-degree rational curves do not interpolate
the expected number of points on
quadric hypersurfaces.  For the normal
bundle, I compute explicit examples of
hypersurfaces $X$ for all possible splitting
types of $N_{C/X}$ when $C$ is the
rational normal curve. Additionally, for
$d\ge 3$, we compute the dimension of
the space of hypersurfaces $X$ such that
$N_{C/X}$ has a given splitting type.



\medskip\noindent

$\mathcal{O}$ with no subindex means
$\mathcal{O}_{\mathbb{P}^1}$.
$$
\xymatrix{
&0\ar[d]&0\ar[d]\\
&T_{\mathbb{P}^1}=\mathcal{O}_(2)\ar[d]\ar[r]
^=
&T_{\mathbb{P}^1}=\mathcal{O}(2)\ar[d]\\
0\ar[r]&T_{X}|_{C}\ar[d]\ar[r]&T_{\mathbb{P}
^n}|_{C}\ar[d]\ar[r]&
N_{X/\mathbb{P}^n}|_{C}\cong
\mathcal{O}(de)\ar[d]_=\ar[r]&0\\
0\ar[r]&
N_{C/X}\ar[d]\ar[r]&N_{C/\mathbb{P}^n}\ar[d]
\ar[r]&
N_{X/\mathbb{P}^n}|_{C}\cong
\mathcal{O}(de)\ar[r]&0\\
&0&0
}
$$
Throughout this talk
$X$ is a degree $d$ hypersurface in
$\mathbb{P}^n$
and $C \subset X$ is a smooth rational curve
of degree $e$.

Since $C$ is rational we can think of it
as a map $f:\mathbb{P}^1 \to X$.
Then we can pullback any bundle over $C$
back to $\mathbb{P}^1$,
and by Birkhoff-? Theorem
we know that vector bundles over
$\mathbb{P}^1$
split as
$$
E\cong \bigoplus_{i=1}^r
\mathcal{O}_{\mathbb{P}^1}(a_i),
\qquad a_1\leq \ldots \leq a_n.
$$
This decomposition is called the
{\it splitting type} of $E$.
We say $E$ is {\it balanced} if
$|a_i-a_j|\leq 1$ for all $i,j$.

Balancedness is an open condition,
i.e. if you have a family of bundles over
$\mathbb{P}^1$
and one of them is balanced, then all of them
are,
(also you can think it's a generic
condition).

We know the normal bundle $N_{C/X}$ has rank
$n-2$
and degree  $e(n-d+1)-2$.
And the tangent bundle $T_{X/C}$
has rank $n-1$ and degree $e(n-d+1)$.

\begin{example}
\label{example-conic}
A conic $C \subset \mathbb{P}^3$.
Since a conic in $\mathbb{P}^3$ is contained
in a plane,
this obliges to have a ``distinguished normal
direction'',
$\mathcal{O}(2)$. Indeed,
$$
\begin{aligned}
N_{C/\mathbb{P}^3}&\cong(N_{Q/\mathbb{P}^3}
\oplus
N_{\mathbb{P}^2/\mathbb{P}^3}|_{C}\\
&\cong(\mathcal{O}(2) \oplus
\mathcal{O}(1))|_{C}\\
&\cong\mathcal{O}(1)\oplus \mathcal{O}(2).
\end{aligned}
$$
\end{example}

\medskip\noindent
Now we discuss interpolation of general
points.
Let $p_1,\ldots,p_n \in \mathbb{P}^1$
and $x_1,\ldots, x_n \in X$
be general points.
Let $f:\mathbb{P}^1 \to X$ be such that
$f(p_i)=x_i$.

Consider the space of morphisms mapping
the  $p_i$ to $x_i$. It's known that its
tangent space
at $f$
is given by
$$
T_{[f]}\Mor(\mathbb{P}^1,X,p_i \mapsto x_i)
\cong H^0(\mathbb{P}^1,f^*T_X(-n)).
$$
\noindent
Deformations of $f$ (mapping $p_i \mapsto
x_i$)
dominate $X$ if $a_i-n_i \geq 0$.
Also, deformations of $f$ interpolate up
to $a_1+f$ general points in $X$.

\medskip\noindent
How many points can deformations of the curve
interpolate?
This means (I think) that the deformations of
the curve
are still such that $p_i \mapsto  x_i$.
``And for the normal bundle, interpolations
means that the curves contain the points''.

If $T_X|_{C}$ is balanced, $C$ can
interpolate up to
$$
\left\lfloor \frac{e(n+1-d)}{n_1}
\right\rfloor +1
$$
general points. The case of $N_{C/X}$ is
similar.

\medskip\noindent
Next we survey some relevant results.

\begin{theorem}[Coskun-Riedl, 2018]
\label{theorem-CR18}
A general Fano $X$ of degree $d \geq 2$
contains a degree $c$ rational smooth curve
 $C$ with balanced $N_{C/X}$
for every $1 \leq e \leq  n$.
\end{theorem}

The idea is to observe that the space of
curves has
balanced normal bundle.

\begin{theorem}[Ran, 2021]
\label{theorem-R21}
Extends the above for
$1 \leq  e \leq 2n-2$
and $d \geq 4$.
\end{theorem}

Idea: general hypersurface contains a curve
of balanced normal bundle.

\begin{theorem}[Ran, 2024]
\label{theorem-R24}
$X$ general Fano hypersurface.
There exist smooth rational curve with  $C$
with balanced tangent bundle $T_{X}|_{C}$
for $e$ in some arithmetic progressions.
\end{theorem}

$e$ is in general very large.
They show there are curves of arbitrarily
large
degrees with balanced tangent bundle.

\medskip\noindent
Now we describe the new results.

\begin{theorem}
\label{theorem-Lucas}
Let $X \subseteq \mathbb{P}^n$ be Fano
hypersurface of degree $3 \leq d \leq n$.

\begin{enumerate}
\item If $C \subset X$ is a rational smooth
curve
of degree $e$ with
$e \leq  \frac{n-1}{n+1-d}$,
then $T_X|_{C}$ is not balanced.

\item A general hypersurface $X$
contains rational curves $C$
of degree $e$ with balanced
$T_X|_{C}$ for every
$e> \frac{n-1}{n+1-d}$.
\end{enumerate}
\end{theorem}

\begin{theorem}
\label{theorem-lucas2}
Let $X \subseteq \mathbb{P}^n$ be a smooth
quadric hypersurface.

\begin{enumerate}
\item If $e \leq 2$ is even,
then exists a curve $C$ of degree $e$
with $T_X|_{C}\cong \mathcal{O}(e)^{n-1}$.
(I.e. we can interpolate $e+2$ points
if $e$ is even.)

\item If $e \geq 1$ is odd,
then there is no curve $C$ of degree $e$
with balanced $T_X|_C$.
The best we can do is the most balanced
bundle
before balancedness, namely
$T_X|_C \cong
\mathcal{O}(e-1)\oplus\mathcal{O}(e)
\oplus\mathcal{O}(e+1)$.
(I.e. we can only interpolate $e$ points if
$e$ is odd.)
\end{enumerate}
\end{theorem}

Next we give the main ideas in the proof.
Let $C$ be the rational normal curve
of degree $e$ in $\mathbb{P}^1$,
$$
\begin{aligned}
C:\mathbb{P}^1  &\longrightarrow \mathbb{P}^n
\\
(s:t) &\longmapsto
(s^2:s^{e-1}t:\ldots:t^e:0:\ldots:0).
\end{aligned}
$$
$X=V(F)$.

Then find some examples. Then,

\begin{proposition}
\label{proposition-proposition}
If $\mu(T_X|_{C})=\frac{e(n+1-d)}{n-1}\leq 1$
then $T_X|_{C}$ is not balanced.
\end{proposition}

\noindent
Use the proposition to show that $T_X|_C$
is not balanced if $d \leq  n$
and $e le \frac{ n-1}{n+1-d}$.

This implies Theorem \ref{theorem-Lucas}.

\begin{lemma}
\label{lemma-split-normal-bundle}
If $N_{C/X}\cong \bigoplus_i
\mathcal{O}(a_i)$
with $a_i < 4$ for all $i$,
then $T_{X}|_C \cong N_{C/X} \oplus
\mathcal{O}(2)$.
\end{lemma}

\noindent
This implies that
$\text{Ext}^1(N_{C/X},\mathcal{O}(2))=0$.

The most important step is the following,
used as the induction step:

\begin{proposition}
\label{proposition-induction}
If $T_Y|_C$ is balanced for some degree $d$
hypersurface $Y \subseteq \mathbb{P}^{n-1}$,
then there exists a degree $d$ hypersurface
$C \subseteq \mathbb{P}^n$ with balanced
$T_X|_C$.
\end{proposition}

\noindent
Combining the examples, Lemma
\ref{lemma-split-normal-bundle}
and Proposition \ref{proposition-induction},
we establish that for $e \leq
\text{max}\{2d-2,n\}$.

The following argument is used to show that
actually what we showed so far is enough for
all degrees.

\begin{lemma}
\label{lemma-smoothing-argument}
$C=C_1 \cup  C_2$, $C_1 \cong \mathbb{P}^1$.
$E$ vector bundle on $C$ such that
$E|_{C_1}$ is balanced,
$E_{C_2}$ is perfectly balanced
(numerical condition that works when
$e=n-1$).
Then if $E$ is the specialization of a family
of vector bundles $E'$ on $\mathbb{P}^1$,
then $E'$ is balanced.
\end{lemma}

\noindent
Basically, glue and smooth low degree curves
(i.e. $e \leq  \text{max}\{2d-2,n\}$)
to get all degrees


\medskip\noindent
Now we prove the result for quadrics,
i.e. $X$ is a quadraic.
For even $e$ compute examples and conclude
 $T_X|_C \cong \mathcal{O}(e)^{n-1}$.

For odd $e$, compute examples,
$T_X|_C\cong \mathcal{O}(c-1) \oplus
\mathcal{O}(2)^{3}\oplus \mathcal{O}(-1)$.

\begin{proposition}
\label{proposition-interpolation2}
$n\geq 3$, $m\leq  e+1$.
There is a degree $e$ curve $\phi_e$
interpolating
$m$ general points if and only if there is a
degree $e-2$ curve $\psi_{e-2}$ interpolating
$m-2$ general points in $X$.
\end{proposition}

\noindent
The idea is to do induction.
Suppose you can interpolate $e$ points,
then show you can interpolate $e+2$.

Chose $y_1,\ldots,y_{m-2}$ on the quadric.
Choose a line in $\mathbb{P}^{m-1}$
intersecting the quadric,
and $x_1,\ldots,x_{m-1}$ points on the line.
Join the points on the line to the points on
the quadric,
Then you get somem points of intersection
with the quadric.
Those points and the two points of
intersection of the line and the quadric
let us construct a scroll, which intersects
the quadric.
The intersection of the scroll and the
quadric
(I guess the scroll contains the lines, and
so the points $y_i$)
gives the curve passing through the $y_i$.

\section{Regular but non smooth curves of
genus 3}
\label{section-regular-but-non-smooth}
%% speaker: Cesar Hilario
%% place: Algebra Seminar, IMPA
%% date: October 15, 2025

\noindent
Cesar Hilario, .
Algebra Seminar, IMPA.
October 15, 2025.

\medskip\noindent
{\bf Abstract.} Regular but non-smooth curves
are a unique feature of
geometry in positive characteristic, that
results from the fact that over an
imperfect field the notion of regularity is
weaker than the notion of
smoothness. In the setting of algebraic
geometry over an algebraically closed
field, these curves correspond to fibrations
by singular curves, which are
fibrations of relative dimension 1 whose
fibers are singular. The most famous
examples are arguably the so-called
quasi-elliptic curves and quasi-elliptic
fibrations, which play a key role in the
Bombieri-Mumford classification of
algebraic surfaces in characteristics 2 and
3. In this talk I will discuss the
case of genus 3 in characteristic 2, with an
eye towards the classification of
regular plane projective quartic curves that
become rational after base change.

\medskip\noindent
arXiv 2409.05464.

\medskip\noindent
$C$ will denote a projective geometrically
integral curve over
a not necessaily closed field $K$.
Geometricallt integral
means that $C \otimes_K \overline{K}$ is
integra.
Suppose that $p=\text{char}(K) \geq 0$
and $g=h_1(\mathcal{O}_C)$ is of genus $C$.

We define
\begin{enumerate}
\item $C$ regular iff local rings of $C$ are
regular (DVR).
\item $C$ smooth iff $C$ regular and $K
\otimes_K \overline{K}$ regular.
\end{enumerate}

\noindent
Thus it's obvious that smoothness implies
regularity.
In this talk we explain why the converse is
not true.

Note that if $K$ is perfect (e.g. $p=0$ or
$K=\overline{K}$) implies
 smooth = regular. Recall that $K$ is
 imperfect iff $p>0$
the powers of $K$, $K^p$, is a proper
subfield of $K$
(it's always a subfield but we ask it's
proper).

\medskip\noindent
Assume that $C$ is regular and $p>0$.
The genus of rthe normalization is less than
or equal to the genus of the curve,
i.e.
$$
\tilde{g}=h^1(\mathcal{O}_{\widehat{C
\otimes_K \overline{K}}})\leq g.
$$
We know that $C$ is smooth iff
$\overline{g}=g$.
So $g-\overline{g}$ is a measure of
non-smoothness.

\begin{theorem}[Tate]
\label{theorem-tate}
$\frac{p-1}{2}$ divides $g-\overline{g}$.
\end{theorem}

As a corollary, we get that if $C$ is
non-smooth,
$g-\overline{g}>0$ and so $p \leq  2g+1$.

\begin{example}
\label{example-attain}
There are cases in which the bound is sharp,
namely $C:y^2=x^p+t$, with $t \in K\setminus
K^p$.
Then $g=\frac{p-1}{2}$ and $\bar{g}=0$.
Over $\overline{K}$ rhere exists a singular
point,
$(x,y)=(-t^{1/p},0)$ which
{\bf is not visible over $K$}, that is we get
that
$t^{1/p}\not \in K$.
\end{example}

\begin{example}
\label{example-elliptic}
$\{\text{elliptic curves}\}=\{g =
\overline{g}=1,\text{$\exists K$-rational
point }\}=\{\text{smooth cubic curves with
$K$-rational point}\}$.
$y^2=x^3+4g_i x+27$.
\end{example}

\begin{example}
\label{example-quasielliptic}
$g=1$, $\overline{g}=0$, $p=2,3$,
quasielliptic curves. (Important in the
classification of surfaces
in positive characteristic.)
\end{example}

\medskip\noindent
Now we make a digression to describe
Kodaira-Enriques (over $\mathbb{C}$)
in positive characteristic. Bombieri-Mumford:
{\bf new objects} appear,
quasielliptic surfaces.

Let $k=\overline{k}$, $S$ smooth elliptic
({\it quasielliptic}) surface over $k$.
Suppose we have a fibration
where the general fiber is an elliptic curve
an elliptic curve (resp. a plane
cuspidal cubic) and the generic fiber an
elliptic (resp. quasielliptic)
curve over $k(B)$.

$$
\xymatrix{
S\ar[d]\\
B
}
$$

\medskip\noindent
Back to our main content, let's assume that
$g=3$.
$$
\xymatrix{
\{g=1,\exists K\text{-rat. pt}\}\ar[r]&\ar[l]
\left\{\substack{\text{reg. cubic curves} \\
\text{with $K$-rat pt}}\right\}
}
$$
Recall that genus-degree formula
$$
g=\frac{(d-1)(d-2)}{2}=1,
$$
So then $d=3,4$ gives  $g=3$ (right?).

No we present our main theorem, which is a
classification result:

\begin{theorem}
\label{theorem-classification-quartics}
\begin{itemize}
\item $C$ regular non-hyperelliptic curve
over $K$.
\item  $g=3$, $\overline{g}=0$.
\item $p=2$.
\end{itemize}
(Not that non-hyperllipticity of $C$ and
$g=3$ imply that $C$ is a quartic
over $K$, and that  $\overline{g}=0$ implies
$C$ is geometrically rational.)
Then $C$ is isomorphic to one of the
following quartics.

\begin{enumerate}
\item $y^4+az^4+xz^3+bx^2z^2+cx^4=0$ where
$a,b,c \in K$, $c \not \in K^2$.
(Notice the imperfectness of $K$ is crucial.)

\item $y^4+az^4+bx^2y^2+cx^2z^2+bx^3z+dx^4=0$
for $a,b,c,d, \in K$,
$a \not \in K^2$, $b\not =0$.

\item
\item
\item
\end{enumerate}
\end{theorem}

\noindent
Now we enumerate properties of these
families.

\begin{itemize}
\item $C$ is a purely inseparable double
cover of a quasielliptic curve.
To understand this consider the induced map
from
Frobenius map $K \to K$, $a \mapsto a^p$ and
do base change:
$$
\xymatrix{
C^{(p)}\ar[r]\ar[d]\ar@{}[dr]|-{\lrcorner}&
C\ar[d]\\
\Spec K\ar[r]^{\substack{\text{induced from}
\\ \text{Frobenius}}}&\Spec K
}
$$
This gives a map $C \to C^{(p)}$ by universal
property of pullback.

$$
\xymatrix{
C \ar[rrrd]^{\text{Frobenius}}
\ar@{-->}[rrd]\ar[rrdd] & & \\
& &  C^{(p)}\ar[d] \ar[r] & C \ar[d] \\
& &  \Spec K\ar[r] & \Spec K
}
$$

Then consider the normalization of
$C^{(p)}=C_1$ which
turns out to be quasielliptic; by universal
property of normalization we can lift:
$$
\xymatrix{
& \widehat{C^{(p)}}\ar[d]\\
C\ar@{.>}[ur]\ar[r]&C^{(p)}.
}
$$
Here the lower arrow is purely inseparable of
degree 2.

\item There exists a unique non smooth point
$x \in C$.

\item Now we iterate the construction of the
first item:
$$
\xymatrix{
\underbrace{C}_{g=3}\ar[r]&\underbrace{C_1}
_{g=1}\ar[r]&
\underbrace{C_2}_{g_2-9}\ar[r]&
\underbrace{C_3}_{g_3=0}\ar[r]&
\underbrace{C_4}_{g_4=0}\ar[r]&\cdots
}
$$
Let $x$ be the non-smooth (non-rational)
point of $C$. Map it under
these maps to $x_1$( non-smooth,rational),
then  $x_2$ (smooth, rational?),
then $x_3 $ (smooth, rational.), then $x_4$
(smooth, rational), etc.

We have
$$
\begin{matrix}
&\text{$x_2$ rational}& \text{$x$ canoical
divisor} &
E=K(c_2)\\
(i)& Y & Y & Y\\
(ii) & N & Y & N\\
(iii) & Y & N & N\\
(iv) & N & N & Y\\
(v) & N & N & N
\end{matrix}
$$
Where the last column is explained next.

\item Canonical field of $C$ = subfield of
$K(C)$ generated
by the quotients of all non-zero holomorphic
differentials
=  $K(C)$ ($C$ non-hyperelliptic).

\item Pseudocanonical field of $C$ = …
non-zero \underline{exact}
holomorphic differential := $E$. So one has
$[K(c),E]=4=p^2$.

$$
\xymatrix{
&K(C)\ar[ld]_{4}\ar[rd]^{4}\\
E&&K(C_2).
}
$$
\end{itemize}

\noindent
In conclusion, we manage to characterize
these families.

\begin{example}
\label{example-chain-of-maps}
In case 1, let $a=b=0$. We get
\begin{equation}
\label{equation-curve-pencil}
C: y^4+xz^3+\underbrace{c}_{\in K\setminus
K^2}x^4=0
\end{equation}
$$
\underbrace{C_1}_{\text{quasiell.}}:xy^2+z^3+
cx^3=0.
$$
Further
$$
\xymatrix{
C\ar[r]&C_1\ar[r]&C_2 \cong\mathbb{P}^1\ar[r]
&C_3 \cong \mathbb{P}^1\ar[r]&C_4\cong
\mathbb{P}^1\ar[r]&\cdots.
}
$$
\end{example}

\medskip\noindent
With $C$ one can construct a pencil
(fibration over $\mathbb{P}^1$)
of quartics.
Let $k=\overline{k}$ (new base field from
what we fixed at the beginning!).
Idea: replace $c$ for $t$ Equation
\ref{equation-curve-pencil}.
 $K=k(\mathbb{P}^1)$,
 $S=V(t_0(y^4+xz^3)+t_1x^4)
\subseteq \mathbb{P}^2_{(x:y:z)}\times
\mathbb{P}_{(t_0:t_1)}$.
This has a singular point, so we blow up:
$$
\xymatrix{
\tilde{S}\ar[d]\ar@/^1em/[dd]^f\\
S\ar[d]\\\mathbb{P}^1
}
$$
The generic fiber is the curve
$y^4+xz^3+tx^4=0$,
so that $t_1/t_0 \in k(\mathbb{P}^1)=K$.
Fibers are plane rational quartics,
and the singular fiber $f_1(0:1)$ is a
configuration of lines
given by a Dynkin diagram
$$
\xymatrix{
E_1 \ar@{-}[r]& E_2\ar@{-}[r]
& E_3 \ar@{-}[r]& E_4\ar@{-}[d]\ar@{-}[r]
& \cdots \ar@{-}[r]& E_{15}\\
& & & E
}
$$
We do the same with $C_1$: define
$S^1=V(t_0(xy^2+z^3)+tx^3) \subseteq
\mathbb{P}^2 \times \mathbb{P}^1$.
Then $S_1$ is also singular, and we blow up
to obtain a smooth
quasielliptic fibration ($\tilde{S^1}$ is
quasielliptic).
The fibers are plane cuspidal cubics.
The singular fiber an arrangement of curves
$$
\xymatrix{
F_1 \ar@{-}[r]& F_2\ar@{-}[r]&
F_3\ar@{-}[d]\ar@{-}[r]& \cdots \ar@{-}[r]&
F_8\\
& & F
}
$$

In fact, there is a correspondence between
the
singular fibers of the two fibrations.

\section{Variations on a theorem of Baum and
Bott}
\label{section-baum-and-bott}
%% speaker: Stéphane Druel, Université Claude
%% speaker: Bernard Lyon
%% place: Algebra Seminar, IMPA
%% date: November 5, 2025

\noindent
Stéphane Druel, Université Claude Bernard
Lyon.
Algebra Seminar, IMPA.
November 5, 2025.

\medskip\noindent
{\bf Abstract.}
In this talk, I will discuss an analogue of
the Baum-Bott vansihing theorem for
the tangent bundle of the space of leaves of
regular foliations together with
applications to transversely projectively
flat foliations of dimension one. This
is work in progress.

\medskip\noindent



Baum-Bott 70'.

Let $E$ be a vector bundle on $X$
complex manifold is a sheaf of 1-jets of $E$:
$$
J'_X(E)= E \oplus \Omega^1_X \otimes E
$$
where equality is the category of sheavs of
abelian groups.
Where the $\mathcal{O}_X$ structure is given
by
$$
\underbrace{f}_{\substack{\text{local
function} \\ \text{on }U}}
(\underbrace{e}_{\substack{\text{local
section} \\ \text{of $E$ on  $U$}}}
, \underbrace{\alpha}_{\substack{\text{local
section} \\
\text{of $\Omega^1_X \otimes E$}}})
=(fe,f\alpha-df \otimes e)
$$
The {\it Atiyah class} of $E$ is the element
$$
\text{at}_X(E) \in H^1(X,\Omega_X^1 \otimes
\text{End}(E))
$$
corresponding to the exact sequence of
$\mathcal{O}_X$-modules
$$
\xymatrix{
0\ar[r]&\Omega_X^1\otimes
E\ar[r]&J_X^1(E)\ar[r]&E\ar[r]&0.
}
$$

\begin{remark}
\label{remark-holomorphic-connection}
 $\text{at}(E)=0$ if and only if there exists
$\nabla:E \to \Omega^1_X \otimes E$ map of
abelian sheaves
such that $e \mapsto  (e-\nabla e)$ is
$\mathcal{O}_X$-linear
$$
-\nabla fe=-f\nabla e - df \otimes e
$$
if and only if $E$ admits a holomorphic
connextion.
\end{remark}

\begin{example}
\label{example-line-bundl-holom-conne}
If $E$ is a line bundle, at
$(E)=(d\text{log}f_{ij})_{ij}
\in H^1(X,\Omega_X^1)$, $E=(f_{ij})$.
\end{example}

\begin{theorem}[Atiyah, '57]
\label{theorem-atiyah-57}
Let $r =\text{rank}E$, $k$ an integer, $X$ be
compact Kähler
and $\Sigma_k$ a polynomial of degree $k$ on
$M_r$
such that  $\Sigma_k(g)$ is the $k$-th
elementary
symmetric function on the characteristic
zeros of
$g \in M_r$. $S_k=$ polynomial given by
polarization.
$$
\begin{aligned}
M_r^{\otimes k}&\to \mathbb{C}\\
\text{End}E^{\otimes k}& \to \mathcal{O}_X\\
\text{Ex., $k=1$}\qquad
\text{End}E&\xrightarrow{\text{Tr}}
\mathcal{O}_X.
\end{aligned}
$$

$$
\xymatrix{
H^1(X,\Omega_X^1 \otimes
\text{End}(E))^{\otimes k}\ar[r]^{\cup }
&H^k((\Omega^1_X)^{\otimes k}\otimes
\text{End}E^{\otimes k}\ar[r]
&H^k(X,\Omega_X^k)\ar[r]
&H^{2k}(X,\mathbb{C})\\
\text{at}(E)^{\otimes k}\ar[rrr]
&&&(2\pi i)^kc_{2k}(E).
}
$$
(Double check the term $(\Omega^1_X)^{\otimes
k}$…
I think there's a typo.)

Partial connection. $\mathcal{F}$ regular
distribution.
$\mathcal{F} \subseteq T_X$ subbundle.
An {\it $\mathcal{F}$-connection} on $E$ is
$\nabla:E \to \mathcal{F}^\vee \otimes E$
such that
$$
\nabla(fe)=f\nabla e+f_{\mathcal{F}}f \otimes
e
$$
where
$d_{\mathcal{F}}:\mathcal{O}_X\xrightarrow{d}
\Omega_X^1
\to\mathcal{F}^\vee$.
\end{theorem}

\begin{proposition}[Bott-Baum, '70]
\label{proposition-bott-baum}
If $E$ admits an $\mathcal{F}$-connection
then all characteristic
classes of degree  $>\text{codim}\mathcal{F}$
vanish.
\end{proposition}

\begin{proof}
There exists an $\mathcal{F}$-connection if
and only if
$$
\xymatrix{
0\ar[r]&\Omega_X^1 \otimes
X\ar[r]\ar[d]&J_X^1\ar[r]&E\ar@{=}[d]\ar[r]&
0\\
0\ar[r]&\mathcal{F}^\vee\otimes E\ar[r]&
\mathcal{F}^\vee \otimes E \oplus E\ar[r]&
E\ar[r]&0
}
$$
if and only if $\text{at}(E) \in
H^1(X,\Omega_X^1\otimes\text{End}E)$
maps to 0 in
$H^1(X,\mathcal{F}^\vee\otimes\text{End}E)$
is and only if $\text{at}(E)$ comes from
$$
H^1(X,N^\vee\otimes\text{End}E)\to
H^1(X,\Omega_X^1\otimes\text{End}E)
$$
where $N=T_{X/\mathcal{F}}$.
\end{proof}

\begin{example}[Bott's partial connection]
\label{example-bott-partial-connection}
Suppose $\mathcal{F}$ is a regular foliation
and $N=T_{X/\mathcal{F}}$ has a
$\mathcal{F}$-connection.
Take $n$ a local section of $N$,  $v$ a local
section of $\mathcal{F}$
(over the same open) and suppose exists a
section $t$ of $T_X$
that is a lift of $t$, i.e. $p(t)=n$ for
$p:T_X \to N$.
Then $\nabla_v^Bn=p([v,t])$ is a section of
$N$.
This defines an $\mathcal{F}$-connection on
$N$.
\end{example}

\begin{lemma}[Beauville]
\label{lemma-beauville}
Let $X$ be a complex manifold and $E$ is a
direct submanifold of $T_X$.
Then $\text{at}(E) \in H^1(X,\Omega_X^1
\otimes \text{End}E)$
comes from $H^1(X,E^\vee \otimes
\text{End}(E))$.
\end{lemma}

\begin{proof}
Consider
$$
T_X= E \oplus F \xrightarrow{p}E.
$$
We claim that there exists an
$\mathcal{F}$-connection on $E$.
Let $v$ be a local section of $F$,
$e$ a local section of $E$.
Define $\nabla_ve=p([v,e])$.
One checks that this satisfies the required
properties.
\end{proof}

\noindent
Notice that no involutivity is imposed on $E$
- it's just a distribution.

\medskip\noindent
Now we will discuss some variations of
Baum-Bott's theorem
\ref{theorem-baum-bott}.

Let $\mathcal{F}$ be  regular foliation, $E$
a vector bundle on $X$ equipped with a flat
$\mathcal{F}$-connection.
For example, $\nabla^B$ flat.

Recall that  $(E,\nabla)$ can be thought of
as
$\text{Mon}(\mathcal{F})$-equivariant vector
bundles
where $\text{Mon}(\mathcal{F})$ is the
monodromy groupoid.

Let $X/\mathcal{F}$ be the space of leaves,
and
$\mathcal{O}_{X/\mathcal{F}}=\mathcal{O}
_X^{d_{\mathcal{F}}}=
\{\text{functions $f$ such that
$d_{\mathcal{F}}f=0$}\}$.

It is a fact that $E^{\nabla}$ is a locally
free $\mathcal{O}_{X/\mathcal{F}}$
sheaf of rank the rank of $E$ and such that
$E=E^{\nabla}\otimes_{F_{X/\mathcal{F}}}
\mathcal{O}_X$.

Now one defines bundles of 1-forms and
tangent bundle as follows.
 $\Omega_{X/\mathcal{F}}^1=(N_{\mathcal{F}}^*
 )^{\nabla B}$
and
$T_{X/\mathcal{F}}=N_{\mathcal{F}}^{\nabla
B}$,
which are locally free sheaves of
$\mathcal{O}_{X/\mathcal{F}}$-modules
of rank $q=\text{codim}\mathcal{F}$.

The Atiyah class of $(E,\nabla)$ is
$$
J^1_{X/\mathcal{F}}(E,\nabla)=E^{\nabla}
\oplus
\Omega_{X/\mathcal{F}}^1
\otimes_{\mathcal{O}_{X/\mathcal{F}}}
E^{\nabla}
$$
with $\mathcal{O}_{X/\mathcal{F}}$-module
structure given by
 $$
(e,\alpha)=(f\alpha,f\alpha-\underbrace{df}
_{\in
\Omega^1_{X/\mathcal{F}}}
\otimes e.
$$
Then
$$
\text{at}_{X/\mathcal{F}}(E,\nabla)
\in
H^1(X,\Omega_{X/\mathcal{F}}
^1\otimes_{\mathcal{O}_{X/\mathcal{F}}}
\text{End}_{\mathcal{O}_{X/\mathcal{F}}}
E^{\nabla})
$$
corresponding to
$$
\xymatrix{
0\ar[r]&\Omega_{X/\mathcal{F}}
^1\otimes_{\mathcal{O}_{X/\mathcal{F}}}\ar[r]
&J_{X/\mathcal{F}}^1(E)\ar[r]&E^{\nabla}
\ar[r]&0.
}
$$

\begin{lemma}
\label{lemma-atiya-class-maps-atiya-class}
$\text{at}_{X/\mathcal{F}}(E,\nabla)$ maps to
$\text{at}_X(E)$ under
$$
H^1(X,\Omega_{X/\mathcal{F}}
^1\otimes_{\mathcal{O}_{X/\mathcal{F}}}
\text{End}_{\mathcal{O}_{X/\mathcal{F}}}
E^{\nabla}
\to H^1(X,\Omega_X \otimes \text{End}E).
$$
\end{lemma}

\noindent
Partial connections. Let $G \subseteq
T_{X/\mathcal{F}}$
be a subbundle a {\it $G$-connection}
$(E,\nabla)$ is
$$
\begin{aligned}
D: E^{\nabla} &\longrightarrow G^\vee\otimes
E^{\nabla} \\
\text{s.t.} D(fe)&=fDe+d_Gf\otimes e
\end{aligned}
$$
where
$d_G:\mathcal{O}_{X/\mathcal{F}}\to\Omega_{X/
\mathcal{F}}^1\to
G^\vee$.

It's a fact that a $G$-connection
$(E,\nabla)$ exists if and only if
$\text{at}_{X/\mathcal{F}}(E,\nabla)$ maps to
0 in
$H^1(X,G^\vee
\otimes_{G_{X/\mathcal{F}}}\text{End}
_{\mathcal{O}_{X/\mathcal{F}}}
E^{\nabla}$.

The analogue of Beauville's lemma
\ref{lemma-beauville} is

\begin{lemma}[Beauville]
\label{lemma-analogue-beauville}
Suppose $N_{\mathcal{F}}=\mathcal{A} \oplus
B$ such that
$\nabla^B\mathcal{A} \subseteq
\mathcal{F}^\vee
\otimes_{\mathcal{O}_X}\mathcal{A}$
and $\nabla^B \mathcal{B} \subseteq
\mathcal{F}^\vee\otimes_{\mathcal{O}_X}
\mathcal{B}$.
Let $\nabla_{\mathcal{A}}$ and
$\nabla_{\mathcal{B}}$ be a flat
 $\mathcal{F}$-connection on $\mathcal{A}$
 and $\mathcal{B}$.

Then
$\text{at}_{X/\mathcal{F}}(\mathcal{A},
\nabla_{\mathcal{A}})
\in
H^1(X,\Omega_{X/\mathcal{F}}
^1\otimes_{\mathcal{O}_{X/\mathcal{F}}}
\text{End}_{\mathcal{O}_{X/\mathcal{F}}}
\mathcal{A}^{\nabla_{\mathcal{A}}}$
comes from
$H^1(X,(\mathcal{A}^{\nabla_{\mathcal{A}}})
^\vee
\otimes_{X/\mathcal{F}}\text{End}
_{\mathcal{O}_{X/\mathcal{F}}}\mathcal{A}
^{\nabla_{\mathcal{A}}}$.
\end{lemma}

\begin{proof}
Same proof but we need a Lie bracket on
$T_{X/\mathcal{F}}=N_{\mathcal{F}}\subseteq
N_{\mathcal{F}}\overset{p}{\leftarrow}T_X$.
Let $v_1,v_2$ be local sections on
$T_{X/\mathcal{F}}$
with $v_1=p(t_1)$ and $v_2=p(t_2)$.
Define
$$
[v_1,v_2]=p([t_1,t_2]).
$$
One can check that this a Lie bracket on
$T_{X/\mathcal{F}}$.
Suppose that $t_1 \in \mathcal{F}$. Then
$$
p([t_1,t_2])\overset{\text{def}}{=}
\nabla_{t_1}^{\mathcal{B}}v_2=0.
$$
For $v_3 \in \mathcal{F}$,
$$
\nabla_{v_3}^{\mathcal{B}}[v_1,v_2]
=p([v_3,[t_1,t_2]])=\pm\nabla_{v_3}
^{\mathcal{B}}v_1\pm
\nabla_{v_3}^{\mathcal{B}}v_2=0.
$$
\end{proof}

\begin{lemma}
\label{lemma-vanishing-chern-class}
Suppose that $\mathcal{F}$ is a regular
foliation on $X$.
Let $N_{\mathcal{F}}=\mathcal{L}_1 \oplus
\ldots \oplus \mathcal{L}_q$
be line bundles on $X$,
$\nabla^{\mathcal{B}}\mathcal{L}_i\subseteq
\mathcal{F}^\vee\otimes\mathcal{L}
_i\rightsquigarrow
\nabla_i$.
Suppose
$(\mathcal{L}_i,\nabla_i)\cong(\mathcal{L}_j,
\nabla_j)$
for all $i,j$.
Then $c_1(\mathcal{L}_i)=0$.
\end{lemma}

\begin{proof}
By the Comparison Lemma,
$$
\begin{aligned}
H^1(X,\Omega_{X/\mathcal{F}}^1
&\longrightarrow H^1(X,\Omega_X^1) \\
\text{at}_{X/\mathcal{F}}(\mathcal{L},
\nabla_i)
&\longmapsto
\text{at}(\mathcal{L}_i),
\end{aligned}
$$

\noindent
it suffices to prove that
$\text{at}_{X/\mathcal{F}}(\mathcal{L}_i,
\nabla_i)$
vanishes.
$$
\text{at}_{X/\mathcal{F}}(\mathcal{L}_i,
\nabla_i)
\in H^1(X,\Omega_{X/\mathcal{F}}^1)
=\bigoplus_{j=1}^9
H^1(X,(\mathcal{L}_j^{\nabla_j})^\vee)
$$
comes from
$H^1(X,(\mathcal{L}_i^{\nabla_i})^\vee)$.
\end{proof}

\begin{example}
\label{example-degree-2-line-bundle-on-e}
$A=E^1 \times B^{n-1}$, an abelian variety.
Let $ L$ be the pullback of a degree 2
divisor on $E$.
Then there exists  $T_X
\to\!\!\!\!\!\to L^{\oplus n-1}$.
\end{example}

\begin{proposition}
\label{proposition-}
$\mathcal{F}$ regular foliation of dimension
1 on $X$ complex projective
of dimension $\geq 3$.
Suppose $\mathcal{F}$ is transversally
projective flat
and that $X$ contains a rational curve.
Then $\mathcal{F}$ is a $\mathbb{P}^1$-bundle
and it is tangent
to $\mathcal{F}$.
\end{proposition}

\medskip\noindent
Now we explain what it means to be
projectively flat.
It is when the projectivization of the normal
bundle $\mathbb{P}(N)$
is induced by a representation of $\pi_1(X)
\to \text{PGL}_g$
it comes with an Ehresmann connection on
$$
\xymatrix{
\mathbb{P}(N)=\mathbb{P}\ar[d]^{f}\\
X
}
$$
$T_{\mathbb{P}}=T_f \oplus g$.

Now, it is also called transversally
projectively flat
if $\nabla^B$ is a partial connection on $N$
and there exists a foliation $\mathcal{H}$ on
$$
\xymatrix{
\mathbb{P}(N)\ar[d]\\ X
}
$$
such that $\mathcal{H}$ projects onto
$\mathcal{F}$.
Transversality means that $\mathcal{H}
\subseteq \mathcal{G}$.

Finally, transversally projectively flat
means that
it is induced on a
$\text{Mon}(\mathcal{F})$-module
by a representation of
$\pi_1(\text{Mon}(\mathcal{F}) \to
\text{PGL}_g$.

\section{Singular schemes of codimension one
distributions on projective space}
\label{section-muniz-distributions}
%% speaker: Alan Muniz
%% place: OPEGA 2026, UFPE
%% date: February 2026

\noindent
Alan Muniz.
OPEGA 2026, UFPE.
February 2026.

\medskip\noindent
No abstract or title available; the title
above is descriptive.

\medskip\noindent
{\bf Sources of these notes.} Daniel's short
LaTeX notes typed during the talk,
transcribed here in October 2026. There are
no slides or recording. The original notes
give the speaker only as ``Alan''; the
surname comes from the two papers cited in
the talk, in which the speaker is a coauthor.
The date is the month of the file date of
the original notes. These notes are
fragmentary; gaps are marked as such. The
minicourses of OPEGA 2026 have their own
notebooks.

\medskip\noindent
{\bf Distributions and their singular
schemes.} A $1$-form on $\mathbb{P}^n$
written in homogeneous coordinates as
$$
\omega=\sum_{i=0}^n A_i\,dx_i,
$$
with the $A_i$ homogeneous polynomials of
degree $d+1$, defines a codimension one
distribution $\mathcal{F}$ of degree $d$ by
$\ker\omega$. Its \emph{singular scheme} is
$$
Z=V(A_0,\ldots,A_n).
$$
The projective space
$\mathbb{P}H^0(\Omega^1_{\mathbb{P}^n}(d+2))$
of such forms is stratified into flat
families according to the Hilbert polynomial
$P_Z(t)$ of $Z$:
$$
\mathbb{P}H^0(\Omega^1_{\mathbb{P}^n}(d+2))
=\coprod_i S_i.
$$
The strata were described for $d=1$ by
Corrêa, Jardim and Muniz (Math. Z., 2022)
and for $d=2$ by Galeano, Jardim and Muniz
(J. Pure Appl. Algebra, 2022). The original
notes write ``Galeane''; references as
recorded, not checked.

\medskip\noindent
{\bf Question 1.} Given a distribution of
degree $d$, can its singular scheme $Z$ be
contained in a hypersurface $V$ of degree at
most $d$?

There is a map $\pi\colon S\to
\operatorname{Hilb}$ from a stratum to the
Hilbert scheme, sending a distribution to
its singular scheme.

Daniel's question in the original notes: is
$S$ the subvariety of
$\mathbb{P}H^0(\Omega^1(d+1))$ cut out by
the integrability condition
$\omega\wedge d\omega=0$? (The twist $d+1$
here differs from the $d+2$ above; one of
them is a slip in the original notes.)

\medskip\noindent
{\bf A diagram.} With
$V=\langle dx_0,\ldots,dx_n\rangle$, the
talk showed maps
$$
\Lambda^2V\otimes H^0(\mathcal{I}_Z(d))
\longrightarrow
H^0(\Omega^1_{\mathbb{P}^n}(d+2)
\otimes\mathcal{I}_Z)
\longrightarrow
\operatorname{Tor}_1(I_Z,\mathbb{C})_{d+2},
$$
the second one surjective. The original
notes write $\mathcal{I}_2$ for the ideal
sheaf, and the formula given on elements is
illegible; I do not reproduce it.

\medskip\noindent
{\bf Question 2.} If $\mathcal{F}$ is a
foliation, is $\deg Z_2\geq d+1$, where
$Z_2$ is the part of $Z$ of codimension
$2$? (The codimension is marked as
uncertain in the original notes.)



\section{Um passeio nos trabalhos de
Gerd Faltings}
\label{section-faltings}
%% speaker: Hossein Movasati, IMPA
%% place: Special talk (palestra especial)
%% date: May 19, 2026

\noindent
Hossein Movasati, IMPA.
Pelestra especial,
May 19, 2026.

\medskip\noindent
{\bf Abstract.}
Nesta palestra vamos dar um passeio nos
trabalhos de G. Faltings, começando por sua
demonstração da conjetura de Mordell em
1983, que lhe trouxe a medalha Fields. Uma
boa parte de sua carreira foi elaborar o
lado aritmético de um dicionário heurístico
proposto por P. Deligne em 1970, que tem
origem na teoria conjetural de motivos de
Grothendieck. Isso inclui versão p-ádica de
teoria de Hodge, domínios de períodos,
correspondência de Simpson, etc. Várias
propriedades aritméticas de variedades
abelianas são outro foco da carreira dele
que vamos abordar nesta palestra. Em 2026
ele ganhou o prêmio Abel por suas
contribuições fundamentais em geometria
aritmética.

\medskip\noindent
{\bf Upshot.} 
Elliptic curves are diophantine equations.
Some are simple and fun to solve,
some are extremely complicated ---
like Fermat's last theorem.
Finding solutions to a diophantine
equation is like finding $k$-points
on an algebraic curve for suitable
choices of fields.
In this sense,
Mordell's conjecture, now Faltings'
theorem is a statement of finiteness
of solutions of diophantine equations.
Other results are adaptations
of results from complex geometry
to finite fields.

\medskip\noindent
{\bf Notes.}

Finding roots of a polynomial $f \in
\mathbb{Q}[x,y]$ becomes the geometric
problem of finding the $\mathbb{Q}$-rational
points of $\operatorname{Spec}(\mathbb{Q}[x,y]/(f))$.
Or, one can pass to the projective version
$\operatorname{Proj}(\mathbb{Q}[x,y,z]/(F))$,
where $F$ is the homogenization of $f$.
We write
$C(\mathbb{Q})=\{[x:y:z] \in
\mathbb{P}^2(\mathbb{Q}) \mid
F(x,y,z)=0\}$
and similarly $C(\mathbb{C})$ for the
analogous set.
 
It turns out that the topology
of the curve, via its genus, strongly
constrains the $\mathbb{Q}$-rational points.
The easiest cases are:
\begin{itemize}
\item ($g=0$.) Here $d=1,2$;
then the curve is projectively equivalent to
a conic, and over $\mathbb{Q}$ it has either
no rational points or infinitely many.

\item ($g=1$.) Here $d=3$; this is more difficult
but still well-studied; see the
Birch--Swinnerton-Dyer conjecture. 

\item ($g\geq 2$.) Here $d\geq 4$.
This is much more complicated!
Faltings'
theorem implies that $C(\mathbb{Q})$ is
finite.
\end{itemize}

\noindent
This is the former Mordell conjecture
from 1962.

\begin{theorem}[Faltings, 1983]
For any smooth curve of genus $\geq 2$ over
$\mathbb{Q}$, we have $\#C(\mathbb{Q})<
\infty$.
\end{theorem}

The same theorem remains true after
replacing $\mathbb{Q}$ by a number field,
that is, a finite field extension of
$\mathbb{Q}$.

\begin{theorem}[Manin 1963, Grauer 1965]
\label{theorem-manin-grauer}
For any smooth non isotrivial curve of
genus $\geq 2$ over $k=\mathbb{C}(t)$
we have
$$
\# C(k) < \infty.
$$
\end{theorem}

\begin{proof}
Using Gauss-Manin connection ---
the name of which was invented by 
Grothendieck.
\end{proof}

\noindent
The main difference between
$\mathbb{C}(t)$
and $\mathbb{Q}$
is that in $\mathbb{C}(t)$
we have a derivation
$$
\frac{\partial }{\partial t}:
\mathbb{C}(t) \to \mathbb{C}(t)
$$
satisfying the Leibniz identity.

\medskip\noindent
The proof is somehow related
to Shafarevich and Tate conjectures.
The former states there's only
finitely many isomorphism
classes of abelian varieties of a
given dimension with
fixed polarization degree, defined
over a field with
good reduction outside a finite set
of places.

Tate conjecture was proved for finite
fields by Tate in 1966,
and by Faltings for number fields
in 1983.
It relates maps between the tangent
spaces of two abelian varieties
and maps between the varieties
themselves:
$$
\Hom_K(A,B) \otimes \mathbb{Z}_p
\cong
\Hom_{Gal(\tilde{K}/K)}
(T_pA,T_pB)
$$
where $K$ is a field of characteristic
$\neq p$, and $A,B$ are abelian over $K$.

Further, the proof uses Arakelov's
theory of arithmetic genus
and Faltings height.

\medskip\noindent
{\bf Abelian varieties.} 

Any curve (over any field)
can be embedded in some variety
called the Jacobian variety.
Over $\mathbb{C}$,
abelian varieties are just tori:
$$
\mathbb{C}^n/\Lambda
\cong_{C^\infty}
S^1\times… S^1,\qquad  2n \text{ times}.
$$

An abelian variety of dimension one
over a field $k$ is called
an {\it elliptic curve}.
They are given by diophantine
equations of degree 3 in two variables:
$$
y^2+a_1xy+a_3y
=
x^3+a_2x^2+a_4x,
\qquad a_1,…,a_6 \in k.
$$

The name was invented by A. Weil.
Before him, many mathematicians like E.
Picard,
G. Humbert, were using ``hyperelliptic
surface'' for Abelian surfaces.

\medskip\noindent
{\bf Shavarevich's conjecture.} 
What is the number of non-isomorphic elliptic
curves over $\mathbb{Q}$ and with bad
reduction only for primes $2,3$?

\medskip\noindent
A consequence of Tate conjecture
and Chebotarev density theorem is:;

\begin{theorem}
\label{theorem-tate-chebot-consequence}
LEt $E_i$, $i=1,2$ be two
elliptic curves over $\mathbb{Q}$.
TFAE:
\begin{enumerate}
\item $E_1$ is isogeneous to $E_2$
over $\mathbb{Q}$.
\item For all except a finite number
of primes,
$$
\# E_1(\mathbb{F}_p)
=\# E_2(\mathbb{F}_p).
$$
\end{enumerate}
\end{theorem}

\begin{definition}
\label{definition-isogeny}
{\it AI-generated.}
An isogeny between elliptic curves over
$\mathbb{Q}$ is a nonzero morphism of
elliptic curves defined over $\mathbb{Q}$
that is also a group homomorphism and has
finite kernel.
\end{definition}

Here isogenous means that there exists an
isogeny $f:E_1 \to E_2$ over $\mathbb{Q}$.

\medskip\noindent
One of the key ingredients
in Falting's proof
is the notion of {\it Faltings height},
which are integrals used to
control the numerator and denominator
or rational solutions of a (diophantine?)
equation.

For example, for $r = n/m$,
$h(r)=n^2+m^2$. For all $M \in \mathbb{R}$,
$$
\{r \in \mathbb{Q}| h(r)<M\}
$$
is finite.

\begin{proposition}[Main property
of height functions]
\label{proposition-faltings-height}
For any positive integer $M$,
the number of elliptic curves $E/\mathbb{Q}$
with $h_{\text{Falt}}<M$ is finite.
\end{proposition}

\noindent
On the curve given by Equation
\ref{equation-?}
consider the 1-form
$$
\omega=\frac{dx}{2y+a_1x+a_3}.
$$
The Falting's height is basically…
[missing! an integral involving $\omega$…]

\medskip\noindent
\begin{itemize}
\item Elliptic integrals $\leq 1900$.
\item Poincaré, Pircard, $\sim 1900$.
\item Singular, de Rham 1900-1950.
\item Weil: cohomology over the integers
(enter arithmetic).
\item Grothendieck: étale, crystalline, etc.
\end{itemize}

\noindent
Then is Deligne's dictionary,
in a 1970 paper, where $\ell$-adic
cohomology is introduced
(I think --- abstract in french, check).

\medskip\noindent
From complex to arithmetic geometry:

\begin{itemize}
\item Endilchkeitssätze
fur abelsche Varietäten über Zahlkörpern.
\item 1984: Calculus on arithmetic surfaces.
\item 1988: p-adic Hodge theory.
\item 1990: Degeneration of abelian
varieties (with Chai Ching-li).
\item 1992: Lectures on the
arithmetic Riemann-Roch theorem.
\item 1999: Does there exist an arithmetic
Kodaira-Spencer class? (!!!)
\item 2005: A p-adic Simpson corresponding.
\item 2020: covering of p-adic period
domains.
\end{itemize}

\medskip\noindent
{\bf p-adic Hodge theory.}

Grothendieck and Deligne
noticed that we can obtain
the Hodge decomposition of cohomology
using filtrations, which sit
in the realm of polynomials ---
no harmonic analysis needed.

The idea of Deligne (no, Faltings?) 
was to use
this proof for the $p$-adic case.

\section{Grothendieck--Teichm\"uller theory: from arithmetic geometry to arithmetic topology}
\label{section-grothendieck-teichmuller-arithmetic-topology}
%% speaker: Marco Boggi
%% place: Sala 228
%% date: May 27, 2026, 15:30

\noindent
Marco Boggi.
Qua 27 mai 2026, 15:30.
SALA 228.

\medskip\noindent
{\bf Abstract.}
Grothendieck--Teichm\"uller theory originated
in arithmetic geometry through the study of
Galois actions on geometric \'etale
fundamental groups. In this talk, we argue
that arithmetic topology provides a natural
setting for applications of
Grothendieck--Teichm\"uller theory. More
specifically, we explain how our theory of
automorphisms of profinite mapping class
groups may lead to new results on profinite
rigidity questions for mapping class groups.

\medskip\noindent
{\bf Warning.}
These are live notes from the talk,
written with unusually heavy AI assistance
since that day I forgot my glasses.
Some theorem
statements may be imprecise; the goal is to
record the ideas and references from the day.\medskip\noindent

\medskip\noindent
The main upshot is that GT theory allows us
to pass the problem to a free group in three
generators, and that the talk aims to
contrast arithmetic and topological
properties.

A more personally significant 
upshot, explained next,
is that the étale fundamental
group of the moduli space of curves
is the profinite completion of
the fundamental group of any curve in 
the moduli, which makes sense since the
underlying topological space of all curves
in the moduli is the same.

\medskip\noindent
Historically, this is very much in the spirit
of Grothendieck's 1984 \emph{Esquisse d'un
programme}. Drinfeld then introduced the
Grothendieck--Teichm\"uller group in 1991,
making precise a framework that was already
expected to have arithmetic applications.

\medskip\noindent
The point is that $\mathcal M_{g,n}$ is
not just an abstract moduli space: it is
built from one fixed topological surface
$S_{g,n}$. A point of $\mathcal M_{g,n}$
is essentially a complex structure on this
same underlying surface.

\medskip\noindent
Thus the topological fundamental group of
the moduli space records how a family of
curves can move around and come back to the
same curve with its underlying surface
changed by a diffeomorphism. This gives the
mapping class group:
$$
\pi_1^{\mathrm{top}}
\bigl(\mathcal M_{g,n}(\mathbb C)\bigr)
\cong
\Gamma_{g,n}.
$$
Here
$$
\Gamma_{g,n}
=
\operatorname{Diff}^+(S_{g,n})/
\operatorname{Diff}_0(S_{g,n}).
$$

\medskip\noindent
The \'etale fundamental group is the
algebraic-geometric version of the
fundamental group. Instead of remembering
all covering spaces, it remembers finite
algebraic covers. Over $\mathbb C$, this
amounts to taking the profinite completion
of the usual fundamental group. Therefore
one gets
$$
\pi_1^{\acute{e}t}
\bigl(\mathcal M_{g,n}
\otimes_{\mathbb Q}\overline{\mathbb Q}\bigr)
\cong
\widehat{\Gamma}_{g,n}.
$$

\medskip\noindent
So the profinite mapping class group is not
an artificial object: it is literally the
\'etale fundamental group of the moduli
space of curves. This is why
Grothendieck--Teichm\"uller theory naturally
acts on profinite mapping class groups.

\medskip\noindent
The starting point is the \'etale
fundamental group of the projective line with
three points removed. Over $\mathbb{C}$, the
space
$$
\mathbb{P}^1_{\mathbb{C}} \setminus \{0,1,\infty\}
$$
is a sphere with three punctures, so its
topological fundamental group is a free group
on two generators:
$$
\pi_1^{\mathrm{top}}
\bigl(\mathbb{P}^1_{\mathbb{C}} \setminus
\{0,1,\infty\}\bigr)
\cong F_2.
$$
After profinite completion one gets
$$
\pi_1^{\acute{e}t}
\bigl(\mathbb{P}^1_{\overline{\mathbb{Q}}} \setminus
\{0,1,\infty\}\bigr)
\cong \widehat{F}_2,
$$
the free profinite group on two generators.

\medskip\noindent
Now let
$$
X=\mathbb{P}^1_{\mathbb{Q}} \setminus
\{0,1,\infty\}.
$$
There is an exact sequence of profinite
groups
$$
1\to
\pi_1^{\acute{e}t}(X_{\overline{\mathbb{Q}}})
\to
\pi_1^{\acute{e}t}(X)
\to
\operatorname{Gal}(\overline{\mathbb{Q}}/\mathbb{Q})
\to 1.
$$
Because the geometric fundamental group is
\(\widehat{F}_2\), the absolute Galois group
acts by outer automorphisms on a free
profinite group:
$$
\operatorname{Gal}(\overline{\mathbb{Q}}/\mathbb{Q})
\longrightarrow
\operatorname{Out}(\widehat{F}_2).
$$
This is one of the first places where
Grothendieck's vision becomes concrete.
The theorem says that this action is
faithful.

\medskip\noindent
In particular, the absolute Galois group can
be studied through its action on the
combinatorics of finite coverings of
\(\mathbb{P}^1\setminus\{0,1,\infty\}\).
This is the origin of dessins d'enfants.
A finite cover of
\(\mathbb{P}^1\setminus\{0,1,\infty\}\)
corresponds to a finite quotient of
\(\widehat{F}_2\), and the Galois action
permutes these quotients in a very rigid way.
That rigidity is what eventually leads to the
Grothendieck--Teichm\"uller group.

\medskip\noindent
The point of the talk is that a similar
picture appears in arithmetic topology.
Mapping class groups also act on profinite
fundamental groups of surfaces, and one can
study their automorphisms in a way that
mirrors the Galois action above. This
suggests profinite rigidity questions: to
what extent does the profinite completion of
a mapping class group remember the original
group? The Grothendieck--Teichm\"uller
framework provides a natural language for
these problems, and may lead to new rigidity
results.

\medskip\noindent
In the arithmetic case there are more
automorphisms than in the topological case,
and this is what makes the theory useful for
applications. The speaker suggested that
these methods should lead to a proof of a
failure of profinite rigidity: more
precisely, if \(d(S)\) denotes the complexity
of the surface, then there should exist
finite-index subgroups \(\Gamma\) and
\(\Gamma'\) of \(\Gamma(S)\) such that
$$
\widehat{\Gamma}\cong \widehat{\Gamma'}
\qquad\text{but}\qquad
\Gamma\not\cong \Gamma'.
$$

\medskip\noindent
For \(g=0\), the speaker can prove that every
pointwise inner automorphism of
\(\widehat{\Gamma}(S)\) is actually inner.
That is, if an automorphism acts on every
single element as conjugation, then the
automorphism itself is conjugation.

\medskip\noindent
As a corollary, the Grothendieck--Teichm\"uller
group acts faithfully on the conjugacy
classes of elements of the profinite
completion of the braid group
\(\widehat{B}_n\). This is interesting from
the topological point of view because of the
Alexander theorem, which characterizes
braids via their relationship with links.

\medskip\noindent
Another application mentioned in the talk is
Furusho's work on finite-type (Vassiliev)
knot theory. The GT action can be used to
organize knot-theoretic information in a way
that is compatible with the profinite braid
group viewpoint.

\medskip\noindent
Let $B$ denote the profinite mapping class
group under consideration. We write
\(\operatorname{Aut}^{\mathbb{I}_0}(B)\) for
the subgroup of \(\operatorname{Aut}(B)\)
consisting of those automorphisms that
preserve the procyclic subgroups generated by
the elements \(C_r\) arising from Dehn
twists. The centralizer of the natural image
of the Grothendieck--Teichm\"uller action in
\(\operatorname{Aut}^{\mathbb{I}_0}(B)\) is
what we call the Grothendieck--Teichm\"uller
group \(G(c,0)\).

\medskip\noindent
For a connected oriented surface
\(S_{g,n}\) of genus \(g\) with \(n\)
punctures and \(\chi(S_{g,n})<0\), the
mapping class group is
the group of orientation-preserving
diffeomorphisms modulo isotopy to the
identity. More precisely, if
\(\Gamma_{g,n}\) denotes the mapping class
group of \(S_{g,n}\), then
\(\Gamma_{g,n}=\mathrm{P}\Gamma(S_{g,n})\)
in the pure case.

\begin{theorem}\label{thm:mapping-class-moduli-space}
There is a natural isomorphism
$$
\Gamma_{g,n}\cong
\pi_1^{\mathrm{top}}
\bigl(\mathcal{M}_{g,n}(\mathbb{C})\bigr),
$$
where \(\mathcal{M}_{g,n}\) is the moduli
space of smooth genus \(g\) curves with
\(n\) marked points.
\end{theorem}

\begin{proof}
This is the classical identification of the
mapping class group with the orbifold
fundamental group of the moduli space. In the
case of ordered marked points, the pure
mapping class group corresponds to the
fundamental group of the pointed moduli
space.
\end{proof}

\medskip\noindent
There is also a standard exact sequence
relating the pure mapping class group to the
full mapping class group:
$$
1\to \widehat{\mathrm{P}\Gamma}_{g,n}
\to \widehat{\Gamma}_{g,n}
\to \Sigma_n
\to 1.
$$
At the discrete level this comes from the
action of the mapping class group on the set
of punctures, and the quotient records the
induced permutation of the punctures.

\begin{remark}
How can we study the Deligne--Mumford
boundary through points on a surface? This
appears to be a standard method discussed in
the talk.
\end{remark}

\medskip\noindent
The theorem implies that the \'etale
fundamental group of the moduli space is the
profinite completion of the mapping class
group:
$$
\pi_1^{\acute{e}t}
\bigl(\mathcal{M}_{g,n}\otimes_{\mathbb{Q}}\overline{\mathbb{Q}}\bigr)
\cong
\widehat{\Gamma}_{g,n}.
$$
In the pure case this is
\(\widehat{\mathrm{P}\Gamma}(S_{g,n})\).

\begin{theorem}[Hoshi--Mochizuki]\label{thm:hoshi-mochizuki-centralizer}
Let \(\Sigma_n\) act on the profinite
mapping class group in the natural way by
permuting the punctures. Then there is a
natural injection
$$
Z_{\operatorname{Out}^{\mathbb{I}_0}
(\widehat{\Gamma}_{g,n})}(\Sigma_n)
\hookrightarrow
Z_{\operatorname{Out}^{\mathbb{I}_0}
(\widehat{\Gamma}_{g,n+1})}(\Sigma_{n+1}),
$$
and this map is an isomorphism for \(n\ge 5\).
\end{theorem}

\begin{proof}
This is the centralizer comparison used in
the study of profinite mapping class groups
and their puncture-permutation symmetries.
The low-puncture cases require separate
analysis, while for \(n\ge 5\) the natural
stabilization map becomes an isomorphism.
\end{proof}

\begin{theorem}[Hoshi--Minamide--Mochizuki]\label{thm:hoshi-minamide-mochizuki}
There is a natural isomorphism
$$
\operatorname{Out}^{\mathbb{I}_0}
\bigl(\mathrm{P}\check\Gamma(S)\bigr)
\cong
\Sigma_n \times \widehat{\operatorname{GT}}.
$$
Equivalently, in the genus zero punctured
sphere case one may write
$$
\operatorname{Out}^{\mathbb{I}_0}
\bigl(\mathrm{P}\check\Gamma_{0,n}\bigr)
\cong
\Sigma_n \times \widehat{\operatorname{GT}}.
$$
\end{theorem}

\begin{proof}
This is the theorem that removes the extra
rigidity hypotheses and identifies the full
outer automorphism group, under the
\(\mathbb{I}_0\)-condition, as the product of
the puncture permutation group and the
profinite Grothendieck--Teichm\"uller group.
\end{proof}

\medskip\noindent
There is also a slogan in the genus-zero
setting:
$$
\operatorname{Out}\bigl(\widehat{\Gamma}_{0,n}\bigr)
\cong \widehat{\operatorname{GT}}
$$
for the appropriate completion and rigidity
condition encoded by the exponent on
\(\operatorname{Out}\). More generally, for
small genus \(g\leq 2\), the speaker proved a
corresponding isomorphism
$$
\operatorname{Out}^{\ast}\bigl(\widehat{\Gamma}_{g,n}\bigr)
\cong
\Sigma_n \times \widehat{\operatorname{GT}},
$$
where the superscript \(\ast\) indicates the
same rigidity condition as in the previous
theorems.

\section{Hodge structures in toric ambient}
\label{section-hodge-structures-toric-ambient}
%% speaker: Roberto Villaflor, USM, Chile
%% place: GADEPS Focused Meeting, IMPA
%% date: June 29, 2026

\noindent
Roberto Villaflor, USM, Chile.
GADEPS Focused Meeting June 29, 2026, IMPA.

\medskip\noindent
Joint with L. Broring, G. Martinez,
A. Viergever.

\medskip\noindent
{\bf Hodge structures over a perfect
field.} Let $X$ be a smooth
projective variety over $\mathbb{C}$.
We have the Hodge-De Rham spectral sequence
$E^{p,q}_1=H^q(X,\Omega_X^p) 
\implies \mathbb{H}^k(X,\Omega_X^\bullet)
\simeq H^k_{dR}(X/\mathbb{C})$.
It degenerates at $E_1$, giving
a non-canonical isomorphism
$$
H^\bullet_{dR}(X/\mathbb{C})
\simeq
\bigoplus_{p+q=k}H^q(X,\Omega_X^p)
$$
We may ask ourselves whether this is also
true for other fields.

\begin{remark}[Deligne-Illusie]
\label{remark-derham-to-hodge-not-degenerate}
In general, for $X/k$ perfect field,
this may not degenerate at $E_1$.
\end{remark}

\noindent
For $X = \{F=0\} \subseteq \mathbb{P}^{n+1}$
smooth degree $d$ hypersurface over
$\mathbb{C}$,
\begin{enumerate}
\item 
$$
H^q(X,\Omega_X^p)
\simeq \begin{cases}
H^q(X,\Omega_X^p)_{\operatorname{prim}}\qquad
&p \neq q \\
H^p(X,\Omega_X^p)_{\operatorname{prim}}\oplus
\mathbb{C}\theta |_X^p\qquad &
\end{cases}
$$
where $\theta$ is the class of a hyperplane.

\item
$H^q(X,\Omega_X^p)_{\operatorname{prim}}=0$
for $p \neq q \neq  n$ 
(Lefschetz hyperplane theorem).

\item Griffith basis theorem for $p+q=n$.
\end{enumerate}

\begin{remark}
\label{remark-lps}
For $X/k$ smooth projective, $k$ perfect
of characteristic $\neq 2$,
Levine-Pepin-Srinivas '24: (1), (2)
and (3) still hold.
\end{remark}

\medskip\noindent
{\bf Quadratic Euler characterstic.} 
Let $X$ be a CW-complex of dimension $n$.
$$
\chi_{\operatorname{top}}(X)
=\sum_{i=1}^n(-1)^i b_i(X),
$$
where
$b_i=\operatorname{rank}H_i(X,\mathbb{Z})$
is the $i$-th Betti number.

\begin{proposition}
\label{proposition-properties-of-euler}
\begin{enumerate}
\item $\chi_{\operatorname{top}}$
is homoropy invariant.
\item $Y \subseteq X$,
$\chi_{\operatorname{top}}(X\setminus Y)
=\chi_{\operatorname{top}}(X)
-\chi_{\operatorname{top}}(Y)+1$.
\item $F \hookrightarrow E \to B$ fiber
bundle, then $\chi_{\operatorname{top}}(E) =
\chi_{\operatorname{top}}(B)
\chi_{\operatorname{top}}(F)$

\item $\hat{\chi_{\operatorname{top}}}(X
\wedge Y) =
\hat{\chi_{\operatorname{top}}}(X)
\hat{\chi_{\operatorname{top}}}(Y)$.
\end{enumerate}
\end{proposition}

\begin{definition}
\label{definition-stably-homotopic}
$X$ is {\it stably homotopic} to $Y$
if
$$
\Sigma^kX \equiv \Sigma^kY
$$
where $\Sigma^k$ is the smash product
of $k$ 1-spheres.
\end{definition}

\noindent
By the reduced Euler characteristic formula,
if $X$ and $Y$ are stably homotopic,
then $\chi_{\operatorname{top}}(X)
=\chi_{\operatorname{top}}(Y)$.

\medskip\noindent
Morel-Boevodky '99: we have
a stable motivic homotopy category.

There exists an analogous construction
of categorical Euler characteristic
$X$ quasi-projective smooth $k$-scheme.
$\chi(X/k) \in GW(k)$
where $GW(k)$ is the Grothendieck-Witten
ring. In characteristic $\neq 2$, 
bilinear symmetric forms = group
generated by the monoid of isometry
classes of quadratic forms with $\oplus$.

The upshot is that instead
of assigning a number, now we put
a quadratic form to our scheme.

\begin{remark}
\label{remark-real-field-euler}
If $k \subset \mathbb{R}$ is a real field,
the rank of $\chi(X/k) =
\chi_{\operatorname{top}}(X(\mathbb{C}))$,
and the signature of $\chi(X/k)$
is
$\chi_{\operatorname{top}}(X(\mathbb{R}))$.
\end{remark}

\begin{theorem}[Levine '20]
\label{theorem-levine20}
$X/k$ smooth projective, $k$ perfect,
char $\neq 2$, then
\begin{enumerate}
\item If $\dim X$ is odd, then
$\chi(X/k) = mH$, $m \in \mathbb{Z}$,
$$
H=\begin{pmatrix}1&0\\ 0&-1\end{pmatrix}
=\langle 1\rangle+ \langle-1\rangle.
$$
\item If $\dim X$ is even, then
$\chi(X/k)=mH+Q$, $m \in \mathbb{Z}$
and $Q$ is determined by the first
page of the Hodge-to-de-Rham spectral
sequence:
$$
Q: H^{n/2}(X,\Omega^{n/2}_X)\times
H^{n/2}(X,\Omega_X^{n/2})
\xrightarrow{\cup }
H^n(X,\Omega_X^n)
\xrightarrow{\operatorname{trace}}k.
$$
\end{enumerate}
\end{theorem}

\noindent
By LPS '24, for $X = \{F=0\}
\subset \mathbb{P}^{n+1}_k$,
$Q = Q_{\operatorname{prim}} \oplus
\langle 1\rangle$ and
$$
Q_{\operatorname{prim}}:
J(X) _{(d-2)(n/2+1)}\times
J(X)_{(d-2)(n/2+1)}\xrightarrow{\cdot}
J(X)_{(d+1)(n+2)}
=\det (\text{Hess}((F)))
$$

\medskip\noindent
{\bf Non-split toric varieties.}
Recall: $X/\mathbb{C}$ is toric
if it admits an algebraic torus action
$T = (\mathbb{C}^\times)^n$, $n = \dim X$
such that there exists $x \in X$
with $T \simeq T\cdot x$ (dense open orbit).
When $X$ is normal there exists a
combinatorial object called the fan
$\Sigma = \bigcup_{k=1}^n \Sigma(k)$
consisting of cones inside $N_{\mathbb{R}}
\simeq M^4_{\mathbb{R}} \simeq \mathbb{R}^n$
where $M = \Hom(T,\mathbb{C}^\times)
\simeq \mathbb{Z}^n$
group of characters of $T$
for $\sigma \in \Sigma(k)$
if is a cone $\sigma \subseteq N_{\mathbb{R}}$
of dimension $k$, rational polyhedral
and strongly convex.

\begin{itemize}
\item for all $\sigma \in \Sigma$,
each face of $\sigma$ is in $\Sigma$,
\item for all $\sigma,\sigma' \in \Sigma$
$\sigma \cap \sigma' \in \Sigma$.
\end{itemize}

\noindent
For each $\sigma \in \Sigma$, we have
$\sigma^\vee \subseteq M_{\mathbb{R}}$
$\implies $ $\sigma^\vee \cap M \subset M$
is a semigroup, which gives
$\Spec \mathbb{C}[\sigma^\vee\cap M]
\supseteq \Spec \mathbb{C}[M]
\simeq T$.
Denote $U_\sigma = \Spec
\mathbb{C}[\sigma^\vee \cap M]$.
$\sigma \subseteq \sigma' \implies
(\sigma')^\vee \subseteq \sigma^\vee \implies 
U_\sigma \subseteq U_{\sigma'}$.

The idea is that the fan gives us the data
of how to glue the affine toric varieties
$\Spec \mathbb{C}[\sigma^\vee \cap M]$
into a toric variety.

When the rays of all cones ar $\mathbb{R}$-
linearly independent, we say
$X_{\Sigma}$ is simplicial.
\begin{itemize}
\item $X_\Sigma$ is an orbifold
(finite quotient singularity).
\item We can put homogeneous coordinates
on $X_{\Sigma}$, the {\it Cox ring}:
$$
\operatorname{Cox}(X_{\Sigma})
= \bigoplus_{\beta \in \text{Cl}(X_{\Sigma}}
H^0(\mathcal{O}_{X_{\Sigma}}(\beta)
$$
since $X_{\Sigma}\setminus T
= \bigcup_{D \in
\text{Div}_{T}(X_{\Sigma}})$.
\end{itemize}

\noindent
Each $\text{Div}_T(X_{\Sigma})
=\{D_{\rho_i}:i 1,…,r\}$,
$\Sigma(1)= \{\rho_1,…,\rho_r\}$.
$\operatorname{Cox}(X_{\Sigma})
=\mathbb{C}[X_{\rho_1},…,X_{\rho_r}]$.
This ring is naturally graded by
$\text{Cl}(X_{\Sigma})$,
$\text{deg}(X_{\rho_i})=
[D_{\rho_i}] \in \text{Cl}(X_{\Sigma})
=\bigoplus_{i=1}^r \mathbb{Z}\cdot
D_{\rho_i}/\text{div}(M)$.

A non-split torus $T$ is a $k$-linear
algebraic group such that
$T_{\overline{k}}\simeq
(\overline{k}^\times)$ is a split torus.

\begin{itemize}
\item $S^1$ is a non-split
$\mathbb{R}$-torus.
Given a $k$-torus $T$, a $T$-toric
variety is a $k$-scheme $T$ acting on $X$
such that $T_{\overline{k}}$ acts on
$X_{\overline{k}}$ becomes a usual
toric variety.
\end{itemize}

\noindent
The Galois group $\Gamma=\operatorname{Gal}
(\overline{k}/k)$ acts on $T$ and $X$.

\begin{theorem}
\label{theorem-toric-primitive-hodge}
$Y \subseteq X$, $T$-toric,
$$
H^q(Y,\Omega_Y^p)_{\operatorname{prim}}
\simeq
(Cox(X_{\overline{k}})
/\langle \frac{\partial F}{\partial x_1}
,…,\frac{\partial F}{\partial x_r}\rangle)
^\Gamma_{[\gamma]}
$$
\end{theorem}

\section{Determinantal surfaces
and Noether-Lefschetz theory}
\label{section-determinantal-surfaces-nl}
%% speaker: Manuel Leal, IMPA
%% place: GADEPS Focused Meeting, IMPA
%% date: June 30, 2026

\noindent
Manuel Leal, IMPA.
GADEPS Focused Meeting June 30, 2026, IMPA.

\medskip\noindent
The talk studies determinantal surfaces
inside the parameter space of smooth
degree $d$ surfaces in $\mathbb{P}^3$.
Such a surface contains natural ACM curves,
which give extra line bundles besides
$\mathcal{O}_X(1)$.
Thus determinantal surfaces lie in the
Noether-Lefschetz locus. The main point is
that, for the degree type $(a,b)$, the
determinantal locus is itself a
Noether-Lefschetz component.

\medskip\noindent
Given a determinantal surface given
by the vanishing of the determinant
of a matrix $S$, we can define two curves:
$\mathcal{C}_+$ and $\mathcal{C}_-$
both given by the vanishing of maximal
rank minors of the matrix of $S$ after
removing everything except the highest/lowest
degree forms of the matrix.

These two curves are arithmetically
Cohen-McCaulay.

\begin{theorem}[L,LH,V, 24]
\label{theorem-llhv24}
$$
\xymatrix{
&\{(\mathcal{C}_+,X):\mathcal{C}_+ \subseteq X\}
\ar[dr]\ar[dl]\\
\mathcal{H}_{a,b}\subset
\operatorname{Hilb}_{\mathbb{P}^3}
& &
\det(a,b) \subseteq U_d
}
$$
where
$U_d$ is the open subset of
$|\mathcal{O}_{\mathbb{P}^3}(d)|$
parametrizing smooth surfaces of degree $d$.
The locus $\det(a,b)\subseteq U_d$
parametrizes surfaces cut out by the
determinant of a matrix whose entries have
degree type $(a,b)$.
The scheme $\mathcal{H}_{a,b}$ is the
Hilbert scheme component parametrizing
the corresponding arithmetically
Cohen-McCaulay curves.
\end{theorem}

\noindent
This was proved using
\begin{theorem}[Ellingsrad 75]
\label{theorem-ellin75}
$\mathcal{H}_{a,b} \subseteq
\operatorname{Hilb}_{\mathbb{P}^3}$
is an irreducible component gen. smooth
and of dimension…
\end{theorem}

\begin{proposition}
\label{proposition-duplicate-2}
\begin{itemize}
\item $\pi^{-1}(X)
=
\coprod_{\mathcal{C}_+ \subseteq X}
|\mathcal{O}_\ell(\mathcal{C}_+)|$.
\item $H^0(X,\mathcal{C}_+)
=
\chi(X,\mathcal{C}_+)$.
\end{itemize}
\end{proposition}

\medskip\noindent
{\bf Noether-Lefschetz theorem.} 
If $d \geq  4$ and $X \in U_d$
is very general,
$$
\text{Pic}\mathbb{P}^3
\xrightarrow{\simeq}
\text{Pic}X 
=
\langle \mathcal{O}_X(1)\rangle
\simeq \mathbb{Z}
$$
where the firs map is a resolution.

\begin{definition}
\label{definition-nl-locus}
The {\it Noether-Lefschetz locus} is
$$
NL(d)
=
\{ X \in U_d: \text{Pic}(X)\not\simeq
\mathbb{Z}\}.
$$
\end{definition}

\noindent
$$
NL(d) = \bigcup_{\operatorname{countable}}
\Sigma
$$
where $\Sigma$ are the NL-compononents.

\begin{theorem}[L,LH,V 24]
\label{theorem-llhv24ii}
$\det(a,b)$ are NL components.
\end{theorem}

\begin{theorem}[Lopez, 91]
\label{theorem-lopez-91}
A ver general degree $d$ surface
$X \supseteq \mathcal{C}_-$,
$$
\text{Pic}(X) = 
\langle \mathcal{O}_X(1),
\mathcal{O}_X(\mathcal{C}_-)\rangle.
$$
\end{theorem}

\noindent
This implies that $\det(a,b) \subseteq
NL(d)$,
and this implies that
$\det(a,b) \subseteq \Sigma$ for some
NL component.

Take $\{ X_s\} \subseteq \Sigma$
and $X_s = X \in \det(a,b)$.
There is a family of line bundles
$\{\mathcal{L}_s \in \text{Pic}(X)\setminus
\langle\mathcal{O}_{X_s}(\ell)\rangle\}$.
Since $\rho(X) = 2$ we can assume that
$\mathcal{L}_0 =
\mathcal{O}_X(\mathcal{C}_+)$.
$H^1(X,\mathcal{O}_X(\mathcal{C}_+))=0$
this implies that there exists
$\{\mathcal{C}_s \subseteq X_s\}$ such that
$\mathcal{C}_0 = \mathcal{C}_{+}$.
By the theorem of Elli, this implies that
$\mathcal{C}_s \in \mathcal{H}_{a,b}$ for
gen. s. This implies that
$X_s \supseteq \mathcal{C}_s$ is in 
$\det(a,b)$. This concludes the ideas
behind the proof.

\section{On a counterexample to a conjecture
of J. Harris for octic surfaces}
\label{section-counterexample-harris-conj}
%% speaker: Hossein Movasati, IMPA
%% place: GADEPS Focused Meeting, IMPA
%% date: June 30, 2026

\noindent
Hossein Movasati, IMPA.
GADEPS Focused Meeting June 30, 2026, IMPA.

\medskip\noindent
Based on works by Hossein Movasati:
2021 - Special components of
Noether-Lefschetz. 2026 - On a
counterexample…

\medskip\noindent
Let $T $ be the parameter space of smooth
degree $d$ surfaces on $\mathbb{P}^3$ over
$\mathbb{C}$.
Consider
$$
x_0^d+x_1^d+x_2^d+x_3^d
+\sum_{\text{deg}x^\alpha =d}t_\alpha
x^\alpha,
$$
where there are 165 parameters $t_\alpha$.
$d=8$.

$$
NL_\alpha
=
\{t \in T|\rho(X_t)>1\}
$$
where $\rho$ is the Picard number.
As explained in Manuel's talk,
$NL_\alpha$ is a union of enumerable
sets and
$$
5=d-3 \leq  \text{codim}(\text{components} )
\leq \binom{d-1}{3}=h^{2,0}
$$
Components which attain the upper bound
are called general components, and those
for which the right-hand inequality
is strict are called special.
General components are dense in $T$.

Harris' conjecture: the number of special
components is finite.
In the 90's, Voisin found a counter-example
for large $d$.
The conjecture is known to hold
for $d=5$. For $d=6,7$ ``the conjecture
must be true''. Thus, the next interesting
case is $d=8$.

\begin{exercise}
\label{exercise-computation-codim}
The codimension of the set of octic
surfaces containing a line is 5.
The codimension of the set of octic surfaces
containing a complete intersection
of type $(3,3)$ is 27.
Consider $x_0^8+x_1^8+x_2^8+x_3^8$
and $C_1:f_1=g_1=0$, $C_2:f_3=g_3=0$ where
the $f_i$ and $g_i$ are factors
of the Fermat equation. Then $V_{C_1}$
and $V_{C_2}$ are smooth and intersect
transversely, $\text{codim}(V_{C_1}\cap
V_{C_2})=5+27=32$.
\end{exercise}

\noindent
Let
$$
\delta_0=[C_1]+r[C_2]
\in H^2(X_0,Q) \cap H^{1,1} 
$$
where $C_1,C_2$ are as in the
exercise and $X_0$ is the Fermat,
and $r \in \mathbb{Q}$ is a parameter,
making the above a pencil of classes. 

Recall that Hodge conjecture for the case
of surfaces is called Lefschetz theorem
on $(1,1)$-classes.

Also,
$$
\delta_t \in H^2(X_t,Q),\qquad 
t \in (T,0).
$$
Local Noether-Lefschetz locus:
$$
V_{\delta_0}:=
\{t \in (T,0):
\delta_t \in H^2(X_t,\mathbb{Q})
\cap H^{1,1}(X_t)\}.
$$

In principle, this $V_{\delta_t}$
is in some component of the
Nother-Lefschetz locus.

\medskip\noindent
Conjecture: for all $r \in \mathbb{Q}$,
$r \neq  0$, except a finite number,
$V_{\delta_0}$ are distinct, smooth
varieties of codimension 31.
(Recall that $31 < 35 = h^{2,0}$.)


\section{Asymptotic behaviour of tropical
rank functions}
\label{section-asymptotic-tropical-rank}
%% speaker: Eduardo Vital, Bielefeld University
%% place: Seminar of Algebra, IMPA, room 228
%% date: August 26, 2026, 15:30

\noindent
Eduardo Vital, Bielefeld University.
Seminar of Algebra, IMPA.
August 26, 2026, 15:30, room 228.

\medskip\noindent
{\bf Abstract.} In this talk, we study the
asymptotic behaviour of two notions of rank
for linear series on tropical curves.
Although these two notions can differ for a
fixed divisor, we show that they have the
same asymptotic growth. This leads to a
natural notion of tropical volume.

We explain the main ideas behind this
result, including the role of uniform bounds
for the Baker--Norine rank and tropical
independence. As a consequence, both notions
of rank satisfy an asymptotic
Riemann--Roch formula. We also discuss
examples where the two ranks differ at
finite level, but nevertheless have
identical asymptotic behaviour. Finally, we
relate the tropical volume to the classical
volume of divisors through tropicalization.
This is joint work with Ana Maria Botero and
Alex K\"uronya.

\medskip\noindent
{\bf Comment.}
Here I learnt that scheme theory
for tropical geometry has not been
developed. The speaker and his collaborators
are interested in developing
Cech cohomology for higher dimensions
(and in developing scheme theory
for tropical geometry).

\section{\texorpdfstring
{$C^\infty$}{C-infinity}-Superrings and
\texorpdfstring
{$C^\infty$}{C-infinity}-Superschemes}
\label{section-rizzo-smooth-superrings}
%% speaker: Pedro Rizzo
%% place: Seminar of Algebra, IMPA, room 228
%% date: September 9, 2026, 15:30

\noindent
Pedro Rizzo.
Seminar of Algebra, IMPA.
September 9, 2026, 15:30, room 228.

\medskip\noindent
{\bf Abstract.}
Superalgebraic and supergeometric structures
generalize classical nongraded structures to
graded ones. This extension introduces an
interplay between commuting and anticommuting
elements, enriching tools of commutative
algebra, differential geometry, and algebraic
geometry. Although originally motivated by
theoretical physics, supergeometry has
evolved into a robust mathematical framework
for studying objects such as supermanifolds
and superschemes.

In this talk, we introduce
$C^\infty$-superrings and their associated
geometric spaces, which we term
$C^\infty$-superschemes. We establish an
equivalence of categories between these
algebraic and geometric structures.
Furthermore, we present a generalization of
Batchelor’s theorem to $C^\infty$-superspaces
extending the classical result that every
smooth supermanifold is (noncanonically)
isomorphic to the sheaf of sections of a
vector bundle over its underlying smooth
manifold.

Conceptually, $C^\infty$-superrings
synthesize the features of $C^\infty$-rings
and superrings. Since smooth supermanifolds
form a core subclass of
$C^\infty$-superscheme, this framework
provides a natural scheme-theoretic
generalization of supergeometry in the smooth
setting. This is a joint work with Danilo
Olarte and Alexander Torres-Gómez (both at
Universidad de Antioquia) and based on the
preprints: arXiv:2508.09900v3 and
arXiv:2605.07169v2



\medskip\noindent
{\bf Introduction.}
$\mathbb{Z}_2$-structures
are inspired by commutative ``bosonic''
and anticommutative ``anti-bosonic''
elements. Its origins lie in theoretical
physics during the 70's and 80's, driven by
developments in supersymmetric field theory.

These notions are important in
supergeometry.
 
Synthetic differential geometry
also emerged in the 80's. Dubuc, Moerdijk,
Reyes and others constructed bridges between
differential and algebraic geometry.

Derived differential geometry attempts to
use techniques of algebraic geometry to
$C^\infty$ rings rather than commutative
rings.

Nishimura and Yetter explored this problem
using Weil superalgebras (which are a
particular case of $C^\infty$-superrings).
Liu and Yau introduced a notion of
$C^\infty$-superring to propose a work in
theoretical physics.

\medskip\noindent
{\bf Context and goal.}

Goal I: to establish structured mathematical
foundation of $C^\infty$ superspaces.

See results by Batchelor and Donagi-Witten.

Goal II: we propose the notion of Batchelor
space, which attempts to generalize
Batchelor's result.

\medskip\noindent
{\bf Category of $C^\infty$-rings}

\begin{definition}[C infinity ring]
\label{definition-cinf-ring}

\end{definition}

\begin{itemize}
\item A morphism s…
\item $\mathfrak{C}$ is {\it finitely
generated} if…
\item An {\it ideal} is…
\item Obvious example: the
$\mathbb{R}$-algebra of smooth functions on
a differentiable manifold.
\item Every finitely-generated
$C^\infty$-ring is isomorphic to a quotient
of $C^\infty(\mathbb{R}^n)$ by some ideal.
\item An {\it $\mathbb{R}$-point} is a
$C^\infty$-ring morphism $x:\mathfrak{C}\to
\mathbb{R}$.
\item {\it Localization} is…
\item $\mathfrak{C}$ is called {\it fair} if
it is finitely generated and
germ-determined.
\end{itemize}

\medskip\noindent
{\bf $C^\infty$-ringed space.}

\begin{definition}
\label{definition-cinf-ringed-space}

\end{definition}

\begin{itemize}
\item The stalk is localization…
\item Other expected results…
\end{itemize}

\noindent $\Spec$ and $\Gamma$ (sections)
functors are…

\begin{theorem}[Joyce]
\label{theorem-joyce-adjoint-functors}
$\Spec$ and $\Gamma$ are adjoint.
\end{theorem}

\medskip\noindent
{\bf $C^\infty$-modules and sheaves.}

\begin{itemize}
\item A {\it module over a $C^\infty$-ring} 
is…
\item 
\end{itemize}

\medskip\noindent
{\bf Superring: basic theory.}

\begin{itemize}
\item A {\it superring} is a
$\mathbb{Z}_2$-graded unitary
supercommutative ring.
\item
$\mathfrak{R}=\mathfrak{R}_{\overline{0}}
\oplus \mathfrak{R}_{\overline{1}}$
such that $\mathfrak{R}_i,
\mathfrak{R}_j \subseteq
\mathfrak{R}_{i+j}$ (mod 2).
\item A {\it morphism} is…
\item An {\it ideal} is…
\item The quotient $\mathfrak{R}/I$
is a superring…
\item The {\it canonical superideal} is…
which gives the {\it superreduced ring}…
\end{itemize}

\noindent
Batchler's result concerns precisely when
can we ``lift'' a property from the reduced
ring to all of it.

\begin{example}
\label{example-polynomial-superring}
Polynomial superring over the field
$\mathbb{K}$…
\end{example}

\begin{itemize}
\item We define $\Lambda(\mathfrak{R})$
as…
\item {\it Localization} is…
\end{itemize}

\medskip\noindent
{\bf $C^\infty$-superrings.}
Now we mix the two constructions so far.

\begin{definition}
\label{definition-cinf-superring}
A {\it $C^\infty$-superring} is…
\end{definition}



\begin{proposition}
\label{proposition-superring}
(Essentially describes the object
we shall call {\it split}.)
\end{proposition}

\noindent
[Missing content] We may thus define
super-versions of the adjoint functors
$\Spec$ and $\Gamma$…

\begin{theorem}[A]
\label{theorem-a-superring}
$s\Spec$ and $s\Gamma$ are adjoint.
\end{theorem}

\begin{theorem}[B]
\label{theorem-b-superring}
\end{theorem}

\medskip\noindent
{\bf Split spaces.}
{\it Split spaces} are… $\mathcal{X}$ is
{\it projected} if…

\begin{theorem}[C]
\label{theorem-c-superring}
Every Batchelor space is split.
\end{theorem}


\medskip\noindent
{\bf Ongoing work.}
With Bilson Castro and Alexander Torres.
Based on Bruzzo et al., understand
$C^\infty$-superstacks.

\section{Quantised Picard Fuchs equations
and quantum periods}
\label{section-garay-quantum-periods}
%% speaker: Mauricio Garay
%% place: GADEPs seminar, IMPA, room 345
%% date: October 14, 2026, 10:30

\noindent
Mauricio Garay.
Seminário de Geometria, Aritmética e
Equações Diferenciais de Períodos
(GADEPs), IMPA.
October 14, 2026, 10:30. Room 345.

\medskip\noindent
{\bf Abstract.}
The founders of quantum mechanics (Born,
Heisenberg, Dirac and Jordan) proposed an
interpretation of the spectrum as a
quantised version of normal forms in
classical mechanics, hence the name quantum
mechanics. The inverse of this normal form,
improperly called the quantum Birkhoff
normal form, defines a function which is a
quantum analog of the period integrals.
We show that like in the classical case,
these satisfy ordinary linear differential
equations whose degree is bounded by the
rank of homology like for ordinary
Picard-Fuchs theory. Computations for the
standard cubic oscillator show that the
coefficients of the quantum-period expansion
are integers, suggesting a new kind of
quantisation of the classical Gromov-Witten
theory of $(K_{\mathbb{P}^2})$.

\medskip\noindent
{\bf Upshot.}
Classically, integrating along a closed
curve in an energy level measures an
action, related to the area it encloses.
As the energy varies, these integrals
satisfy linear differential equations:
Picard--Fuchs equations. The talk replaces
commuting coordinates by quantum
operators and constructs a corrected
version of this area as a formal series.
Its main claim is that these quantum
periods still satisfy linear differential
equations controlled by the topology of
the classical fibres. Thus the familiar
geometry of periods survives quantization.
The observed integer coefficients suggest
a further enumerative connection; the
abstract presents this as a suggestion.

\medskip\noindent
{\bf Source.}
\href{https://impa.br/seminarios/}
{IMPA seminar announcement},
consulted October 14, 2026.
Abstract reproduced verbatim from the
text supplied by Daniel; mathematical
notation is typeset in LaTeX.

\medskip\noindent
{\bf Lecture notes (provisional).}
He first inverted the Birkhoff function.
The precise normalization is not yet
recorded. The next step studies an action
integral on a moving cycle in $f=t$:
$$
A(t)=\int_{\gamma_t}y\,dx.
$$
The Gauss--Manin derivative differentiates
its cohomology class as the fibre varies:
$$
A'(t)=\int_{\gamma_t}\eta,
\qquad
df\wedge\eta=dy\wedge dx.
$$
This explains the board notation
$(dy\wedge dx)/df$: it denotes a
Gelfand--Leray form, not ordinary division.
The wedge order fixes the sign.

He computes $A,A',A''$ to obtain a
Picard--Fuchs equation: a linear relation
among these functions with coefficients
depending on $t$. The idea is to find a
relation among the corresponding
cohomology classes, then integrate.
The exact equation and the link with
the inverted Birkhoff function remain
to be recorded. This is a reconstruction
of the reported board notes, not a
verified transcription.

\medskip\noindent
{\bf Quantum Birkhoff normal form.}
Born--Heisenberg--Jordan (1925),
as mentioned in the talk.
Now $q,p$ do not commute:
$$
[q,p]=qp-pq=h,\qquad h=i\hbar,
$$
$$
H=qp+q^3+p^3.
$$
The board uses $\mathbb{C}[[q,p,h]]$,
with multiplication subject to this
commutation relation.

The reported normal-form statement
involves a series
$$
\psi(qp)=a_0qp+a_1(qp)^2+\cdots.
$$
The idea is the noncommutative analogue
of Birkhoff normal form: transform $H$
into a function of $qp$.

{\it Formula to check.}
The transcription reads
$$
e^{-f/h}=\psi(qp).
$$
The role of $H$ and the transformation
are missing from this formula; it is
retained as an incomplete board note,
not asserted as a theorem.

\medskip\noindent
{\bf Quantum area and the classical limit.}
The next definition clarifies the role
of the inverse Birkhoff function:
$$
A(t)=\psi^{-1}(t).
$$
This is called the quantum area.
The reported expansion is
$$
A(t)=A_0(t)+hA_1(t)+h^2A_2(t)+\cdots.
$$
Our provisional understanding is that
$A_0$ recovers the classical action,
up to the normalization used in the talk,
and the higher terms are quantum
corrections. The aim is a quantum
Picard--Fuchs equation for this series.

An integrality statement was reported as
$$
V\in\mathbb{Z}[q,p]
\quad\Longrightarrow\quad
A\in\mathbb{Z}[t][[\hbar]].
$$
The definition of $V$, hypotheses,
normalization, and order of completion
still need checking. This is a board
transcription, not a verified theorem.
The distinction between $h$ and $\hbar$
also matters since $h=i\hbar$.

A possible spectrum formula was noted as
$$
\operatorname{Sp}_{BHJ}(H)=\psi(n\hbar).
$$
The first factor may be an integer $n$;
its range and the spectral convention
remain unconfirmed. This does not yet resolve
the earlier exponential formula.

\medskip\noindent
{\bf Quantum differential form
(reported theorem, unverified).}
The hypotheses include
$f\in\mathbb{C}[q,p]$ with Morse critical
points. A further hypothesis was heard
as “Infinite sing ifolita”, possibly
“isolated singularities at infinity”;
this interpretation is unconfirmed.
The reported conclusion is a formal form
$$
\alpha=\alpha(f)\in\Omega^1[[h]],
\qquad
\alpha=y\,dx+h\alpha_1+h^2\alpha_2+\cdots.
$$
The ambient space and allowed poles
have not yet been recorded.
The integral formula was transcribed as
$$
\frac{\partial A_i}{\partial t}
=\int_{\gamma_1}\alpha.
$$
The indices and derivative need checking:
it is unclear whether this concerns
$A$, its coefficients, or its derivative.
We do not identify these alternatives.
The intended picture is that quantum
data can be expressed using a formal
series of differential forms on the
classical fibres.

The reported Picard--Fuchs “degree” is
$2g(f^{-1}(t))$. Here degree presumably
means differential order. The announcement
states a bound by homology rank.
Equality, compactification, and possible
contributions at infinity remain to check.
See the background reference below.

\medskip\noindent
{\bf Deformation quantization and ordering.}
This calculation uses $pq-qp=h$,
the opposite sign to the commutator
recorded earlier. With $h$ central,
$$
qpq=q(qp+h)=q^2p+hq.
$$
Normal ordering places $q$ before $p$.
On commuting polynomial symbols, the
corresponding star product satisfies
$$
q\ast_N p=qp,\qquad p\ast_N q=qp+h,
$$
$$
q\ast_N p\ast_N q=q^2p+hq.
$$
Symmetric ordering gives the Moyal
convention
$$
q\ast_M p=qp-\frac{h}{2},
\qquad
p\ast_M q=qp+\frac{h}{2}.
$$
Both have $p\ast q-q\ast p=h$.
The terms $hq$ and $h/2$ correct
algebraic slips in the transcription.

An ordering map $Q$ encodes operator
multiplication on symbols by
$$
Q(F\ast G)=Q(F)Q(G).
$$
Compare Sections
\ref{quantum-mechanics-%
section-dirac-rule} and
\ref{symplectic-geometry-%
section-kostant-souriau-formula}.
Prequantization supplies a commutator
representation; it is not itself the
invertible symbol map $Q$ used here.

The analogous reordering of field modes
appears in {\it Vertex algebras},
Definition
\ref{vertex-algebras-%
definition-creat-annih-opera}
and Exercise
\ref{vertex-algebras-%
exercise-most-important}:
$$
:a(z)b(z):=a(z)_+b(z)+b(z)a(z)_-.
$$
The common mechanism is reordering
noncommuting operators and retaining
commutator corrections. The field
construction involves infinitely many
modes and is not an identification of
the two products.

\medskip\noindent
{\bf Trace formula
(reported theorem, unverified).}
The final formula was transcribed as
$$
A'(t)=\operatorname{Tr}
\left(\frac{k}{f-t}\right).
$$
The meaning of $k$, the trace, the inverse
of $f-t$, and the hypotheses were not
recorded. The formula is retained as a
board transcription, without identifying
the trace with an ordinary matrix or
operator trace. No precise theorem
reference has been verified.

\medskip\noindent
{\bf Background and connections.}
See {\it Variations of Hodge Structures},
Section \ref{variations-of-hodge-structures-%
section-quantum-periods-background},
for periods, Gauss--Manin,
Picard--Fuchs equations, and links to
the quantization and Seiberg--Witten notes.

\section{Irrationality loci and
Noether--Lefschetz divisors in the moduli
space of K3s}
\label{section-chen-irrationality-k3}
%% speaker: Nathan Chen
%% place: Algebra Seminar, IMPA, room 228
%% date: October 14, 2026, 15:30

\noindent
Nathan Chen.
Algebra Seminar, IMPA.
October 14, 2026, 15:30.
Room 228.

\medskip\noindent
{\bf Abstract.}
The goal of this talk is twofold. In the
first part, we will present a family of
surfaces where the locus of points in the
base with constant degree of irrationality is
not constructible. In the second part, we
will explore the notation of gonality loci in
the moduli space of K3 surfaces. Among other
things, we prove that the polarized gonality
loci are always unions of Noether--Lefschetz
divisors.

\medskip\noindent
{\bf Source.}
\href{https://impa.br/seminarios/}
{IMPA seminar announcement},
consulted October 14, 2026.


\medskip\noindent
{\bf Setting.}
We work over $\mathbb{C}$.
The base $T$ is quasi-projective and
irreducible. Maps
$\mathcal{C}\to T$ or $\mathcal{S}\to T$
are called families, with flatness,
properness, smoothness, etc., specified
as needed. Below the curves and surfaces
are smooth and projective.

\medskip\noindent
{\bf Motivation: curves.}
Let $\mathcal{C}\to T$ be a smooth
projective family of curves.
The gonality is
$$
\operatorname{gon}(C)=\min_f\deg(f),
\qquad f:C\to\mathbb{P}^1,
$$
where $f$ ranges over nonconstant
morphisms.
For $\dim T=1$, gonality can only drop
at special points:
$$
\operatorname{gon}(C_0)
\leq\operatorname{gon}(C_t)
$$
for general $t$.
Constant-gonality loci are locally closed.
These are the motivating facts from
curve theory reported in the talk;
an exact reference remains to be added.
The special fibre belongs on the left,
correcting the transcription.

\medskip\noindent
{\bf Rationality: surfaces.}
For higher dimensions, start by asking
about rationality.
Recall the invariants
$$
q(S)=h^1(S,\mathcal{O}_S),
\qquad
P_2(S)=h^0(S,\omega_S^{\otimes 2}).
$$

\begin{theorem}[Castelnuovo]
\label{theorem-chen-castelnuovo}
A smooth projective complex surface $S$
is rational if and only if
$$
q(S)=P_2(S)=0.
$$
\end{theorem}

\noindent
Reference: Beauville,
{\it Complex Algebraic Surfaces},
second edition, Chapter V, Theorem V.1.
This corrects the garbled equality
in the board transcription.

\medskip\noindent
{\bf Consequence reported in the talk.}
In a smooth projective family of surfaces
over a connected smooth complex base,
if $S_0$ is rational, then every $S_t$
is rational. The proof uses constancy
of $q$ and $P_2$ in the family.
The deformation-invariance reference
for this step remains to be added.

\medskip\noindent
{\bf Higher dimensions.}

\begin{theorem}[Hassett--Pirutka--Tschinkel]
\label{theorem-chen-hpt}
There exists a smooth family
$\mathcal{X}\to T$ of projective
fourfolds such that:
\begin{enumerate}
\item A very general $X_t$ is not stably
rational, hence is not rational.
\item The locus of $t$ for which $X_t$
is rational is dense in the Euclidean
topology.
\end{enumerate}
Each fourfold is a quadric surface
bundle over $\mathbb{P}^2$.
\end{theorem}

\noindent
Reference: Hassett, Pirutka and Tschinkel,
{\it Stable rationality of quadric
surface bundles over surfaces},
Acta Mathematica 220 (2018), 341--365,
Theorem 1;
\href{https://arxiv.org/abs/1603.09262}
{author manuscript}.

\medskip\noindent
{\bf Degree of irrationality.}
If rationality behaves this way, what
happens to a numerical measure of
irrationality?

\begin{definition}
\label{definition-chen-irrationality}
For a smooth projective variety $X$
of dimension $n$, its
{\it degree of irrationality} is
$$
\operatorname{irr}(X)=\min_f\deg(f),
\qquad f:X\dashrightarrow\mathbb{P}^n,
$$
where $f$ ranges over dominant
generically finite rational maps.
\end{definition}

\noindent
Thus $\operatorname{irr}(X)=1$ precisely
when $X$ is rational.
For a smooth family of surfaces,
can this invariant drop?

\begin{theorem}[Chen--Stapleton]
\label{theorem-chen-stapleton-specialization}
Let $\mathcal{S}\to T$ be a smooth
projective family of surfaces over a
smooth irreducible complex curve,
with $0\in T$ and $q(S_0)=0$.
Then, for very general $t$,
$$
\operatorname{irr}(S_0)
\leq\operatorname{irr}(S_t).
$$
\end{theorem}

\noindent
Reference: Chen and Stapleton,
{\it Fano hypersurfaces with arbitrarily
large degrees of irrationality},
Forum of Mathematics, Sigma 8 (2020),
e24, Proposition C;
\href{https://arxiv.org/abs/1908.02803}
{author manuscript}.
The published statement assumes
$q(S_t)=0$ for very general $t$;
constancy of $q$ gives the formulation
above. The name transcribed as
Stephensen is Stapleton.

\medskip\noindent
{\bf Consequence of today's talk
(reported, unverified).}
Attribution recorded as B\_CD;
the full names and exact reference
remain to be checked.
There exists a smooth family of
projective surfaces $\mathcal{S}\to T$
and an integer $d\geq3$ such that:
\begin{enumerate}
\item For very general $t$,
$\operatorname{irr}(S_t)=d$.
\item There are infinitely many parameters
with $\operatorname{irr}(S_t)<d$,
and this locus is Euclidean dense.
\end{enumerate}

\medskip\noindent
{\bf Bootstrapping.}
The speaker says that bootstrapping a
certain family of K3 surfaces produces
families of surfaces of general type.
Here this means using the K3 examples
as input to a further construction,
transferring their irrationality
phenomenon to the new surfaces.
This is an interpretation of the phrase;
the actual construction has not yet
been recorded.


\medskip\noindent
The speaker also mentioned products with
$\mathbb{P}^n$ for higher-dimensional
examples. The required control of
irrationality was not recorded.

\medskip\noindent
{\bf The degree-two case.}
Chen--Esser prove that irrationality
at most two specializes in families
of integral klt surfaces over a smooth
complex curve. In a smooth projective
family with very general irrationality
two, all fibres have irrationality two:
a rational fibre would force rationality
throughout, by the earlier argument.

Reference: Chen and Esser,
{\it Group actions and irrationality
in surface families},
Selecta Mathematica 31 (2025), article 96.
Theorem 1.2 and Corollary 3.1 were
checked in the
\href{https://arxiv.org/abs/2410.18427}
{author manuscript}, version 1.

\medskip\noindent
{\bf K3s and hyperelliptic curves.}
The following equivalence was reported
in the seminar; its exact source
remains to be verified.
For a projective K3 surface $X$:
\begin{enumerate}
\item $\operatorname{irr}(X)=2$.
\item $X$ contains a smooth hyperelliptic
curve of genus at least two.
\item $X$ contains a smooth genus-two
curve, or admits a genus-one fibration
$X\to\mathbb{P}^1$ with a multisection
of degree at most two.
\end{enumerate}
The implication from (1) to (2) was
proved in the seminar; the proof was
not transcribed. In (3), $X$ replaces
the accidental $C$ in the notes.

\medskip\noindent
{\bf Lattice criterion
(board transcription, unverified).}
A further equivalence for
$\operatorname{irr}(X)=2$ used a line
bundle or one of these intersection
matrices in $\operatorname{Pic}(X)$:
$$
\begin{pmatrix}0&2\\2&0\end{pmatrix},
\qquad
\begin{pmatrix}0&1\\1&-2\end{pmatrix},
\qquad
\begin{pmatrix}0&2\\2&-2\end{pmatrix}.
$$
The other alternative was transcribed
as an ample $L$ with $L^2=0$.
This is impossible: an ample line
bundle on a surface has positive square.
A genus-two curve has square two by
adjunction, suggesting a square-two
condition. The precise positivity and
base-point-freeness hypotheses must
be checked before replacing the claim.
Embedding and primitivity conditions
for the lattices also remain to check.

\medskip\noindent
{\bf Quasi-polarized moduli.}
Write $\mathcal{F}_g$ for the moduli
space of quasi-polarized K3 surfaces
$(X,L)$, with $L$ primitive, big and nef,
and
$$
L^2=2g-2.
$$

\medskip\noindent
{\bf Corollary reported in the seminar.}
For each $g\geq3$, $\mathcal{F}_g$
contains infinitely many divisors whose
generic member has irrationality two.
The attribution and published reference
have not yet been identified.
The phrase about finitely many covers
of $\mathbb{F}_n$ was also recorded;
its intended quantifiers are unclear.

\medskip\noindent
{\bf AI-reconstructed proof outline
(incomplete, not a verified proof).}
Impose an extra divisor class which
produces a degree-two map to a rational
surface. A possible route is to seek
rank-two Picard lattices generated by
$L,H$ with intersection matrix
$$
\begin{pmatrix}
2g-2&a\\
a&2
\end{pmatrix}.
$$
Require $H$ nef and base-point-free.
Riemann--Roch and vanishing then give
$h^0(X,H)=3$; since $H^2=2$, the map
defined by $|H|$ has degree two onto
$\mathbb{P}^2$. A K3 is not rational,
so its irrationality is exactly two.

For an admissible primitive embedding,
requiring $H$ to remain algebraic gives
a Noether--Lefschetz divisor in
$\mathcal{F}_g$. To conclude, produce
infinitely many such embeddings with
$L,H$ nef in the same chamber and prove
that their divisors are distinct.
The varying determinants
$$
4g-4-a^2
$$
suggest a way to distinguish lattices,
but do not prove the embedding and
positivity claims. These lattice steps
were not transcribed and remain a gap
in this reconstruction.

A square-zero pencil $|E|$ instead
gives a genus-one fibration over
$\mathbb{P}^1$. A degree-two
multisection supplies a degree-two
linear series on its generic fibre,
leading to a rational map to a ruled
surface, birational to $\mathbb{F}_n$.
The pencil alone is not a degree-two
map of surfaces.

Relevant classical reference:
Saint-Donat, {\it Projective models
of K-3 surfaces}, American Journal
of Mathematics 96 (1974), 602--639,
Theorem 5.2 on hyperelliptic systems.
Its use is cross-checked in
Garbagnati--Sarti,
{\it Projective models of K3 surfaces
with an even set}, Section 4,
\href{https://sarti.pages.math.cnrs.fr/%
home/ottonodali1312.pdf}
{author manuscript}.
This does not verify the infinite-divisor
assertion.

\medskip\noindent
{\bf Higher gonality loci.}
We now relate irrationality to the
existence of higher-gonality curves.

The unpolarized locus
$G_{e;g}\subseteq\mathcal{F}_g$
consists of $(X,L)$ admitting a smooth
curve $C\subset X$ with
$$
\operatorname{gon}(C)=e,\qquad g(C)\geq2.
$$
The curve need not lie in a multiple
of the chosen polarization.

The polarized locus
$P^m_{e;g}\subseteq\mathcal{F}_g$
requires a smooth $C\in|mL|$ with
$$
\operatorname{gon}(C)=e.
$$
Here $m>0$. Smoothness is made explicit
to match the preceding definition;
if singular curves were intended, the
gonality convention needs checking.

\medskip\noindent
{\bf Example reported in the seminar
(unverified).}
For every
$$
g\equiv7\pmod{9},
$$
the locus $G_{3;g}$ has an irreducible
component of codimension two not
contained in $G_{2;g}$.
The congruence sign replaces the
approximation sign in the transcription.
The exact range of $g$ and the lattice
criterion were not recorded.

\medskip\noindent
{\bf Moretti example
(reported, not fully verified).}
A polarized K3 surface with
$$
\operatorname{Pic}(X)=\mathbb{Z}L,
\qquad L^2=10
$$
was said to have
$\operatorname{irr}(X)=3$ but no
smooth trigonal curve. Here $g=6$.
Thus a degree-three rational map need
not arise from a smooth trigonal curve.
Related source: Moretti,
{\it The polarized degree of
irrationality of K3 surfaces},
\href{https://arxiv.org/abs/2303.07289}
{arXiv:2303.07289}.
An exact theorem covering both parts
of this example remains to be checked.

\medskip\noindent
{\bf Final theorem
(reported, unverified).}
The loci $P^m_{e;g}$ are either empty,
all of $\mathcal{F}_g$, or finite
unions of Noether--Lefschetz divisors.
The exact hypotheses and reference
were not recorded. Check whether the
statement uses gonality equal to $e$
or at most $e$, and the loci themselves
or their closures. The announcement
confirms the union-of-divisors theme,
but not these finer distinctions.



\section{Determinantal surfaces in
$\mathbb{P}^3$ as Hodge loci, and quartic
threefolds}
\label{section-lozano-determinantal-nl}
%% speaker: César Lozano Huerta, UNAM Oaxaca
%% place: Online Seminar, Moduli Spaces of
%% place: Sheaves and Complexes (Zoom)
%% date: October 27, 2026, 12:00 (BRT)

\noindent
César Lozano Huerta, UNAM Oaxaca.
Online Seminar, Moduli Spaces of Sheaves
and Complexes.
October 27, 2026, 12:00 (BRT), Zoom.

\medskip\noindent
{\bf Abstract.}
Determinantal surfaces of degree $d$ in
$\mathbb{P}^3$ form families with several
irreducible components. I will discuss the
following theorem: each of these components
is also a Noether-Lefschetz component. Many
aspects of the global geometry of such
components remain mysterious. In one
dimension higher, however, things seem to
simplify: I will present results on the
geometry of the family of determinantal
quartic hypersurfaces in $\mathbb{P}^4$.

\medskip\noindent
{\bf Sources of these notes.} Screenshots of
the handwritten slides, an automatic
transcript of the Zoom audio (poor in
places), and the two papers below. Remarks
marked as spoken come from the transcript.

\medskip\noindent
{\bf Papers.}
\begin{enumerate}
\item M. Leal, C. Lozano Huerta,
M. Vité, \emph{The Noether--Lefschetz locus
of surfaces in $\mathbb{P}^3$ formed by
determinantal surfaces}, Math. Nachr. 297
(2024), arXiv:2303.09028.
\item M. Leal, C. Lozano Huerta,
M. Vité, \emph{Explicit birational geometry
of determinantal quartic 3-folds},
arXiv:2504.14461 (v3, June 2026).
Spoken: to appear in a Brazilian journal.
\end{enumerate}
The slides cite these as ``Leal--Vite--LH,
24'' and ``Leal--Vite--LH, 26''.

\medskip\noindent
{\bf Determinantal surfaces.} Take an
$n \times n$ matrix $A$ whose entries are
homogeneous polynomials in four variables,
with degrees chosen so that $\det A$ is a
homogeneous polynomial of degree $d$. The
surface
$$
Y = \{ \det A = 0 \} \subset \mathbb{P}^3
$$
is a \emph{determinantal surface}. To get
families, fix the degrees of the entries and
vary the coefficients.

\begin{example}
\label{example-lozano-3x3}
Let $l_1,\ldots,l_6$ be linear forms and
$q_1,q_2,q_3$ quadratic forms on
$\mathbb{P}^3$, and
$$
A =
\begin{pmatrix}
l_1 & l_2 & l_3 \\
l_4 & l_5 & l_6 \\
q_1 & q_2 & q_3
\end{pmatrix}.
$$
Then $\det A$ has degree $1+1+2 = 4$, so
$Y = \{ \det A = 0 \}$ is a quartic surface.
\end{example}

The degrees of the entries are recorded by
two vectors $a = (a_1,\ldots,a_n)$ and
$b = (b_1,\ldots,b_n)$: the entry in
position $(i,j)$ has degree $a_i - b_j$.
For instance
$$
a = (1,1,2),
\quad
b = (0,0,0)
$$
gives the matrix of degrees
$$
\begin{pmatrix}
1 & 1 & 1 \\
1 & 1 & 1 \\
2 & 2 & 2
\end{pmatrix},
$$
which is Example~\ref{example-lozano-3x3},
and $a = (2,2)$, $b = (0,0)$ gives a
$2 \times 2$ matrix of quadrics. The degree
of the surface is
$$
d = \sum_i a_i - \sum_j b_j.
$$
Pairs satisfying the needed conditions are
called \emph{admissible pairs of degree $d$};
the talk did not state the conditions.
Write $\operatorname{det}(a,b)$ for the
closure in $\mathbb{P}^N$, the space of
surfaces of degree $d$, of the set of
determinantal surfaces of type $(a,b)$.

\begin{example}
\label{example-lozano-five-families}
For $d = 4$ there are five families, with
matrices of degrees
$$
\begin{pmatrix}
1 & 1 \\
3 & 3
\end{pmatrix},
\quad
\begin{pmatrix}
1 & 2 \\
2 & 3
\end{pmatrix},
\quad
\begin{pmatrix}
2 & 2 \\
2 & 2
\end{pmatrix},
$$
the $3 \times 3$ matrix above, and the
$4 \times 4$ matrix of linear forms.
\end{example}

\begin{theorem}[Leal--Vité--Lozano Huerta]
\label{theorem-lvl-irreducible}
Fix an admissible pair $(a,b)$. Then
$\operatorname{det}(a,b)$ is irreducible,
and its dimension is given by an explicit
formula in $a$ and $b$.
\end{theorem}

The formula was not shown; spoken, it is
long and is in the first paper.

\medskip\noindent
{\bf Noether--Lefschetz.} For $d \ge 4$
the Noether--Lefschetz theorem says that a
very general surface $S \subset
\mathbb{P}^3$ of degree $d$ has
$\operatorname{Pic}(S) = \mathbb{Z}h$, with
$h$ the hyperplane class. The
\emph{Noether--Lefschetz locus}
$NL(d) \subset \mathbb{P}^N$ is the set of
smooth surfaces of degree $d$ with Picard
number at least 2. It is a countable union
of components $\Sigma$, and the component
records the lattice $\operatorname{Pic}(S)$.
Spoken, in answer to a question: there are
infinitely many components because
$H^2(S,\mathbb{Z})$ is huge (rank 55 already
for sextics, the speaker said) and a rank 2
lattice of Hodge classes can sit in it in
infinitely many ways. For example, the
general quartic containing a twisted cubic
lies in one component.

\begin{theorem}[Leal--Vité--Lozano Huerta]
\label{theorem-lvl-nl-component}
Fix an admissible pair $(a,b)$ of degree
$d$, and let
$\operatorname{det}(a,b)^{sm}$ be its open
subset of smooth surfaces. Then
\begin{enumerate}
\item $\operatorname{det}(a,b)^{sm}$ is a
component of $NL(d)$, and
\item a general $Y \in
\operatorname{det}(a,b)^{sm}$ has
$\operatorname{Pic}(Y) \cong \mathbb{Z}^2$.
\end{enumerate}
\end{theorem}

Spoken: each family has a smooth member, so
the smooth locus is non-empty. Entering the
determinantal locus is the first time the
Picard number jumps, and it jumps by the
least possible amount, from 1 to 2.

\medskip\noindent
{\bf Idea of the proof, for $d = 4$.} Take
$A$ of type
$a = (1,1,2)$, $b = (0,0,0)$ and a
$3 \times 2$ submatrix $A_0$ of $A$. The
$2 \times 2$ minors of $A_0$ cut out a curve
$C$. Every $2 \times 2$ minor of $A_0$ also
appears in the expansion of $\det A$, so $C
\subset S = \{ \det A = 0 \}$. The claim is
$$
\operatorname{Pic}(S) = \langle h, C
\rangle,
$$
with $C$ not a complete intersection, so
$S \in NL(4)$. Let $\Sigma^4_C$ be the
Noether--Lefschetz component of quartics
containing a curve of the class of $C$.
Then
$$
\operatorname{det}(a,b) \subset
\Sigma^4_C,
$$
and it remains to rule out that
$\Sigma^4_C$ is bigger. For this one shows
that the tangent spaces agree,
$$
T\Sigma^4_C = T\operatorname{det}(a,b),
$$
and that $\operatorname{det}(a,b)$ is
generically smooth; an irreducible
subvariety of the same dimension with equal
tangent spaces is then the whole component.
Spoken: the tangent space computation is
deformation theory of the pair, deforming
$S$ together with the line bundle
$\mathcal{O}_S(C)$.

\begin{remark}
\label{remark-lozano-minors}
The sentence ``every $2 \times 2$ minor of
$A_0$ appears in $\det A$'' is my
justification of $C \subset S$: expanding
$\det A$ along the remaining column shows
$\det A$ lies in the ideal of $2 \times 2$
minors of $A_0$. The slide boxes the
submatrix but does not say which columns.
\end{remark}

\medskip\noindent
{\bf Global geometry, $d = 4$.} Quartic
surfaces form
$$
\mathbb{P}^{34} = |\mathcal{O}_{\mathbb{P}^3}
(4)|.
$$
Each of the five families of
Example~\ref{example-lozano-five-families}
is a hypersurface in $\mathbb{P}^{34}$.

\begin{theorem}[Leal--Vité--Lozano Huerta]
\label{theorem-lvl-degrees}
The five hypersurfaces
$\operatorname{det}(a,b) \subset
\mathbb{P}^{34}$ of determinantal quartics
have degrees $320$, $2508$, $38475$,
$136512$ and $320112$.
\end{theorem}

The slide lists the degrees in this order;
the abstract of arXiv:2303.09028 lists
$38475$ and $136512$ the other way round, so
which degree goes with which family was not
checked. Spoken: to decide whether a given
quartic is determinantal of the first type,
plug its 35 coefficients into a polynomial
of degree 320; nobody expects to write such
polynomials down.

\medskip\noindent
{\bf Questions on the slides.} What is the
global geometry of
$\Sigma = \operatorname{det}(a,b)$? Does the
result generalize to $\mathbb{P}^4$?
Spoken: the authors expected no Hodge theory
in $\mathbb{P}^4$, which turned out wrong.

\medskip\noindent
{\bf Linear determinantal quartic
threefolds.} Let $A$ be a $4 \times 4$
matrix of linear forms on $\mathbb{P}^4$ and
$$
Y = \{ \det A = 0 \} \subset \mathbb{P}^4,
$$
a \emph{linear determinantal quartic}.
Spoken: a general $Y$ is singular, with 20
nodes, and its topology is as simple as
possible; this is why the surface theorem
cannot be copied verbatim. A small resolution
of the 20 nodes is a smooth threefold.

\begin{theorem}[Leal--Vité--Lozano Huerta]
\label{theorem-lvl-quartic-threefolds}
Let $C \subset \mathbb{P}^3$ be a general
curve of degree 10 and genus 11. Let
$\widetilde{X} = \operatorname{Bl}_C
\mathbb{P}^3$. Then there are a flop
$$
\widetilde{X} \dashrightarrow X,
$$
to a second blow-up
$X = \operatorname{Bl}_{C'}\mathbb{P}^3$ of
$\mathbb{P}^3$ along a curve $C'$, and small
contractions
$$
\widetilde{X} \to Y,
\quad
X \to Y
$$
onto the same linear determinantal quartic
$Y \subset \mathbb{P}^4$.
\end{theorem}

On the slide both blow-ups are written
$\operatorname{Bl}_C$; I wrote $C'$ for the
second curve, see
Remark~\ref{remark-lozano-residual-series}.
Spoken: $Y$ is the anticanonical model of
the blow-up, so no special linear system is
needed. The slide also draws the cones of
$X$ in the plane spanned by $H$ (pullback of
the hyperplane) and $E$ (exceptional
divisor): the ample cone contains $-K_X$,
and the movable cone $\operatorname{Mov}(X)$
is drawn with the ray $40H - 11E$ as a
boundary. Spoken: there are no other
birational models beyond those in the
diagram. The arXiv abstract states this as
four stable base locus chambers in the
effective cone of $X$, for $C$ general in
each of three families of curves.

\medskip\noindent
{\bf Moduli.} Classically, the intermediate
Jacobian of the blow-up is
$$
IJ(\operatorname{Bl}_C\mathbb{P}^3)
\cong \operatorname{Jac}(C),
$$
a principally polarized abelian variety of
dimension 11. Let
$$
\mathcal{W}^3_{10} =
\{ (C, g^3_{10}) \}
$$
be the space of genus 11 curves with a
linear series $g^3_{10}$, let
$\mathcal{M}^3_{11,10} \subset
\mathcal{M}_{11}$ be the locus of curves
that have one, and let
$$
\mathcal{M}^{det} =
\operatorname{det}(a,b) /
\operatorname{PGL}_5
$$
be the moduli of linear determinantal
quartic threefolds. There are maps
$$
\mathcal{M}^3_{11,10}
\leftarrow \mathcal{W}^3_{10}
\rightarrow \mathcal{M}^{det},
$$
the right one sending $(C,g^3_{10})$ to the
quartic $Y$ of
Theorem~\ref{theorem-lvl-quartic-threefolds}.
The question on the slide is whether
$$
\mathcal{M}^3_{11,10} \cong
\mathcal{M}^{det}.
$$
Spoken: this was the goal of the paper and
is not proved; the speaker suspects it is
true. A general curve of genus 11 has no
$g^3_{10}$, so $\mathcal{M}^3_{11,10}$ is a
proper subvariety of $\mathcal{M}_{11}$.

\begin{remark}
\label{remark-lozano-residual-series}
One slide marks the left map as $2:1$, a
later one marks the right map as $2:1$.
Riemann--Roch suggests why: on $C$ of genus
11, $\deg K_C = 20$, and for $D$ a
$g^3_{10}$,
$$
h^0(K_C - D) = h^0(D) - \deg D - 1 + g
= 4 - 10 - 1 + 11 = 4,
$$
so the residual $K_C - D$ is again a
$g^3_{10}$. Each curve thus has two such
series, which explains the $2:1$ map to
$\mathcal{M}^3_{11,10}$. Plausibly the two
sides of the flop are the embeddings by $D$
and by $K_C - D$, which would explain the
other $2:1$. This reading is mine and was
not checked.
\end{remark}

\medskip\noindent
{\bf Questions (spoken).} Asked how a
quartic gives a point in a moduli of
sheaves, the speaker described the cokernel
of $A$, a sheaf supported on $Y$; whether it
is a line bundle is unclear because $Y$ is
singular. The speaker's closing case was
that determinantal varieties are a concrete
way into moduli spaces of sheaves.

\section{The Hilbert polynomial of linked
projective spaces}
\label{section-de-leon-linked-spaces}
%% speaker: Felipe de León, IMPA
%% place: Algebra Seminar, IMPA, Sala 228
%% date: November 4, 2026, 15:30

\noindent
Felipe de Le\'on, IMPA. Algebra Seminar
(Semin\'ario de \'Algebra), IMPA. Wednesday 4
November 2026, 15:30, Sala 228.

\medskip\noindent
{\bf Abstract.} Linked projective spaces are
quiver Grassmannians of one dimensional
subrepresentations, introduced to study
limiting divisors in degenerations of linear
series. It is unknown whether every linked
projective space is a flat degeneration of
the small diagonal in the corresponding
product of projective spaces. Previous work
showed that they have the Chow class of this
diagonal. In this talk, I will present a
refinement: linked projective spaces have the
multivariate Hilbert polynomial of this
diagonal. The proof has three ingredients.
First, local Mustafin models yield simple
normal crossings and an inclusion and
exclusion formula. Second, intersections of
irreducible components are multiplicity free
varieties, with Hilbert polynomials encoded
by certain polytopes. Third, the underlying
representation determines a polymatroidal
tiling of a simplex. This tiling reduces the
remaining calculation to nerves of convex
covers, which are contractible by the nerve
theorem. This is a joint work with E.Esteves
and E.Vital.

\medskip\noindent
{\bf Upshot.} Take a family of curves (or
varieties) that degenerates to a special
fibre with several components, and a linear
series on the general fibre. One wants to
know what the divisors of the linear series
become in the limit. Since the line bundle
can be twisted by the components of the
special fibre, there is not one limit line
bundle but a whole lattice of them, indexed
by the vertices of a $\mathbb{Z}^n$-quiver,
with maps between them. The spaces of
sections form a quiver representation, a
\emph{linked net}. A \emph{linked projective
space} is the space of compatible choices of
one section, up to scalar, at every vertex.
It lives inside a product of copies of
$\mathbb{P}^r$, it has dimension $r$, and it
parametrizes the limit divisors. When the
degeneration really comes from a family, it
is a flat degeneration of the \emph{small
diagonal}
$\{(x,x,\dots,x)\}\cong\mathbb{P}^r$; the
question is whether this holds for every
linked projective space. The theorem is that
it at least looks like one numerically: it
has the same multigraded Hilbert polynomial
as the small diagonal. Earlier work of
Esteves and de Le\'on had shown the weaker
statement that it has the same Chow class.
The proof is a count. The linked projective
space is a union of components meeting like
normal crossings, so its Hilbert polynomial
is an alternating sum over intersections of
components. Each intersection is a
multiplicity free variety whose Hilbert
polynomial counts the lattice points of a
polytope. The polytopes tile a simplex, and
the alternating sum then reduces to Euler
characteristics of contractible nerves, which
are $1$. Spoken: the answer to the
degeneration question seems to be negative,
so the equality of Hilbert polynomials is the
best possible numerical statement (see the
end of these notes).

\medskip\noindent
{\bf Sources of these notes.} The abstract,
copied from the IMPA seminars page (assembled
in pieces, since the page loads it
dynamically). An automatic transcript of the
audio (Whisper), of very poor quality: the
formulas on the board are mostly lost. The
paper of Esteves and de Le\'on below, for
definitions and the earlier theorem. There
are no photographs. Statements marked
``spoken'' come from the transcript; what I
(Claude) reconstructed from the paper or
added is marked as such. Earlier notes on the
same objects, from a talk by the same
speaker: Section
\ref{16th-alga-meeting-2026-%
section-mustafin-varieties}.

\medskip\noindent
{\bf Reference.} E. Esteves and F. de Le\'on,
\emph{Polymatroidal tilings and the Chow
class of linked projective spaces},
arXiv:2508.14799 (August 2025). Checked:
title, authors and Theorem 4.5 read from
arXiv.

\medskip\noindent
{\bf $\mathbb{Z}^n$-quivers.} (Definitions as
in the paper; the talk gave a simplified
version.) A $\mathbb{Z}^n$-quiver is a quiver
$Q$ with a partition of its arrows into $n+1$
\emph{arrow types} $A_0,\dots,A_n$, such that
each vertex has exactly one outgoing arrow of
each type, any two vertices are joined by an
\emph{admissible} path (one that omits at
least one arrow type), and two paths from the
same vertex end at the same vertex exactly
when their numbers of arrows of each type
differ by a constant. Spoken: the basic
example has vertex set
$\mathbb{Z}^{n+1}/\mathbb{Z}(1,\dots,1)$,
with an arrow of type $i$ from $u$ to
$u+e_i$. For $n=1$ this is a line of vertices
with arrows in both directions; for $n=2$ a
triangular lattice with three arrow types.

\medskip\noindent
{\bf Linked nets.} A \emph{linked net} is a
representation $\mathfrak{V}$ of $Q$ in
finite dimensional vector spaces, with maps
$\varphi_\gamma$ along paths $\gamma$, such
that:
\begin{enumerate}
\item for admissible paths
$\gamma_1,\gamma_2$ with the same endpoints,
$\varphi_{\gamma_1}$ is a nonzero multiple of
$\varphi_{\gamma_2}$;
\item $\varphi_\gamma=0$ for every minimal
circuit $\gamma$ (a closed path with one
arrow of each type);
\item if $\gamma_1,\gamma_2$ are admissible
paths leaving the same vertex with no arrow
type in common, then
$\ker\varphi_{\gamma_1}\cap
\ker\varphi_{\gamma_2}=0$.
\end{enumerate}
It is \emph{exact} if
$\mathrm{Im}\,\varphi^{u_1}_{u_2}=
\ker\varphi^{u_2}_{u_1}$
for every pair of neighbours $u_1,u_2$
(vertices joined by a simple path, with at
most one arrow of each type); exactness
implies that all the $V_v$ have the same
dimension, say $r+1$. It is \emph{finitely
generated} if there is a finite set $H$ of
vertices such that every vertex receives a
surjection along an admissible path from some
vertex of $H$. Spoken and in the paper: a
nontrivial exact finitely generated linked
net has a minimum set of generators $H$,
which is convex, and $\varphi^v_u\neq 0$ for
$u,v\in H$.

\medskip\noindent
{\bf Linked projective spaces.}
$\mathbb{LP}(\mathfrak{V})$ is the quiver
Grassmannian of subrepresentations of pure
dimension $1$. Using $H$ it is identified
with
$$
\mathbb{LP}(\mathfrak{V})
=\{(s_v)_{v\in H}\in
\prod_{v\in H}\mathbb{P}(V_v) \mid
\varphi^v_w(s_v)\wedge s_w=0
\ \forall v,w\in H\},
$$
which gives it a scheme structure. Spoken:
this works because $H$ determines the whole
representation, even though $Q$ is infinite.

\medskip\noindent
{\bf The theorem.} Let $\mathfrak{V}$ be a
nontrivial exact finitely generated linked
net of dimension $r+1$, with minimum set of
generators $H$. Then
$\mathbb{LP}(\mathfrak{V})$ has the same
multigraded Hilbert polynomial as the small
diagonal of $\prod_{v\in H}\mathbb{P}^r$. My
remark (standard, not from the talk): the
small diagonal is $\mathbb{P}^r$, and
$\mathcal{O}(d_v)_{v\in H}$ restricts to
$\mathcal{O}(\sum_v d_v)$ on it, so the
theorem says
$$
\chi\left(
\mathcal{O}_{\mathbb{LP}(\mathfrak{V})}
(d_v)_{v\in H}\right)
=\binom{r+\sum_{v\in H}d_v}{r}.
$$
The earlier result (Theorem 4.5 of the paper,
checked) is the Chow class:
$$
[\mathbb{LP}(\mathfrak{V})]
=\sum_{q}\prod_{v\in H}h_v^{\,r-q(v)},
$$
the sum over integer vectors $q\geq 0$ on $H$
with $\sum_v q(v)=r$, where $h_v$ is the
hyperplane class from $\mathbb{P}(V_v)$. This
is also the class of the small diagonal.

\medskip\noindent
{\bf Motivation: limits of divisors
(spoken).} The setting comes from work of
Esteves, Santos and Vital (names from the
transcript, not checked). Let $\mathcal{X}\to
B$ be a flat family over a disc or the
spectrum of a DVR, with smooth total space
and smooth general fibre, and with special
fibre $X$ reduced, with irreducible
components $Y_0,\dots,Y_n$. On the general
fibre take a line bundle $L$ and a linear
series $V\subset H^0(L)$. Each section gives
a divisor; the closure of the family of these
divisors meets the special fibre in the limit
divisors, and the question is how to
parametrize them. The line bundle $L$ extends
to the total space, and since the $Y_i$ are
Cartier divisors, so do all twists
$\mathcal{L}\otimes\mathcal{O}(\sum_i
a_iY_i)$; twisting by the whole special fibre
changes nothing, so the extensions are
indexed by
$\mathbb{Z}^{n+1}/\mathbb{Z}(1,\dots,1)$.
These are the vertices of the
$\mathbb{Z}^n$-quiver. Twisting by one
component gives maps between extensions,
which are the arrows. Restricting to $X$
gives a representation by line bundles, and
the images of $V$ in their spaces of sections
give an exact finitely generated linked net.
A point of $\mathbb{LP}(\mathfrak{V})$ is a
choice of sections, compatible along the
arrows, and the intersection of their zero
loci is a divisor on $X$: this defines a map
from $\mathbb{LP}(\mathfrak{V})$ to the
Hilbert scheme of $X$ whose image is the set
of limit divisors, with reduced structure.
When $\mathfrak{V}$ comes from such a
degeneration, $\mathbb{LP}(\mathfrak{V})$ is
a flat degeneration of the small diagonal
(via Mustafin varieties). The question is
whether this holds for every exact finitely
generated linked net.

Spoken, near the end of this part (garbled):
first the authors computed the Chow class,
which is easier, and then the Hilbert
polynomial; afterwards they found that the
answer to the question is negative: there is
a linked projective space that does not lie
in the main component of the Hilbert scheme,
although it has the Hilbert polynomial of the
diagonal. Not checked; I could not find this
in writing.

\medskip\noindent
{\bf Sketch of the proof (spoken, heavily
garbled; reconstructed with the abstract).}

Step 1. Local structure. Each point of
$\mathbb{LP}(\mathfrak{V})$ has a
neighbourhood isomorphic to an open subset of
the special fibre of a Mustafin variety
(theorem of Esteves and de Le\'on, see
Section
\ref{16th-alga-meeting-2026-%
section-mustafin-varieties}).
So the irreducible components $Y_1,\dots,Y_m$
of $\mathbb{LP}(\mathfrak{V})$ meet like
simple normal crossings, and the Hilbert
polynomial is given by inclusion and
exclusion:
$$
P_{\mathbb{LP}(\mathfrak{V})}
=\sum_{\emptyset\neq I}(-1)^{|I|+1}P_{Y_I},
\qquad Y_I=\bigcap_{i\in I}Y_i .
$$

Step 2. Components and their intersections.
My reading of the paper, not checked in
detail: the irreducible components are
indexed by vertices of $H$: for $v\in H$ the
component is the closure of the image of the
rational map
$\mathbb{P}(V_v)\dashrightarrow\prod_{u\in
H}\mathbb{P}(V_u)$,
$s\mapsto(\varphi^v_u(s))_u$. Spoken: an
intersection of components over a subset of
$H$ is nonempty exactly when the vertices are
pairwise neighbours, and it is again the
image of a rational map from a product of
projective spaces. Such images are
multiplicity free varieties, whose
multigraded Hilbert polynomials are sums of
products of binomial coefficients indexed by
the lattice points of a polytope (spoken:
computed by earlier authors, names lost in
the transcript; Li's work is cited in the
paper for the Chow classes). The polytope is
a truncation of a polymatroid, a standard
construction in combinatorics.

Step 3. Tiling. The polytopes attached to the
vertices of $H$ (base polytopes of
polymatroids defined from the kernels of the
maps $\varphi^v_u$) tile a standard simplex
(Theorem 4.2 of the paper). Reorganizing the
inclusion and exclusion sum by lattice points
of the simplex, the coefficient of each point
is the Euler characteristic of a simplicial
complex, the nerve of a cover of a convex set
by convex pieces. By the nerve theorem this
complex is contractible, so its Euler
characteristic is $1$. Spoken: the remaining
sum is evaluated with the Chu--Vandermonde
identity and the hockey stick identity, and
gives the Hilbert polynomial of the small
diagonal.

\medskip\noindent
{\bf Questions.} Spoken: describing the
irreducible components becomes much harder
for quiver Grassmannians of higher
dimensional subrepresentations, which is the
natural next case.


\section{Slopes of irregular fibrations
over surfaces}
\label{section-barja-irregular-fibrations}
%% speaker: Miguel A. Barja, Universitat
%% speaker: Politècnica de Catalunya
%% place: Seminário de Álgebra, IMPA, room 228
%% date: November 11, 2026, 15:30

\noindent
Miguel A. Barja, Universitat Polit\`ecnica
de Catalunya.
Semin\'ario de \'Algebra, IMPA.
November 11, 2026, 15:30, room 228.
Joint work with Mart\'i Lahoz and
Andr\'es Rojas.

\medskip\noindent
{\bf Abstract.}
In the first part of the talk, I will
introduce problems and results on the
geography of irregular varieties and
varieties fibred over curves. Then we will
study varieties fibred over surfaces, where
a new invariant appears---the ``positive
second Chern character''---explaining its
construction, properties, and relation to
invariants in lower dimensions. Finally, I
will present a theorem concerning the
geography of these fibrations in the case
of irregular varieties. This is joint work
with Mart\'i Lahoz and Andr\'es Rojas
(University of Barcelona).

\medskip\noindent
{\bf Sources of these notes.}
Daniel's 42 photographs, from
\texttt{IMG\_4381} to \texttt{IMG\_4422},
two English recordings of about 46 and
25 minutes, and his written comments.
The recordings were machine-transcribed
with faster-whisper small; names and
formulas in the raw transcripts are often
wrong. The photographs, including
progressive versions of slides, determine
the displayed formulas below. Published
references were checked where specified.
The new joint results and their proof
outline are recorded from the talk;
they have not been checked against a paper.
Daniel's comments are preserved separately
at the end, with labelled editorial answers.

\medskip\noindent
{\bf The problem in one paragraph.}
A line bundle $L$ on a variety $X$ has a
volume measuring the asymptotic growth of
its sections. For a fibration
$f\colon X\to T$, one wants to compare
that volume with a numerical invariant of
$f_*L$. The question changes with the
dimension of $T$: over a point one counts
sections, over a curve one uses a positive
degree, and over a surface the talk
proposes a positive second Chern character.
The numerator remains the volume on $X$;
the new invariant is on the base side.

\medskip\noindent
{\bf Geography and the absolute case.}
Let $X$ be an integral projective variety
of dimension $n$, and $L$ a line bundle.
Here integral means reduced and
irreducible. Recall that
$$
\operatorname{vol}_X(L)
=\limsup_{k\to\infty}
\frac{n!h^0(X,L^{\otimes k})}{k^n}.
$$
It is nonnegative, and equals $L^n$ when
$L$ is nef. A geography problem asks for
relations between numerical invariants,
and for the geometric meaning of equality.
For positive base invariant $\psi(f_*L)$,
one studies the ratio
$$
s_f(L)=
\frac{\operatorname{vol}_X(L)}{\psi(f_*L)}.
$$
The inequalities themselves still make
sense when the denominator is zero.

Over a point, $f_*L$ is a vector space,
so the relevant invariant is $h^0(X,L)$.
The slides recall Clifford's inequality
for a line bundle on a smooth curve of
genus at least two, with
$0\leq\deg L\leq 2g-2$:
$$
\deg L\geq 2(h^0(X,L)-1).
$$
They also recall Noether's inequality for
a smooth minimal surface of general type,
$$
K_X^2\geq 2p_g(X)-4,
$$
and the threefold analogue
$$
\operatorname{vol}_X(K_X)
\geq \frac43p_g(X)-\frac{10}{3}.
$$
The last formula is transcribed from the
slide, which lists several contributors;
its precise generality and attribution
have not been independently checked here.

\medskip\noindent
{\bf Fibrations over curves.}
For a relatively minimal fibration
$f\colon X\to B$ from a smooth surface
to a smooth curve, put $b=g(B)$ and let
$g\geq2$ be the genus of a general fibre
$F$. In characteristic zero, the classical
slope inequality is
$$
K_{X/B}^2
\geq \frac{4(g-1)}{g}
\deg f_*\omega_{X/B}.
$$
The numerical quantities are
$$
\begin{aligned}
\deg f_*\omega_{X/B}
&=\chi(\omega_X)
-(b-1)(g-1),\\
K_{X/B}^2
&=K_X^2-8(b-1)(g-1).
\end{aligned}
$$
Checked: the inequality and its classical
attribution to Xiao and Cornalba--Harris
are reviewed in the introduction to
\cite{zhang-slope-surfaces}, equation (1.1).
The discussion of positive characteristic
in the talk emphasizes that the analogous
raw invariants can be negative and that
there are counterexamples to unrestricted
extensions. The slide attributes these to
Gu--Sun--Zhou; this recent attribution has
not been checked here.

The proposed replacement is to compare
$\operatorname{vol}_X(L)$ with
$\deg^+ f_*L$. For a curve bundle, this
second invariant has a concrete
Harder--Narasimhan description.

\begin{definition}
\label{definition-barja-positive-degree}
Let $E$ be a vector bundle on a smooth
projective curve $B$, with
Harder--Narasimhan filtration
$$
0=E_0\subset E_1\subset\cdots
\subset E_r=E.
$$
Write $G_i=E_i/E_{i-1}$ and
$$
\mu_i=\frac{\deg G_i}
{\operatorname{rk}G_i},
\qquad \mu_1>\cdots>\mu_r.
$$
Let $m$ be the last index with
$\mu_m\geq0$, taking $m=0$ if none exists.
The \emph{positive part},
\emph{positive degree}, and
\emph{positive rank} are
$$
\begin{aligned}
E^+&=E_m,\\
\deg^+ E&=\deg E^+,\\
\operatorname{rk}^+E
&=\operatorname{rk}E^+.
\end{aligned}
$$
\end{definition}

Checked: this is the nonnegative-slope
convention of
\cite[Definition 2.2]
{barja-irregular-slopes}.
The slide mixes a strict inequality with
$\mu_m\geq0$; retaining the zero-slope
pieces is necessary for the stated
section-preservation property.
It does not change $\deg^+ E$.

The positive part has maximal degree
among subsheaves of $E$, and
$$
H^0(B,E^+)=H^0(B,E).
$$
In characteristic zero it is the largest
nef subbundle. Checked: for the direct
images used in the talk, see
\cite[Lemma 2.3]{barja-irregular-slopes}.
On $\mathbb{P}^1$, the nonnegative
summands are exactly the part generated
by global sections.

Formally twisting the slopes by a real
number $x$ gives the convex piecewise
linear function
$$
\deg^+_E(x)
=\sum_i\operatorname{rk}(G_i)
\max\{\mu_i+x,0\}.
$$
Away from its break points, its derivative
is the rank of the positive part after
twisting. At a break point, including the
zero-slope factors affects the rank and
requires a choice of one-sided derivative.

\medskip\noindent
{\bf Irregular varieties and twisting.}
Let $a\colon X\to A$ be a map to an
abelian variety, and write
$\widehat A=\operatorname{Pic}^0(A)$.
One says that $X$ has maximal
$a$-dimension when
$$
\dim a(X)=\dim X.
$$
For a fibration with general fibre $F$,
maximal $a$-dimension of the fibration
means $\dim a(F)=\dim F$.
These two conditions should be
distinguished.

The continuous rank is the generic,
equivalently minimal, value
$$
h_a^0(X,L)=
\min_{P\in\widehat A}
h^0(X,L\otimes a^*P).
$$
For $f\colon X\to B$, the continuous
positive degree is the generic value of
$$
\deg^+f_*(L\otimes a^*P),
$$
denoted $\deg_a^+f_*L$.
Checked:
\cite[Definition 2.4]
{barja-irregular-slopes}.
For smooth $X$ of maximal $a$-dimension,
generic vanishing gives
$$
h_a^0(X,\omega_X)=\chi(\omega_X).
$$
For the relative canonical bundle, the
corresponding degree is
$$
\deg_a^+f_*\omega_{X/B}
=\chi(\omega_X)
-\chi(\omega_B)\chi(\omega_F).
$$
Checked: see
\cite[Example 2.5]{barja-irregular-slopes}.
The relative canonical bundle is essential
in this formula; a slide writes $K_X$
instead of $K_{X/B}$.

For maximal $a$-dimension, the talk recalls
Clifford--Severi inequalities of the form
$$
\operatorname{vol}_X(L)
\geq n!h_a^0(X,L),
$$
with $2n!$ in the subcanonical case.
Checked in the smooth nef case:
\cite[Main Theorem]{barja-clifford-severi}.
For fibrations over curves, the analogous
bounds are
$$
\operatorname{vol}_X(L)
\geq n!\deg_a^+f_*L
$$
and $2n!$ for numerically subcanonical
$L$, with the hypotheses of
\cite[Corollaries 4.15 and 4.16]
{barja-irregular-slopes}; those corollaries
assume $h_a^0(X,L)\neq0$.
Numerically subcanonical means that
$K_X-L$ is pseudo-effective, in the
normal $\mathbb{Q}$-Gorenstein setting.

\medskip\noindent
{\bf The surface-base problem.}
Now let $f\colon X\to S$ have a smooth
surface as base. A degree with respect to
an auxiliary ample divisor only sees the
first Chern character. For families of
curves, the dimensionally natural candidate
is instead the number
$\operatorname{ch}_2(f_*\omega_{X/S})$.
Throughout this section Chern characters
in top degree are evaluated on the base.

\begin{theorem}[Zhang]
\label{theorem-barja-zhang}
Let $f\colon X\to S$ be a flat,
Cohen--Macaulay, relatively minimal
fibration from a normal projective
threefold to a smooth projective surface,
with connected fibres and general fibre
a smooth curve of genus $g\geq2$.
In characteristic zero,
$$
K_{X/S}^3\geq
\frac{12(g-1)}{g+1}
\operatorname{ch}_2(f_*\omega_{X/S}).
$$
\end{theorem}

Checked:
\cite[Definition 1.1 and Theorem 1.3]
{zhang-slope-surfaces}.
Here relatively minimal means that
$K_{X/S}$ is nef. The published theorem
also has a positive-characteristic case
under an extra semipositivity hypothesis.
It does not assert that the second Chern
character on the right is nonnegative.

\begin{remark}
\label{remark-barja-chern-character}
The relevant invariant is the second
Chern \emph{character}, not $c_2$:
$$
\operatorname{ch}_2(E)
=\frac12c_1(E)^2-c_2(E).
$$
For curves over an $m$-dimensional base,
the analogous candidate is
$\operatorname{ch}_m(f_*\omega_{X/T})$.
If the total space has dimension $n$, this
index is $n-1$, not $n$.
This dimensional proposal is discussed
after Theorem 1.3 of
\cite{zhang-slope-surfaces}.
\end{remark}

\medskip\noindent
{\bf The positive second Chern character.}
Fix a polarization $H$ on a smooth complex
projective surface $S$. There are two
slopes involved:
$$
\mu_H(E)=\frac{Hc_1(E)}
{\operatorname{rk}E},
\qquad
\nu_{0,0}(E)=
\frac{\operatorname{ch}_2(E)}{Hc_1(E)}.
$$
The first is the usual sheaf slope.
The second belongs to a weak stability
condition on the tilted heart
$\operatorname{Coh}^0(S)$ inside
$D^b(S)$. When the denominator of the
second slope is zero, one uses the
infinite-slope convention of weak
stability, rather than ordinary division.

An object $E$ of this heart has ordinary
cohomology only in degrees $-1$ and $0$.
Its degree $-1$ sheaf has all ordinary
HN slopes at most zero; its degree $0$
sheaf has all ordinary HN slopes strictly
positive, allowing torsion.
Checked:
\cite[Section 2, Remark 2.5]
{lahoz-rojas-chern-degree}.
This is an instance of the stability
framework recalled in Section
\ref{algebraic-geometry-section-stability}.

\begin{definition}
\label{definition-barja-positive-chern}
For $E\in\operatorname{Coh}^0(S)$, take
its HN filtration for $\nu_{0,0}$.
Let $E^+$ be its initial segment whose
HN factors have strictly positive
$\nu_{0,0}$-slope. The
\emph{positive second Chern character} is
$$
\operatorname{ch}_2^+(E)
=\operatorname{ch}_2(E^+).
$$
For a general sheaf $E$, first take its
degree-zero cohomology for the tilted
heart, and apply this construction.
\end{definition}

The notation suppresses possible dependence
on $H$. Spoken, second recording near
00:14: Barja said that they did not have
examples demonstrating actual dependence
on the polarization.
Checked: this is the degree-zero Chern
degree function at zero of
\cite[Definition 3.1]
{lahoz-rojas-chern-degree}.
For a sheaf, the preceding tilted
cohomology is its ordinary HN part with
strictly positive $\mu_H$-slopes.
It is distinct from the curve convention
that retains the zero-slope factors.

A key point of the slides is that
$E^+\to E$ is a subobject in the tilted
heart. It need not be an injection of
ordinary sheaves. If $E$ is a sheaf in
that heart and $Q$ is its tilted quotient,
one has an ordinary exact sequence
$$
0\to\mathcal H^{-1}(Q)\to E^+
\to E\to\mathcal H^0(Q)\to0.
$$
The degree $-1$ term explains why an
ordinary-subsheaf picture is insufficient.
The slides also state that $E^+$ is
locally free when $E$ is locally free.
This last property is recorded from the
talk, without an independent proof here.

Rational twists are treated using
$$
\operatorname{ch}^{\beta}(E)
=e^{-\beta H}\operatorname{ch}(E),
\qquad \beta\in\mathbb{Q},
$$
and the corresponding tilted hearts and
slopes. This is a formal twisted Chern
character; a fractional tensor power is
not required to be a genuine line bundle.

\begin{theorem}[Lahoz--Rojas]
\label{theorem-barja-chern-functions}
The Chern degree functions on a smooth
complex polarized projective surface
extend continuously to real arguments.
At every rational argument they have a
polynomial expression on each side.
On an abelian surface they coincide with
the cohomological rank functions.
\end{theorem}

Checked:
\cite[Theorems A and D]
{lahoz-rojas-chern-degree}.
The paper uses argument $x=-\beta$ for a
Chern character twisted by $-\beta H$.
These notes use the paper's convention;
the talk's twist parameter has the
opposite sign.
In particular, for an abelian surface,
$$
\operatorname{ch}_2^+(E)
=h^0_{\operatorname{id}}(S,E).
$$
The slides further state a comparison
with continuous rank when $S$ maps
finitely to an abelian variety and $H$
is pulled back from it. That comparison
is recorded as part of the talk and has
not been checked in a published source.

For $f\colon X\to S$ and $a\colon X\to A$,
the talk states that
$$
\operatorname{ch}_2^+
\bigl(f_*(L\otimes a^*P)\bigr)
$$
is constant for $P$ in a nonempty open
subset of $\widehat A$, in the stated
maximal $a$-dimension setting.
Its generic value is denoted
$\operatorname{ch}_{2,a}^+(f_*L)$.
This generic-constancy assertion is
transcribed from the slides, not
independently verified here.

\medskip\noindent
{\bf The new inequalities.}

\begin{theorem}[Barja--Lahoz--Rojas,
as stated in the talk]
\label{theorem-barja-surface-slope}
Let $X$ be a complex integral projective
variety of dimension $n$, and let
$a\colon X\to A$ be a generating map to
an abelian variety, with $X$ of maximal
$a$-dimension.
Let $L\in\operatorname{Pic}(X)$,
and let $f\colon X\to S$ be a flat
fibration onto a smooth projective
surface with fixed polarization $H$.
Then
$$
\operatorname{vol}_X(L)
\geq n!\operatorname{ch}_{2,a}^+(f_*L).
$$
If $L$ is subcanonical, the announced
bound is
$$
\operatorname{vol}_X(L)
\geq 2n!\operatorname{ch}_{2,a}^+(f_*L).
$$
\end{theorem}

This is transcribed from the photographed
main-theorem slide. No paper containing
this exact statement was located.
The slide does not spell out the
singularity hypotheses implicit in the
word subcanonical; no additional
singular generalization is asserted here.
The generating-map terminology is retained
from the slide.

The next slide announces a refinement of
Theorem~\ref{theorem-barja-zhang}, under
its characteristic-zero hypotheses:
$$
K_{X/S}^3\geq
\frac{12(g-1)}{g+1}
\operatorname{ch}_2^+
\bigl(f_*\omega_{X/S}\bigr).
$$
This refinement is also a talk statement,
unchecked against a paper. It uses the
untwisted positive invariant, not the
continuous invariant indexed by $a$.

\medskip\noindent
{\bf Proof strategy from the slides.}
This is an outline of the presentation,
not a complete proof of the new theorem.

\begin{enumerate}
\item Spread the complex data over a
finitely generated $\mathbb{Z}$-algebra
$R\subset\mathbb{C}$. The original
fibration is recovered from the generic
fibre by extending scalars to
$\mathbb{C}$. Reduce at closed points.
Spread the chosen
positive part and the cohomology sheaf of
its tilted quotient along with $X,S,L$.
Thus a characteristic-zero construction
is retained in a model before reducing.
\item Control the volume under reduction.
The slide states that the complex volume
is the infimum of the geometric
special-fibre volumes over any dense set
of closed points of the spread.
Its attribution is Kim and
Barja--Lahoz--Rojas, based on Jiao.
The exact reference was not identified.
\item After restricting the spread,
cohomology dimensions and the numerical
second Chern character of the spread
positive part agree with their values
in characteristic zero.
The slide writes $X$ in a cohomology
expression for a sheaf on $S$; the base
surface is intended.
\item Apply iterated Frobenius on the
reduced base surface. Its degree is
$p^{2e}$, the scale appearing in the
volume comparison. A suitable second
cohomology error is
$o(p^{2e})$. Divide by $p^{2e}$ and let
$e$ grow; then use the volume comparison
to return to characteristic zero.
\end{enumerate}

Spoken, second recording near 14:00:
equality examples are understood when the
base is an abelian surface, where the
invariant is a continuous section count.
Barja said that computing the invariant
and understanding equality on other
surfaces remain difficult.

Spoken, in the final discussion: an
analogue over a threefold base would need
a replacement for the positive third
Chern character. Barja described a
candidate via asymptotic section counts
after Frobenius base change in positive
characteristic. He did not give a
characteristic-zero construction or a
proved higher-base-dimensional theorem.

The final slide explicitly says that this
proof does not work in positive
characteristic. It also announces variants
using the geometry of $a$, including
nonmaximal $a$-dimension, and a general
principle transferring suitable absolute
inequalities in sufficiently large
characteristic to positive-Chern slope
inequalities in characteristic zero.
These are recorded as the speaker's final
remarks, without a proof here.

\medskip\noindent
{\bf Daniel's notes, with answers.}
The quoted comments below are Daniel's.
The answers are editorial explanations by
ChatGPT, not quotations from the speaker
or independently reviewed proofs.

\begin{remark}
\label{remark-barja-daniel-positivity}
{\bf Daniel's question.}
\begin{quote}
In slide 11 he told that he cares only
about \emph{positive} invariants.
Like---why? What's wrong with those little
numbers being negative?
\end{quote}
{\bf Editorial answer.}
Nothing makes a negative invariant
illegitimate. Here the aim is a lower
bound for a nonnegative volume in terms
of another measure of positivity.
If $\lambda>0$ and $\psi<0$, then
$\operatorname{vol}(L)\geq\lambda\psi$
is automatic and gives no useful bound.
Also, dividing an inequality by a negative
number reverses it. Positive parts retain
the contribution the desired slope is
meant to measure. They do not remove
negative invariants from geometry.
\end{remark}

\begin{remark}
\label{remark-barja-daniel-main-problem}
{\bf Daniel's formulation.}
\begin{quote}
Main problem: what is the right invariant
to play the role of volume when we study
fibrations over surfaces instead of over
curves?
\end{quote}
{\bf Editorial answer.}
The volume remains on the total space.
The missing invariant replaces the degree
of the direct image on the base. The
proposed surface-base replacement is
$\operatorname{ch}_2^+$, or its generic
twisted value in the irregular setting.
\end{remark}

\begin{remark}
\label{remark-barja-daniel-zhang}
{\bf Daniel's note.}
\begin{quote}
Zhang's theorem: the inequality contains
the volume $K^3$ and also the
\emph{second} Chern class.
\end{quote}
{\bf Editorial answer.}
The formula uses $K_{X/S}^3$ and the second
Chern character of $f_*\omega_{X/S}$.
The first equals the relative canonical
volume under the nef hypothesis.
The character is the combination of
$c_1^2$ and $c_2$ in
Remark~\ref{remark-barja-chern-character}.
\end{remark}

\begin{remark}
\label{remark-barja-daniel-index}
{\bf Daniel's note.}
\begin{quote}
In fact, the natural generalization for a
variety of dimension $n$ is the $n$th
Chern character. However, we don't know
if $\operatorname{ch}_2
(f_*\omega_{X/S})$ is nonnegative.
\end{quote}
{\bf Editorial answer.}
Use the dimension of the base as the
index. For a total $n$-fold fibred in
curves it is $n-1$.
The slides raise nonnegativity of the raw
second Chern character as a difficulty;
it is not a hypothesis-free positivity
assertion in Zhang's theorem.
\end{remark}

\begin{remark}
\label{remark-barja-daniel-wrong-invariant}
{\bf Daniel's conclusion.}
\begin{quote}
In fact, the second Chern character is
\emph{not} the right invariant for this
work.
\end{quote}
{\bf Editorial answer.}
More precisely, the talk replaces the raw
second Chern character by a positive part
constructed using stability. The raw
character remains central to the
construction and to Zhang's inequality.
\end{remark}

\begin{remark}
\label{remark-barja-daniel-hn}
{\bf Daniel's note.}
\begin{quote}
degree for curves: take the most
positive-degree chunk in the
Nardar-Hara\-shiman filtration of the line
bundle on a curve
\end{quote}
{\bf Editorial answer.}
The name is Harder--Narasimhan.
It is the vector bundle $E=f_*L$ on the
base curve that is filtered. Keep the
whole initial segment with nonnegative
HN slopes, not just the first or highest
slope piece. A line bundle itself has
only one HN factor.
\end{remark}

\begin{remark}
\label{remark-barja-daniel-integration}
{\bf Daniel's note.}
\begin{quote}
It looks like he has defined the
``positive degree'', which recovers the
second chern character by integration.
\end{quote}
{\bf Editorial answer.}
You did hear an integration relation.
Spoken, second recording at 02:54--03:22:
Barja relates the surface invariant to
integration of the positive degree on a
curve, just as the curve invariant relates
to section counts over a point.
The photographed slides do not give the
precise integration formula, and it cannot
be safely reconstructed from this audio.
The definition actually written on the
slides is the stability construction
above. Separately, an ordinary numerical
second Chern character is the integral of
its cohomology class over the surface.
These two uses of integration should not
be conflated. No unverified integral
formula is asserted here.
\end{remark}

\begin{remark}
\label{remark-barja-daniel-characteristic}
{\bf Daniel's correction.}
\begin{quote}
They cannot prove the theorem in
characteristic zero. No, wait, I think
it's the other way around (see final
remarks): the proof does not work in
positive characteristic.
\end{quote}
{\bf Editorial answer.}
Your correction is confirmed by the final
slide: the announced result is in
characteristic zero. Frobenius is used
after spreading out and reducing modulo
primes as a tool in that proof. This does
not establish the same theorem over an
arbitrary field of positive
characteristic.
\end{remark}

\bibliography{my}
\bibliographystyle{amsalpha}

\section{About simplicity of the Cremona
group}
\label{section-simplicity-Cremona-group}
%% speaker: Sokratis Zikas, IMPA
%% date: February 27, 2026

\medskip\noindent
\textbf{Título.}
On simplicity of the Cremona groups.

\medskip\noindent
\textbf{Palestrante.}
Sokratis Zikas (IMPA).

\medskip\noindent
\textbf{Data.}
27 Fevereiro 2026.

\medskip\noindent
\textbf{Resumo.}
The Cremona group or rank $n$ is the group of
birational transformations of the projective
space of dimension $n$.
One of the central questions regarding these
groups, dating all the way back to the 19th
century and settled only recently, is whether
they are simple.
In this talk I will review the history of the
problem, as well as give a brief overview of
the theory of elementary relations of
Blanc-Lamy-Zimmermann, that led to the proof
of the non-simplicity of higher-rank Cremona
groups.
We will also see how this theory was later
used to prove even more exotic properties of
these groups.

Recall that the {\it Cremona group} is
defined by
$$
\text{Cr}_n(k)=
\{f:\mathbb{P}^n \dashrightarrow
\mathbb{P}^n:
f\text{ is birational}\}.
$$
Question by Enriques (1865):
is $\text{Cr}_n(k)$ simple?

We know that $\text{Cr}_n(k)$ is not simple
in the following cases:
\begin{enumerate}
\item $n=2$, $k$ any field.
Contat-Lamy '13, Lonjou '16.

\item $n=2$, $k$ perfect field
such that there exists a Galois orbit of size
8.
Lamy-Zimmermann '20.

\item $n \geq 3, k=\mathbb{C}$.
Blanc-Lamy-Zimmermann '21.
\end{enumerate}

[The second and third item they use a
different method
to define the automorphisms.]

\begin{theorem}
\label{theorem-cr3}
There is an isomorphism
$\text{Cr}_3(\mathbb{C}) \simeq
G \ast (\ast_J \mathbb{Z}/2\mathbb{Z})$
where $J$ is uncountable.
\end{theorem}

With corollaries:

\begin{corollary}
\label{lemma-corollary}
$\text{Aut}(\text{Cr}_3(\mathbb{C}))$
is not generated by inner automorphisms
nor automorphisms of the field $\mathbb{C}$.
\end{corollary}

Déserti '06:
$\text{Aut}(\text{Cr}_2(\mathbb{C}))$
is generated by such automorphisms.

\medskip\noindent
[We shall prove Theorem \ref{theorem-cr3}.]

Define a groupoid by
$$
\underbrace{\text{BirMori}(\mathbb{P}^n)}
_{\text{Cr}_n(\mathbb{C})}
=
\left\{f:X \dashrightarrow Y
:\substack{X,Y \text{ Mfs}, X,Y
\sim_{\text{bir}}\mathbb{P}^n \\
f\text{ birational}}\right\}.
$$
\begin{definition}
\label{definition-rank-r-fibration}
Let $r \geq 1$.
A morphism  $\eta:X \to B$ is a {\it rank $r$
fibration}
if
\begin{enumerate}
\item $\dim X > \dim B, \rho(X)-\rho(B) = r$.
\item $X$ is a MDS over $B$
\item $-K_X$ is $\eta$-big.
\item The result of any
$D$-MMP over $B$ is
$\mathbb{Q}$-factorial is terminal.
\item $\exists \Delta_B \geq 0$
such that $(B,\Delta_B)$ is kld.
\end{enumerate}
\end{definition}

In the first case there exists a 2-ray game.
In case (2) the 2-ray game ends.
In the last 3 cases the 2-ray game
stays in the category of MMP.

\begin{itemize}
\item Rank 1 fibration = Mfs.
\item Rank 2 fibration = Sarkisov link.
\end{itemize}

\begin{theorem}[BLZ]
\label{theorem-blz2}
The relations among Sarkisov links are
generated by
by the dominated relations of rank 3
fibrations.
\end{theorem}

$$
\text{BirMori}=
\left\langle{}
\substack{\text{Sarkisov} \\ \text{links}}
:\substack{\text{relations given by} \\
\text{rank-3 fibrations}}\right\rangle{}.
$$

Let $(g,d) \in
\{(2,8),(6,9),(10,10),(11,14)\}$,
$C \in U_{g,d}\subseteq H_{g,d}^{\leq}$,
$J=\bigcup U_{g,d}$,
$X=Bl_C \mathbb{P}^3 \dashrightarrow
\mathbb{P}^3$.

\begin{proposition}
\label{proposition-rank-2}
Let $X \to \mathbb{P}^3 \to
\Spec(\mathbb{C})$
be a rank-2 fibration.
$$
\xymatrix{
X\ar@{-->}[rr]\ar[d]
&&X\ar[d]\\
\ar[dr]\mathbb{P}^3\ar@{-->}[rr]^{\chi_C}
&&\mathbb{P}^3 \ar[dl]\\
&\Spec(\mathbb{C}).
}
$$
with $\chi_C^2=\text{id}$.
\end{proposition}

$$
\xymatrix{
X\ar[dd]\ar[dr]&&& X\ar[dl]\ar[dd]\\
&Z \ar[r]^{t \mapsto -t}& Z\\
\mathbb{P}^3 \ar@{-->}[rrr]&&&\mathbb{P}^3.
}
$$
\begin{proposition}
\label{proposition-no-rank-3}
There are no rank 3 fibrations
dominated $X \to \Spec(\mathbb{C})$.
\end{proposition}








\clearpage

\section{Automorphisms of quartic surfaces
and Cremona transformations}
\label{section-quartic-automorphisms}
%% speaker: Carolina Araujo
%% place: IMPA
%% date: September 24, 2024

\medskip\noindent
\textbf{Speaker.}
Carolina Araujo.

\medskip\noindent
\textbf{Place.}
IMPA.

\medskip\noindent
\textbf{Date.}
24 September, 2024.

\medskip\noindent
\textbf{Motivation.}

$X$ a smooth hypersurface of degree $d $,
$X\subset \mathbb{P}^{n+1}$. We want to
understand the group $\operatorname{Aut}(X)$.
These are invertible polynomial maps from $X$
to itself. $X$ is defined by a single
polynomial equation of degree $d$.

\begin{theorem}[Matsumura-Mousky,
1964]\leavevmode
	Except in two special cases, all
	automorphisms of $X$ come from automorphisms
	of the ambient space:
	
	If $(n,d)\neq (1,3),(2,4)$, there is a
	surjective map
	$$
	\operatorname{Aut}(\mathbb{P}^{n+1},X)
	\overset{\pi}{\longrightarrow}
	\operatorname{Aut}(X)
	$$
	where $\operatorname{Aut}(X,Y)$ means
	automorphisms of $X$ that stabilize  $Y$.
\end{theorem}

\medskip\noindent
\textbf{Exceptional cases.}

Let's look at the exceptionals cases.

\begin{enumerate}
	\item $(n,d)=(1,3)$. In this case
	$C=X_3\subset \mathbb{P}^2$ is an elliptic
	curve and we have
		$\operatorname{Aut}(C)$ is a
		semidirect product of $C$ and
		$\mathbb{Z}_m$, with $m=2,4,6$.
and
 $$\operatorname{Aut}(\mathbb{P}^2,C)\text{
 is finite.} $$

 Elements in $C$ are translations of the
 torus. The {\it translation of
 $x$ with respect to $p$ and $(x:y:z)$} is
 done by intersecting the curve with the line
 that joins $p$ and $(x:y:z)$ and reflecting.
 So we have a map
  $$
  t_p(x:y:z)=
  (F_1(x,y,z),F_2(x,y,z),F_3(x,y,z)
  $$
  We have created a {\it Cremona
  transformation} (=biholomorphic birational
  map?).
 
  \begin{definition}[Cremona group]
	$$\operatorname{Bir}(\mathbb{P}^n)=\{\varphi
	:\mathbb{P}^n
	\overset{\operatorname{bir}}
	{\longrightarrow}
	\mathbb{P}^n:\text{bimeromorphic
	map} \}$$
\end{definition}

And we then have the surjective map
$$
\operatorname{Bir}(\mathbb{P}^2,C)
\overset{\pi}{\longrightarrow}
\operatorname{Aut}(C)
$$


	\item $(n,d)=(2,4)$. Here
	$S=X_4=\mathbb{P}^3$ is a smooth quartic
	surface.

		\paragraph{Problem} (Gizatullin) Which
		automorphisms of $S$ are induced (not
		necesarily by automorphisms of
		$\mathbb{P}^3$) but at least by Cremona
		transformations of $\mathbb{P}^3$? ie. are
		restrictions of
		$\varphi\in
		\operatorname{Bir}(\mathbb{P}^3,S)$

\begin{remark}
	Related to K3 surface structure,
	$$
	\operatorname{Bir}(\mathbb{P}^3,S)
	\overset{\pi}{\longrightarrow}
	\operatorname{Bir}(S)\cong
	\operatorname{Aut}(S)$$
\end{remark}

$S$ is a K3 surface.
$H^{2}(S,\mathbb{Z})\cong \mathbb{Z}^{22}$
and the Picard group acts on this lattice.

\begin{definition}[Picard rank of $S$]
	$$
	\rho(S)=
	\operatorname{rk}(\operatorname{Pic}(S))
	\in\{1,\ldots,20\}
	$$
\end{definition}

\begin{itemize}
\item If $S$ is very general, then
$\rho(S)=1$ and
$\operatorname{Aut}(S)=\{1\}$.

\item We are interested in $\rho(S)\geq 2$.
\end{itemize}
\end{enumerate}

\begin{example}[Oguiso, 2012]\leavevmode 
	\begin{enumerate}
		\item $\rho(S)=2$. Any cremona
		transformation that stabilizes the quadric
		is the identity:
			$$\operatorname{Aut}(S)=\mathbb{Z}\qquad
			\operatorname{Bir}(\mathbb{P}^3,S)=\{1\}$$
	
			
		\item $\rho(S)=3$,
			$$
			\operatorname{Aut}(S)=
			\mathbb{Z}_2*\mathbb{Z}_2*\mathbb{Z}_2
			$$
			All automorphisms are induced by Cremona
			transformations, ie. there is a surjective
			map
			$$
			\operatorname{Bir}(\mathbb{P}^3,S)
			\overset{\pi}{\longrightarrow}
			\operatorname{Aut}(S)
			$$

		\item (Paiva, Quedo 2023) Constructed
		surfaces with $\rho(S)=2$,
		$\operatorname{Aut}(S)=\mathbb{Z}_2$ and
		$\operatorname{Bir}(\mathbb{P}^3,S)=\{1\}$.
	\end{enumerate}
\end{example}

\begin{theorem}[A-Paiva-(Socrates)
Zika]\leavevmode
	Solution of Giztallin S problem for
	$\rho(S)=2$.
\end{theorem}

\begin{remark}
	Not exactly, but "there is a moduli space
	for K3 surfaces of dimension $20-\rho(S)$".
	The generic case is $\rho(S)=1$.
\end{remark}

\medskip\noindent
\textbf{Last application.}

\begin{definition}
	A {\it K3 surface} is a smooth
	projective surface that is simply connected
	and has a nowhere vanishing symplectic form
	$\omega\in H^{0}(S,\Omega^2_S)$
\end{definition}

\medskip\noindent
\textbf{Lattices of $S$.}

A {\it lattice} is a
finitely-generated abelian group with a
nondegenerate pairin.

\begin{enumerate}
	\item $H^{2}(X,\mathbb{Z})\cong \mathbb{Z}
	\overset{\pi}{\hookleftarrow}
	\operatorname{Pic}(S)
	$. And if we tensor this with $\mathbb{C}$
	we get $H^{2}(C,\mathbb{C})$, which admits a
	Hodge decomposition.

		Let's study automorphisms of a K3 surface.
		Let $g\in\operatorname{Aut}(S)$. It yields
		an element $g^*$ that acts on cohomology,
		ie $g^*\in\mathcal{O}(H^{2}(X,\mathbb{Z}))$
		preserving the Hodge decomposition. This is
		called {\it Hodge isometry}.

\begin{theorem}[Global Torelli
theorem]\leavevmode
\label{theorem-global-torelli-k3}
	\begin{itemize}
	\item The correspondence $g\mapsto g^*$ is
	injective
	
	\item If
	$\varphi\in\mathcal{O}(H^{2}(X,\mathbb{Z}))$
	is a Hodge isometry preserving the ample
	classes, then $\varphi=g^*$ for some
	$g\in\operatorname{Aut}(S)$.
	\end{itemize}
\end{theorem}
\end{enumerate}

\begin{example}
	$S\subset \mathbb{P}^3$ smooth quartic
	surface with $\rho(S)=2$. Using the
	equivalence of line bundles module
	isomorphism and curves modulo intersection,
	$$\operatorname{Pic}(S)= \left\langle{}H,C\right\rangle{} $$
	where $H$ is a hyperplane section of
	$\mathbb{P}^3$. Then
	$$Q=\begin{pmatrix} H^2& H\cdot C\\H\cdot C&
	C^2 \end{pmatrix} =\begin{pmatrix} 4&b\\b&2c
	\end{pmatrix} $$
	using that $2g(c)-1$, so the number in the
	lower-right on RHS is even.
\end{example}

\begin{remark}
	Hodge Index Them $Q$ is rank (1,1).
\end{remark}

\begin{definition}[Discriminant]
	$$r=\operatorname{disc}(S)=-\det Q$$
\end{definition}

\begin{proposition}
	$S$ general K3 surface (not needed that it
	is a quartic) with $\rho(S)=2$. Then
	$$\operatorname{Aut}(S)=\begin{cases}
		\{1\}\qquad &\text{(finite)}  \\
		\mathbb{Z}_2\qquad &\text{(finite,
		Dani-Ana)} \\
		\mathbb{Z}\qquad&\text{(infinite, Oguiso
		1)}  \\
		\mathbb{Z}_2*\mathbb{Z}_2\qquad
		&\text{(infinite, Dani-Ana)}
	\end{cases}$$
	where the first two are characterized by
	containing a en elliptic curve or a
	rational? curve, ir. $\exists
	D\in\operatorname{Pic}(S)$ such that
	$D^2=0,-2$. On the other hand, the last two
	cases are distinguised by $\not\exists
	D\in\operatorname{Pic}(S)$ such that
	$D^2=0,-2$.

	$$\exists
	\sigma\in\operatorname{Aut}(S)\text{ of
	order 2}\iff \exists  \text{ ample
	}A\in\operatorname{Pic}(S) \text{ such that
	}A^2=2.$$
\end{proposition}

\begin{remark}
	In low Picard numbers there are no
	symplectic involutions…?
\end{remark}

So to understand the surface we want to
understand those bundles and that is all in
the discriminant (not in the lattice
itself!).

\begin{remark}
	$\exists D\in\operatorname{Pic}(S)$ s.t.
	$D^2=k \iff x^2-ry^2=4k$  has integer
	solutions.
\end{remark}

Given the quadratic form we can find the
automorphisms.

\medskip\noindent
\textbf{Main theorem.}

\begin{theorem}[A-Paiva-Zika]\leavevmode
	$S\subset \mathbb{P}^3$ general smooth
	quartic surface with $\rho(S)=2$ and
	$\operatorname{disc}(S)=r$.
\begin{enumerate}
	\item (Negative answer to Gizatulla's
	problem) If $r>57$ or  $r=52$ then
		 $$
		 \operatorname{Bir}(\mathbb{P}^3,S)=\{1\}
		 $$
		 So we cannot realize any automorphism as a
		 Cremona transformation.

	\item If $r\leq 57$ and $r \neq 52$, then we
	get the full automorphism group, ie. a
	surjective map
		$$
		\operatorname{Bir}(\mathbb{P}^3,S)
		\overset{\pi}{\longrightarrow}
		\operatorname{Aut}(S)
		$$
\end{enumerate}
\end{theorem}

\medskip\noindent
\textbf{Birational geometry.}

Take the case of
$$
\xymatrix{
\operatorname{Bir}(\mathbb{P}^2)=
\left\langle{}\operatorname{Aut}(\mathbb{P}^2),
q\right\rangle{}\\
\operatorname{Aut}(\mathbb{P}^2)
\ar[u]_{\hookrightarrow}
}
$$
and the map
$$
\begin{aligned}
	q: \mathbb{P}^2
	&\overset{\operatorname{bir}}
	{\longrightarrow}
	\mathbb{P}^2
	\\
	(x:y:z) &\longmapsto
	\left(\frac{1}{x}:\frac{1}{y}:
	\frac{1}{z}\right)=(yz:xz:xy)
\end{aligned}
$$
which is well-known (Noether-Castelnuovo). So
that is a decomposition of automorphisms and
quadratics. Now in greater dimension, $n\geq
3$ we have

\begin{theorem}[Sakisov Program]\leavevmode
	(The theorem is much more general) Any
	$\varphi
	\in\operatorname{Pic}(\mathbb{P}^n)$ can be
	factorized with  {\it Sarkisov
	links}  $\varphi_i$:
	$$\xymatrix{
		\mathbb{P}^n
		\ar[r]^-{\varphi_1}
		\ar@/^2pc/[rrrr]^-{\varphi}
		&
		X_1\ar[r]^-{\varphi_2}\ar[d]
		&
		X_2\ar[r]\ar[d]
		&
		\cdots
		\ar[r]
		&
		X_k=\mathbb{P}^n\ar[d]
		\\
		&T_1&T_2&&
	}$$
\end{theorem}

Now look at $n=3$,
$\operatorname{Bir}(\mathbb{P}^3,S)\subset
\operatorname{Bir}(\mathbb{P}^3)$. This is a
Calabi-Yau pair:

\begin{definition}%
[Calabi-Yau pair]\leavevmode 
	A pair $(X,D)$
	\begin{itemize}
	\item Terminal projective variety.
	\item $K_X+D\sim0$ that is, a meromorphic
	top form that does not vanish on
	hypersurface and has simple pole on $D$.
	Then $(X,D)$ is called {\it log
	canonical}.
	\end{itemize}

Now take two Calabi-Yau pairs $(X,D_X)$ and
$(Y,D_Y)$.

 $$
\begin{aligned}
	\operatorname{div}(\omega_{D_X})&=-D_X\\
	\operatorname{div}(\omega_{D_Y})&=-D_Y
\end{aligned}
$$
We say that a birrational map $f:X\to Y$ is
{\it volume preserving} if
$f_*\omega_{D_X}=\omega_{D_Y}$.

\end{definition}


\begin{theorem}[Volume-Preserving]\leavevmode
	Everything like in the Sarkisov theorem but
	now maps are volume-preserving.
\end{theorem}

In our case, $(\mathbb{P}^3,S)$, we can
classify the v.p. Sakisov links from
$(\mathbb{P}^3,S)$. It starts by blowing up a
curve $C\subset S$. But this curve has genus
and degree very restricted, it's something
like
$$
(g(C),\operatorname{deg}(C))\in
\{(0,1),(0,2),\ldots,(11,10),(14,11)\}
$$

So for the main theorem, it was checked that
if $r>57$ there are no curves from the list.
And in the second item of the main theorem,
there exist these curves, for instance curve
$(14,11)$ for rank 56, then produce a link
that starts by blowing it up and magically
gives you the Cremona transformation that
restricts with automorphism with which you
started.

\begin{remark}
	So perhaps we expect the answer to G.
	problem to be almost never.
\end{remark}

\begin{remark}
	How does that blowing-up work?
\end{remark}

$$
\xymatrix{
X\ar[r]^-{\text{flops}}
\ar[d]_-{\operatorname{Bl}_C}
&
X\ar[d]^-{\text{contract?}}
\\
C\subset S\subset
\mathbb{P}^3\ar@{-->}[r]&
\mathbb{P}^3\supset
S'\supset C'
}
$$
So for example in case $(2,8)$ you obtain
something of the same type.

\section{Birational geometry of Calabi-Yau
pairs}
\label{section-calabi-yau-pairs}
%% speaker: Carolina Araujo
%% place: IMPA
%% date: January 31, 2025

\medskip\noindent
\textbf{Speaker.}
Carolina Araujo.

\medskip\noindent
\textbf{Place.}
IMPA.

\medskip\noindent
\textbf{Date.}
January 31, 2025.

\vspace{2em}
\medskip\noindent
\textbf{Pairs $(X,D)$.}
Itaca's program (1970's): $U$ complex
algebraic variety. We want to find invariants
like Kodaira dimension. Compactify $U
\rightsquigarrow X$, and consider
$X\setminus U:=D$ (the boundary). The
$\Omega^q_X(\operatorname{log}D)$

\begin{theorem}[Itaca, 1977]\leavevmode
Kodaira dimension of the pair $(X,D)$ does
not depend on the choice of compactification
(and it exists ;)
\end{theorem}

\begin{definition}[Lu-Zhang, 2017]\leavevmode
$(X,D)$ is
{\it Brody-hyperbolic}
{\it (Mori-hyperbolic)} if there
is no nonconstant holomorphic (analytic)
morphism  $f:\mathbb{C} \to X\setminus D$
and the same holds for the open strata of
$D$.
\end{definition}

\begin{remark}
$(X,D)$ is Brody-hyperbolic then $K_X +
D$ is ample.
\end{remark}

\begin{theorem}[Sraldi, 2019]\leavevmode
$(X,D)$ Mori-hyperbolic then $K_X+D$ is
nef.
\end{theorem}

\begin{remark}\leavevmode
How to make sense of ``the birational
geometry of $(X,D)$"
\end{remark}

\medskip\noindent
\textbf{Calabi-Yau pairs.}

\begin{definition}%
[Calabi-Yau pair]\leavevmode
We relax a condition: we won't ask that $X$
is smooth, but that it is a terminal
projective variety. And also that $K_X+D
\sim 0$ (I think this ``would be equivalent
to" Kodaira dimension 0). Also ask that
$(X,D)$ is log canonica.
\end{definition}

\begin{remark}\leavevmode
The condition $K_X+D \sim 0$ says that
there is unique up to scaling volume form
$\omega_D \in \Omega^m_{\mathbb{C}(X)}$
that does not vanish and has a simple pole
along $D$, i.e.
$\operatorname{div}(\omega_D)=-D$.
\end{remark}

Now we study the birational geometry of CY
pairs:

\begin{definition}\leavevmode
Let $(X,D_X),(Y,D_Y)$ CY pairs and $f: X
\dashrightarrow Y$ a birational map. We get
a pullback
$f_*:\Omega^n_{\mathbb{C}(X)}
\overset{\cong}{\to}
\Omega^n_{\mathbb{C}(Y)}$. Then $f:(X,D_X)
\dashrightarrow (Y,D_Y)$ is
{\it volume preserving} if
$f_*(\omega_{D_X})=\lambda \omega_{D_Y}$.
\end{definition}

\begin{remark}[Valuative
characterization]\leavevmode

\end{remark}

\begin{definition}\leavevmode
$(X,D)$ CY pair,
$$\operatorname{Bir}(X,D):=\{f \in
\operatorname{Bir}(X) : f:(X,D)
\dashrightarrow (X,D) \text{ is volume
preserving} \}$$
\end{definition}

\medskip\noindent
\textbf{The Sarkisov-Program.}

$\operatorname{Bir}(\mathbb{P}^n)$ Cremona
group.

\begin{theorem}[Noether-Castelnuovo,
1870-1901]\leavevmode
The Cremona group of $\mathbb{P}^2$ has a
nice set of generators: maps of degree 1 and
$q$:
$$\operatorname{Bir}(\mathbb{P}^2) =
\left\langle{}\operatorname{Aut}(\mathbb{P}^2),
q\right\rangle{}$$
where $q$ is the {\it standard
quadratic transformation} given by $(x:y:z)
\overset{q}{\mapsto }(yz:xz:xy)$
\end{theorem}

You might think that because it has such a
nice group generators it'd be easy to study
this group. It's not the case.

Now in dimension 3:

\begin{theorem}%
[Hilda Hudson (1927)]\leavevmode
$\operatorname{Bir}(\mathbb{P}^n)$ does not
admit a set of generators of bounded degree.
\end{theorem}

\begin{theorem}[Sarkisov Program, Corti 1995;
Hazon-McKunam 2013 for $n \geq
4$]\leavevmode
	$$\xymatrix{
		\mathbb{P}^n
		\ar[r]^-{\varphi_1}
		\ar[d]
		\ar@/^2pc/[rrrr]^-{\varphi}
		&
		X_1\ar[r]^-{\varphi_2}\ar[d]
		&
		X_2\ar[r]\ar[d]
		&
		\cdots
		\ar[r]
		&
		X_k=\mathbb{P}^n\ar[d]
		\\
	\{\operatorname{pt}\}
	&T_1&T_2&&\{\operatorname{pt}\}
	}$$
	The intermediate varieties $X_i/T_i$ are
	called {\it Mori-fiber spaces},
	and $\varphi_i$ are
	{\it Sarkisov links}. So the
theorem is that a map (what map?) can be
factorized by Sarkisov links.
\end{theorem}

\begin{theorem}%
[Corti-Kalog(?) 2016]\leavevmode
Any volume-preserving birrational map between
Mori-fiber CY pairs is a composition of
volume-preserving Sarkisov links. Now every
intermidiate variety admits a divisor making
	it a CY pair:
	$$\xymatrix{
		(X,D_X)\ar[r]^-{\varphi_1}\ar[d]
		\ar@/^2pc/[rrrr]^-{\varphi}
		&
		(X_1,D_1)\ar[r]^-{\varphi_2}\ar[d]
		&
		(X_2,D_2)\ar[r]\ar[d]
		&
		\cdots
		\ar[r]
		&
		(Y,D_Y)\ar[d]
		\\
	T_X	&T_1&T_2&&T_Y
	}$$
\end{theorem}

\begin{remark}\leavevmode
\begin{itemize}
\item The Sarkisov links
$$
\xymatrix{
 \mathbb{P}^n\ar@{-->}[r]\ar[d]
 &
 X_1\ar[d]
 \\
 \operatorname{pt}&T_1
}
$$
are not classified.

\item Depending on $D$, the
volume-preserving Sarkisov links
$$
\xymatrix{
 \mathbb{P}^n\ar@{-->}[r]\ar[d]
 &
 X_1\ar[d]
 \\
 \operatorname{pt}&T_1
}
$$
can be classified.
\end{itemize}
\end{remark}

\begin{theorem}[A-Coti-Massareti]\leavevmode
$D \subset \mathbb{P}^n$ a \textit{general}
hypersurface of degree $n+1$. If you want
to study the birational group of the pair
$(\mathbb{P}^n,D)$,
$\operatorname{Bir}(\mathbb{P}^n,D)$, then
this group is not interesting. Because
$\operatorname{Bir}(\mathbb{P}^n,D)=
\operatorname{Aut}(\mathbb{P}^n,D)$.

And then also consider $D$ smooth instead
of general, and with Picard rank
$\rho(D)=1$.
\end{theorem}

\begin{theorem}[A-C-M]\leavevmode
$D \subset \mathbb{P}^3$ general (A1
singularity and $\rho(S)=1$) quartic with a
singularity.
$\operatorname{Bir}(\mathbb{P}^3,D)$ is a
semidirect product of $\mathbb{G}$ and
$\mathbb{Z}/2\mathbb{Z}$.
Here $\mathbb{G}$ is ``form of
$\mathbb{G}m$ over $\mathbb{C}(x,y)$".
\end{theorem}


\section{Degenerations of the canonical
series for curves}
\label{section-esteves-canonical-series}
%% speaker: Eduardo Esteves, IMPA
%% place: Seminário das Sextas, PUC-Rio
%% date: September 20, 2024

\noindent
Eduardo Esteves, IMPA.
Semin\'ario das Sextas, PUC-Rio.
September 20, 2024.

\medskip\noindent
No abstract available.

\medskip\noindent{\bf Part 1/2.}

\medskip\noindent{\bf Introduction.}

What is the canonical series? Consider $C$ a
projective smooth connected curve over a
closed field $k$. The cotangent bundle
$\omega_C$ is the same as the bundle of
differentials, and the canonical bundle. It
has complex dimension 1. The local sections
are differentials.

A differential $\alpha \in\Gamma(C,\omega_C)$
is holomorphic, it can be meromorphic. It is
written as $\alpha=fdt$ for some $f\in k(C)$.
This leads to the notion of a
\textit{\textbf{divisor}} of a differential,
which in this case is

$$
\operatorname{div}(\alpha)=\sum
\operatorname{ord}_p(t_p)p
$$

and it is also the \textit{\textbf{canonical
divisor}}.

Now

$$
\omega_C=\mathcal{O}_C(k)
$$

The \textit{\textbf{genus}}, which is a
topological invariant, is also
$\dim_k\Gamma(C,\omega_C)$.

$$
\operatorname{deg}(\omega_C)=
\operatorname{deg}(K)=2g-2
$$

And

$$
\mathbb{H}=\Gamma(C,\omega_C)
$$

is the \textit{\textbf{canonical series}}.

\begin{example}
\label{example-esteves1-1}
Smooth quartic plane curve

$$
C=V(F)\subset \mathbb{P}^2
$$

We have that $\operatorname{deg}(F)=4$, and
$$
\begin{aligned}
\omega_C&=\mathcal{O}_C(1)=\mathcal{O}_{
\mathbb{P}^2}(1)|_{C}\\
g&=\dfrac{(d-1)(d-L)}{2}=3\\
\mathcal{O}_C(1)&=\mathcal{O}_C(L\cap C)\\
\operatorname{deg}(\mathcal{O}_C(1))&=4=
2g-2\\
\dim_k\Gamma(G\mathcal{O}_C(1))&=3=g
\end{aligned}
$$
\end{example}
\begin{quotation}
Think of the canonical series as a space of
linear sections of a line bundle, but also as
a collection of divisors parametrized by
$\mathbb{P}^2$
\end{quotation}

\begin{example}[A singular quadric]
\label{example-esteves1-2}
\end{example}

\medskip\noindent{\bf What we do.}

Given a nodal curve $X$ (at a node there are
two "branches" that intersect) which is
general for its topology ($G=(V,E)$ dual
graph) where
\begin{itemize}
\item $V$ is the set of irreducible
components). There is a correspondence of the
vertices in this graph and curves in $X$ :
$v\in V\iff X_v \subset X$.
\item $E$ is the set of nodes. Here $e \in
E\iff N_e\in X$. So an edge is a pair of
points if the node belongs to the
intersection of the corresponding curves:

$$
e=\{u,v\} \iff N_e\in X_u\cap X_v
$$

\item The \textit{\textbf{genus function}}
associates to every component its geometric
genus:
$$
\begin{aligned}
g: V &\longrightarrow \mathbb{Z}_{\geq 0} \\
g(v) &=\text{geometric genus of }X_v
\end{aligned}
$$
(I think the geometric genus is the genus of
the normalization of the variety.)
\end{itemize}
This is the combinatorial data attached to
the curve.

We
We look for a general curve with respect to
the geometric genus, and say it is
\textit{\textbf{general for its position}} if
the nodal points are in general position.

Then we have a stratification of

$$
\overline{\mathcal{M}}=\{\text{stable curves
} X \text{ with finite automorphism group}
\}
$$

And here \textit{stable} means nodal. So for
example you can have stability if the degree
of $\omega_C|_{X}$ is positive.

So the stratification is:

$$
\overline{\mathcal{M}}_{g}=\bigsqcup
\mathcal{M}_{(G,g)}
$$

where

$$
\mathcal{M}_{G,g}=\{X\text{ s.t. $G$ is the
dual graph of $X_\omega$ genus function $g$.}
\}
$$

And we have that

$$
\operatorname{codim}M_{G,g}=|E|,
$$

the number of edges.

\begin{remark}
\label{remark-esteves1-1}
These are graph curves. All their components
are $\mathbb{P}^1$ and they intersect in
prescribed way.
\end{remark}

\medskip\noindent{\bf Some combinatorics.}

We have the \textit{\textbf{genus formula}}:

$$
P_a(X)=\sum g_v+g(G)
$$

$$
g(G)=|E|-|V| +1
$$

$$
g=3g-3-(2g-2)+1
$$

$$
\operatorname{deg}(F)=4
$$

\medskip\noindent{\bf Remark.}
When you have maximum number of edges, you
force everything to be of a particular kind
by the genus formula. (This follows from
stability condition.)

\begin{remark}
\label{remark-esteves1-2}
We may check stability looking if the
canonical bundle is trivial.
\end{remark}

\medskip\noindent{\bf Remark.}
Stable $\iff$ every component which is
$\mathbb{P}^1$ has at least 3 special points.

\medskip\noindent{\bf What we do.}

Now let's finish the statement we started
before:

Given a nodal curve $X$ which is general for
its topology, we describe all limits of the
canonical series in any degeneration to $X$,
and construct a parameter space for them.

\begin{exercise}
\label{exercise-esteves1-1}
Let $X$ be a smooth projective variety such
that $K_X$ is ample. Prove that its
automorphism group is finite.
\end{exercise}

We would like to study the moduli of stable
curves. So we have a parameter for the
objects we want to classify. Diaz-Cutievman
described the locus of curves with special
Weierstrass points.

Weierstrass points are such that the line
(what line?) intersect the curve in at least
(some bound). So for a quartic,

$$
P\text{ is a W point }\iff I(P;T_pC\cap
C)\geq 3
$$

and in fact

$$
g^3-g=24\text{ W points.}
$$

\begin{remark}
\label{remark-esteves1-3}
A general smooth curve of genus 3 has exactly
24 Weierstrass points.
\end{remark}

\begin{exercise}
\label{exercise-esteves1-2}
The genus $g$ curve has $g^3-g$ Weierstrass
points.
\begin{enumerate}
\item[a.] For $g=3$ and plane quartics.
\item[b.] For hyper-elliptic curve of genus
$3$ (also define W point in this case).
\item[c.] For non-hyperelliptic genus 4
curve.
\item[d.] Etc.
\end{enumerate}
\end{exercise}

\medskip\noindent{\bf The additional data
(?).}

Take your variety. The drawing is a bunch of
blue lines. Take another line (red). How do
the y intersect?

$$
\lim_{t\to 0}X_t\cap L=X\cap L.
$$

And intersection with another curve $F$?

$$
\lim_{t \to 0} X_t\cap L_0=F\cap L_0
$$

Let $L=L_0$. We have
$$
\begin{aligned}
L_0L_1L_2L_3+tF&=0\\
L_0&=0
\end{aligned}
$$
Dividing by $t$,
$$
\begin{aligned}
LL_1L_2L_3+F&=0\\
L_0&=0
\end{aligned}
$$
Now look at linear series generated by
$LL_1L_2L_3\forall L$ and $F$ on $L_0=0$.
$L_0=X$.
$$
\begin{aligned}
(\alpha Y+\beta Z)L_1L_2L_3+\gamma
F&=0,\qquad (\alpha,\beta,\gamma
)\in\mathbb{P}^2\\
L_0&=0
\end{aligned}
$$

\medskip\noindent{\bf After break.}

Now we explain how these systems of divisor
appear and how we are going to handle them.

The limit of the $\overset{\vee
}{\mathbb{P}}^2$ is some divisors.

We are considering a smoothing

$$
\xymatrix{\mathfrak{X}\ar[d]\\
B=\Delta_0=\operatorname{Spec}k[[t]]\subset
\mathbb{C}}
$$

And we have
$$
\begin{aligned}
\omega_{\mathfrak{X} /B}&\text{ is the
relativa canonical bundle}\\
\omega_{\mathfrak{X}
/B}\Big|_{\mathfrak{X}_\eta}=\cap_{
\mathfrak{X}_\eta}\\
\omega_{\mathfrak{X}
/B}\Big|_{\mathfrak{X}_\sigma}&=\omega_{X}
\subseteq
\Omega_X=\bigoplus_{v\in V}\Omega_v
\end{aligned}
$$
where $\Omega_v$ is the space of meromorphic
differentials over $C_v=\tilde{X}_v$.

So it's a family of bundles that is the
canonical bundle on the general fibers and on
the exeptional fiber it is the canonical
bundle too.

\medskip\noindent{\bf Regular differentials
(Rosenlicht).}

$$
(\eta_v)_{v}\in\bigoplus_{v}\Omega_v
$$

$$
\operatorname{res}_{p_a}\eta_v+
\operatorname{res}_{p_{\bar{a}}}\eta_\omega
=0
$$

Now $\eta_v$ can only have poles at the
branches, and they should be simple.

\begin{remark}
\label{remark-esteves1-4}
Looks like we have been computing how the
bundles $\mathcal{O}_L(1)$,
$\mathcal{O}_{L_i}$ look like when restricted
to different subvarieties. So for example

$$
\mathcal{O}_{\mathfrak{X}}(1)(-L_0)|_{L_i}=
\mathcal{O}_{L_i}(1)(-1)
$$

is just a skyscraper sheaf.
\end{remark}

In the end we concluded that $(L_h,W_h)$ has
infinitely many linear series in $X$ where

$$
W_h=
\Gamma(\mathfrak{X},\mathcal{L}_{h})\Big|_{
\mathfrak{X}_0}\cap
\Gamma(X,L_h)
$$

So, importantly,

$$
\{0\neq  s\in W_h\qquad Z(s)\subseteq X\qquad
|Z(s)| <\infty\} =\qquad \text{lim divisors}
$$

$h$ vary.

$$
\begin{aligned}
\omega_X&\subseteq \Omega
=\bigoplus\Omega_v\\
\omega_{\mathfrak{X}}\left( -\sum h(v)X_{v}
\right) \Big|_{X}&\subseteq \Omega \\
\Omega_{X_{ \omega}}\left( \sum p_u-\sum p_u
\right)
\end{aligned}
$$

\medskip\noindent{\bf Canonical case.}
$$
W_h\subseteq \Omega =\bigoplus\Omega_v
$$

where $\Omega_v$ is the space of meromorphic
differentials on $C_v=\tilde{X}$. $\dim
W_h=g$.

\medskip\noindent{\bf Kapranov.}

Take some $k$-vector spaces $U_v$ and
consider

$$
U=\bigoplus_{v\in V} U_v
$$

and take the grassmanian of subspaces of
dimension $g$, and the torus action:

$$
\operatorname{Gr}(g,U)\curvearrowleft
\mathbb{G}^\vee_m=\{\psi:V\to k^*\}
$$

and the projection maps:

$$
\theta_I:\bigoplus_{v\in V}
U_v\longrightarrow \bigoplus_{v\in I}
U_v,\qquad I\subseteq V
$$

$W$ general, $\theta_I\Big|_{W}$ has maximal
rank.

So consider the orbit, it is a Chow variety:

$$
\overline{[\mathbb{G}^\vee_m\cdot
W}\in\operatorname{Chow}(
\operatorname{Gauss}(g,U))
$$

And we have the Chow quotient/Hilbert
quotient/Mumford quotient (studied first by
Thaddeus):

$$
\overline{\{\overline{[\mathbb{G}^\vee_m\cdot
W]},W\text{ general} \}} \subseteq
\operatorname{Chow}(\operatorname{Gauss}(g,
V))
$$

And then

$$
\partial\overset{\{*\}} =\sum
[\mathbb{G}^\vee_m\cdot W_i],\qquad
\operatorname{Gauss}(g,V)
$$

The polytopes asrouled to $W_i$ form a
polyhedral decomposition of a certain
polytope. Now
$$
\begin{aligned}
W\subseteq U\implies
\mu_W:2^V\longrightarrow \mathbb{Z}
\end{aligned}
$$
which is a submodular function,

$$
\mu=\mu_W(I)=\dim _k\theta_I(W)
$$

$$
\mu(I)+\mu(S)\geq \mu(I\cap J)+\mu(J\cup J)
$$

$$
P_\mu=\{q\in\mathbb{R}^{\vee }:q(I)\leq
\mu(I)\;\forall I,\; q(V)=\mu(V)\}
$$

\medskip\noindent{\bf What Kapranov
observed.}
$$
K:\bigcup P_{W_i}=P_{W\text{ general} }
$$

\bigskip
\medskip\noindent{\bf Part 2/2.}

\medskip\noindent{\bf Reminder on first
part.}

We have been parametrizing canonical series.

We have $X$ a projective connected nodal
curve over a closed field of characterstic
zero. We created its dual graph, where
\begin{itemize}
\item Vertices correspond to $C_{v}$,
normalization of components of $X$ associated
to $v$.
\item Edges correspond to $p^e$ node of $C$.
\item $e$ connects $u$ and $v$ iff  $p^e\in
X_u\cap X_v$.
\end{itemize}

Associated to the dual graph there's also the
arrow set $\mathbb{E}$. There's a 2-1 map
from the set of arrows to the set of edges.
To each arrow we associate a branch $p^a \in
C_v$, where $a=uv\in\mathbb{E}$. We denote
$$
\begin{aligned}
\mathbb{E}_v=\{a\in \mathbb{E}:t_a&=\{\text{
tail of $a$}\}\text{-}u  \}
\end{aligned}
$$
so that $\mathbb{E}=\bigsqcup_{v \in
V}\mathbb{E}_v$.

We also have the \textit{\textbf{genus
function}} $g$ mapping $V$ to its genus, and
the \textit{\textbf{genus formula}}

$$
p_a(X)=g(V)+g(G)=\sum_{v\in V}g(v)+|E| +|V|
$$

\medskip\noindent{\bf Today.}

We want to consider \textit{smoothings}
$\pi:\mathfrak{X} \longrightarrow B$ of $X$:

$$
\xymatrix{\mathfrak{X}\ar[d]^{\pi}\\
B={\begin{cases}A_0\\
\operatorname{Spec}(k[\![E]\!])\end{cases}}}
$$

So
$\mathfrak{X}_\mathcal{O}\xrightarrow{\cong
}X$ $\mathfrak{X}_\eta$ smooth (over Laurent
series $k(\!(t)\!)$. $\mathfrak{X}$ is
regular away from $p^e\in X$.

$$
\hat{\mathcal{O}}_{\mathfrak{X},p^e}=
\dfrac{k[\!\![
t,u,v]\!\!]}{uv-t^{\ell_e}}\ell_e\in
\mathbb{Z}_{>0}
$$

where $\ell:E\longrightarrow Z_{>0}$ is the
\textit{\textbf{edge lenght function}}. We
also have $\Gamma=(G,\ell)$ the
\textit{\textbf{metric graph}}, and
\begin{itemize}
\item $\omega_{\mathfrak{X} /B}$ the
\textit{\textbf{relative canonical bundle
}}on  $\mathfrak{X}$.
\item $\omega_{\mathfrak{X}
|B}|_{\mathfrak{X}_\eta}=\Omega_{
\mathfrak{X}_\eta}$
the \textit{\textbf{bundle of
differentials}}.
\item $\omega_{\mathfrak{X}
|B}|_{X}=\omega_X$ the canonical bundle of
$X$. $\omega_X$ is the \textit{\textbf{bundle
of regular differentials}} (Rosenlicht 50's)
\end{itemize}

\medskip\noindent{\bf Notation.}
$\Omega_v$ is the space of meromorphic
differentials of $C_v$,

$$
\Omega=\bigoplus_{v\in V} \Omega_v
$$

$\alpha=(\alpha_{v})_{v}\in\Omega$ is regular
if
\begin{enumerate}
\item $\forall a\in\mathbb{E}$,
$\alpha_{t_a}$ has at most a simple pole at
$p^a$.
\item $\forall a\in\mathbb{E}$,
$\operatorname{res}_{p_a}(\alpha_{t_a}+
\operatorname{res}_{pa}(\alpha_{t_a})=0$.
\end{enumerate}

\medskip\noindent{\bf Abelian differentials.}

$C_v$,  $D\in D_N(C_v)$, $D=P-N$,  $P,N\geq
0$, $\operatorname{supp}(P)\cap
\operatorname{supp}(N)=\varnothing$.

$$
\mathbb{H}_v(D)=\{\alpha\in\Omega_v|
\operatorname{div}_\infty(\alpha)\leq
P,\operatorname{div}_0(\alpha)\geq N\}
$$

$$
\mathbb{H}_v\left(0 \right) =\text{Abelian
differentials }
$$

$$
W_0=\Gamma(X,\omega_X)\subseteq \bigoplus_{v}
\mathbb{H}_v\left( \sum_{a\in\mathbb{E}}p_a
\right)
$$

where $W_0$ is the space of global regular
differnetials. $\dim
W_0=g_{\mathfrak{X}_\eta}=p_a(X)$, the
arithmetic genus of $X$.

If $D$ is general and characteristic of  $k$
is zero, then
$$
\begin{aligned}
\dim_k
\mathbb{H}_v(D)&=\max(g_v+
\operatorname{deg}D-1,0)\qquad
\text{if }p\gneq 0\\
&=max(g_v+\operatorname{deg}D,0)\qquad
\text{if} p=0
\end{aligned}
$$

\medskip\noindent{\bf Back to our objective.}

$$
\mathcal{L}_{h}=\omega_{\mathfrak{X}
/B}\otimes"\mathcal{O}_{\mathfrak{X} }\left(
-\sum h_vX_v \right)
$$

where $h:V\longrightarrow \mathbb{Z}$.

$$
\mathcal{L}_{h}|_{X_\eta}=\Omega^1_{
\mathfrak{X}_\eta}
$$

$$
\mathcal{L}_{h}|_{C_v}=\omega_{C_v}\left(
\sum_{a\in\mathbb{E}_{v}}
(1+\partial_{\ell}h(a)p^a \right)
$$

[$*$some missing formulas $*$]

$$
W_h=\operatorname{Im}(\Gamma(\mathfrak{X},
\mathcal{L}_{h})\longrightarrow
\Gamma(X,\mathcal{L}_{h}|_{X})
$$

$$
\dim W_h=p_a(X)
$$

$$
W_h\subseteq \bigoplus_{v\in
V}\mathbb{H}_v\left(
\sum_{a\in\mathbb{E}_v}(1+\partial_\ell
h(a)p^a \right) \subseteq \bigoplus_{u\in
V}\Omega_u    =\Omega
$$

\medskip\noindent{\bf The goal.}
is to describe and parametrize the collection
of subspaces

$$
\mathcal{C}=\{W_h\in\Omega|h\in\mathbb{Z}^
\vee\}
$$

when the points $p^a$ on  $C_v$ for  $a \in
\mathbb{E}$ are in general position.

\medskip\noindent{\bf A space that
parametrizes $W_h$.}

This is work by Kapranov. First of all, $W_h$
is not uniquely defined:

$$
V\twoheadrightarrow V_h=V/\sim_h
$$

where $u\sim_h v\iff h(u)=h(v)$.

We have a torus action

$$
\mathbb{G}^{V_h}_m=\{\psi:V_h\longrightarrow
K^*\} \subseteq \mathbb{G}^V_m
$$

$$
W_h\subseteq \Omega,\qquad \psi\cdot
W_h=\{(\psi_v\alpha_v)_v|(\alpha_v)_v\in
W_h\}
$$

$\mathbb{G}^{V_h}_m\cdot W_h$ orbit is
well-defined.

The case of $h=0\implies V_h=\{V\}$, so $W_0$
is well-defined.

So actually we want to parametrize not the
$W_h$ but their orbits, $\mathbb{G}^V_m\cdot
W_h\in
\operatorname{Gr}(g,\Omega)/\mathbb{G}^V_m$.

$$
W\subseteq U=\bigoplus_{v} U_v\subseteq
\bigoplus_{v} \Omega_v=\Omega
$$

$$
\mathbb{G}^\vee_m\curvearrowright
\operatorname{Gr}(g,U)\subset
\operatorname{Gr}(g,\Omega)\curvearrowleft
\mathbb{G}^\vee_m
$$

\medskip\noindent{\bf Polyhedral approach.}

Basically what you realize is that the only
orbits that matter are the maximal dimension
ones, and these correspond to maximal
dimesional polytopes. (These appeared in the
work of Kapranov and was later generalized by
others.

The idea is as follows. Fix a decomposition
of your space $W\subseteq
\Omega=\bigoplus_{v\in N} \Omega_v $. A
\textit{\textbf{submodular function}}
$$
\begin{aligned}
\nu_W: 2^V &\longrightarrow \mathbb{Z} \\
V\supseteq I &\longmapsto \dim_k\theta_I(W)
\end{aligned}
$$
$*$ some formulas $*$ Definition of base
polytope. The orbit is maximal dimensional if
and only if the polytope is. Definition of
bricks.

\begin{theorem}
\label{theorem-esteves2-1}
The polytope associated to the $W_h$ is a
union of bricks of maximal dimension.
\end{theorem}

\medskip\noindent{\bf Editorial note.} The
two parts were typed as two separate files;
here they are joined in one section. Both
files have an empty abstract field. Some
diagrams in the originals were commutative
diagrams drawn with tikz-cd; they are
rewritten as xymatrix diagrams with the same
arrows and labels.

\section{On the algebraic hyperbolicity of
projective hypersurfaces}
\label{section-mioranci-alg-hyperbolicity}
%% speaker: Lucas Mioranci, IMPA
%% place: Geometric Structures on Manifolds
%% place: Seminar, IMPA
%% date: November 14, 2024

\noindent
Lucas Mioranci, IMPA.
Geometric Structures on Manifolds Seminar,
IMPA.
November 14, 2024.

\medskip\noindent
{\bf Abstract.} A complex projective manifold
$X$ is said to be algebraically hyperbolic if
every integral curve $C$ of $X$ satisfies the
inequality $2g(C)-2>\varepsilon \cdot deg(C)$
for a fixed positive $\varepsilon$ and ample
divisor on $X$. This talk aims to review some
techniques used to prove the algebraic
hyperbolicity of very general hypersurfaces
of degree $d>2n-2$ in $\mathbb{P}^n$.

\medskip\noindent
{\bf Upshot.} Recall that last week we talked
about Kobashi hyperbolicity, which means that
Kobayashi pseudo-distance is non-degenerate.
We saw that this implies Brody hyperbolicity,
which means that $X$ has no non-constant
holomorphic maps. If $X$ is compact this is
an equivalence.

Today we shall see that this implies
algebraic hyperbolicity. Demially's
conjecture is that this last implication is
actually iff.

\begin{definition}
\label{definition-mioranci-1}
A complex projective variety $X$ is
\textit{\textbf{algebraically hyperbolic}} if
there exists $\varepsilon>0$ and ample
divisor $H$ such that, for every integral
curve $C \subset X$ of geometric genus
$g(C)$,

$$
2g(C)-2 \geq \varepsilon
\operatorname{deg}_H(C)
$$

\end{definition}

\begin{remark}[Vitorio]
\label{remark-mioranci-1}
Definition of \textit{\textbf{degree}} is
$H.C$.
\end{remark}

\begin{remark}[Misha]
\label{remark-mioranci-2}
Is there not an associated metric here? Take
the metric in all agebraic curves inside $X$.
Perhaps it is equivalent?
\end{remark}

\begin{remark}
\label{remark-mioranci-3}
In particular, $X$ does not contain any
rational or elliptic curves.
\end{remark}

In this talk: $X \subset \mathbb{P}^n$ very
general projective hypersurface of degree
$d$.

\medskip\noindent{\bf Conjecture
(Kobayashi).} $X$ is algebraically hyperbolic
for $d$ sufficiently large ($d\geq 2n-2$,
$n\geq 4$ or $d\geq 2$, $n=3$).

\medskip\noindent{\bf $n\geq 3$.}

\begin{theorem}[Clemens '86, E. '88]
\label{theorem-mioranci-1}
$n \geq  4$, $d\geq 2n$.
\end{theorem}

\begin{theorem}[Voisin, '96, 98]
\label{theorem-mioranci-2}
$n \geq 4$, $d \geq  2n-1$.
\end{theorem}

\begin{theorem}[Pocienza '04, Clemens-Ron
'04]
\label{theorem-mioranci-3}
$n\geq 6$, $d \geq  2n-2$.
\end{theorem}

\begin{theorem}[Yeong '22]
\label{theorem-mioranci-4}
$n\geq 5$, $d\geq 2n-2$.
\end{theorem}

Open: $(n,d)=(4,6)$.

\medskip\noindent{\bf $n=3$.}

\begin{theorem}[Xu, '94]
\label{theorem-mioranci-5}
$n=3$, $d \geq  6$.
\end{theorem}

\begin{theorem}[Coskan-Ried, '99]
\label{theorem-mioranci-6}
$n=3$, $d \geq 5$.
\end{theorem}

And $n=3$, $d=4$ is a  K3 so (or is it
\textit{because}?)  it has rational curves.

\medskip\noindent{\bf Proof for $n=4$, $d\geq
2n-2$.}

\begin{proof}
Main reference  \textit{Algebraic
hyperbolicity of the very general quintic
surface in $\mathbb{P}^3$}, Coskun and Reid
'99.

\textit{Algebraic hyperbolicity of very
general surface}, Coker and Reid '22.

\textit{Algebraic hyperbolic of very general
hypersurface in
$\mathbb{P}^n\times\mathbb{P}^n$}, Young '22.
\end{proof}

\textbf{Open}: $\mathbb{P}^2 \times
\mathbb{P}^1$.

\begin{proof}
Let
$S_d=\mathbb{P}H^{0}(\mathbb{P}^n,
\mathcal{O}_{\mathbb{P}^n}(d))$
and let

$$
\mathcal{X}=\{(p,[X]):p\in X\}\subseteq
\mathbb{P}^n\times s_d
$$

which is a universal degree $d$ hypersurface
in $\mathbb{P}^n$.

We can find a generically injective map
$h:Y\to X$ over $U
\overset{\operatorname{op}}{\subseteq} S_d$
of curves of geometric genus $g$ and degree
$e$. {$Y$ is a family of curves $Y_t$ inside
each hypersurface $X_t$ for every parameter
$t\in U$}.

Define the \textit{\textbf{vertical tangent
bundles}} are the kernels

$$
0\to T_{X/\mathbb{P}^n}\to T_{\mathcal{X}}\to
\pi^* _2T_{\mathbb{P}^n}\to 0
$$

$$
0\to T_{Y/\mathbb{P}^n}\to T_Y\to  h^*\pi^*
_2T_{\mathbb{P}^n}\to 0
$$

\begin{itemize}
\item $Y$ dominates $U$ under $\pi_1\circ h$.
\item We can assume $Y$ is stable under the
$\mathsf{GL}(n+1)$-action.
\end{itemize}
So $\pi_2\circ h$ dominates $\mathbb{P}^n$
and $T_Y \longrightarrow h^*\pi^*
_2T_{\mathbb{P}^2} $ is surjective.

Define the \textit{\textbf{normal bundles}}
as cokernels of

$$
0\to T_Y\to h ^* T_{\mathcal{X}}\to
N_{h/\mathcal{X}}\to 0
$$

$$
0\to T_{Y_t}\to h_t^* T_{\mathcal{X}_t}\to
N_{ht/\mathcal{X}_t}\to 0
$$

Denote

$$
i_t:\mathcal{X}_t\to\mathcal{X}
$$

$$
\begin{aligned}
j_t:Y_t &\longrightarrow Y \\
p &\longmapsto (p,t)
\end{aligned}
$$

\medskip\noindent{\bf (technical) Lemma 1.}
$N_{h_t/\mathcal{X}_t}\equiv j^* _t
N_{h/\mathcal{X}}$

\begin{proof}
Big commutative diagram.
\end{proof}

\medskip\noindent{\bf Lemma 2.}
$N_{h_t/\mathcal{X}_t}\cong j^*_t K$ where
$K=\operatorname{coker}(T_{Y/\mathbb{P}^n}\to
T_{Y/\mathbb{P}^n}$.

\begin{definition}
\label{definition-mioranci-2}
The \textit{\textbf{Lazarsfeld-Mukai bundle}}
$M_d$ is the kernel of

$$
0\to M_d\to \mathcal{O}_{\mathbb{P}^n}\otimes
H^{0}(\mathbb{P}^n ,
\mathcal{O}_{\mathbb{P}^n}(d))\xrightarrow{
\operatorname{ev}}\mathcal{O}_{
\mathbb{P}^n}(d)\to
0
$$

\end{definition}

\begin{remark}
\label{remark-mioranci-4}
The fiber over $p \in \mathbb{P}^n$ is $M_d$
is the space of sections vanishing at $p$.
\end{remark}

\medskip\noindent{\bf Lemma 3.}
$T_{\mathcal{X}/\mathbb{P}^n}\cong \pi
^*_2M_d$

By Lemma 2 and 3:

$$
j^*_sh^*
T_{\mathcal{X}/\mathbb{P}^n}\longrightarrow
j^* _t K\cong N_{ht/\mathfrak{X}_t}
$$

So we have a surjection

$$
M_d|_{Y_t}\twoheadrightarrow N_{ht/X_t}.
$$

\begin{definition}
\label{definition-mioranci-3}
Let $\mathcal{E}$ be a vector bundle on
$\mathbb{P}^n$, let $L$ be a globally gen
line bundle. We say $L$ is
\textit{\textbf{section-dominating}} for
$\mathcal{E}$ if $\mathcal{E} \otimes
\mathcal{L}^\vee$ is globally gen. and

$$
H^{0}(L \otimes I_p)\otimes H^{0}(\mathcal{E}
\otimes \mathcal{L}^\vee)\longrightarrow
H^{0}(\mathcal{E} \otimes I_p)
$$

is surjective for all $p \in \mathbb{P}^n$.
\end{definition}

\begin{example}
\label{example-mioranci-1}
in $\mathbb{P}^n$,
$\mathcal{E}=\mathcal{O}_{\mathbb{P}^n}(d)$,
$L=\mathcal{O}_{\mathbb{P}^n}(1)$.
\end{example}

\begin{proposition}
\label{proposition-mioranci-1}
There is a sujection
\end{proposition}

\end{proof}

\medskip\noindent{\bf Editorial note.} The
notes stop at this point: the statement of
the last proposition is unfinished in the
original. The short exact sequences, drawn
with tikz-cd in the original, are written as
plain displays with the same terms and maps.

\section{Introduction to K-stability of Fano
varieties}
\label{section-montero-k-stability-fano}
%% speaker: Pedro Montero, UTFSM-Chile
%% place: minicourse, IMPA Verão
%% date: February 17-28, 2025

\noindent
Pedro Montero, UTFSM-Chile.
IMPA Verão (minicourse).
February 17-28, 2025.

\medskip\noindent
No abstract available.

\medskip\noindent{\bf Lecture 3, Oops!}

\medskip\noindent{\bf Intro for today.}

{\bf Aim.} To study the K-stabilty using only
the birational geometry of $(X,L)$ instead of
all $\mathfrak{X},\mathcal{L})$. We will
study a corresponding

$$
\left\{ \text{T. C. } \mathfrak{g}, (X,L)
\right\} \leftrightsquigarrow \left\{
\text{Filtrations } \mathfrak{g},
R(X,L)=\oplus_{m \in
\mathbb{N}}H^{0}(X,L^{\otimes m}) \right\}
$$

$$
\leftrightsquigarrow \left\{
\text{(divisorial) valuations on
$\mathbb{C}(x)$}  \right\}
$$

\begin{remark}
\label{remark-montero-1}
Most of what follows is based on
\textit{Uniform K-stability, DH measures and
singularity of pairs} by Buckson Hisamoto,
Jonsson (2017). Barely 100 pages. Super nice.
\end{remark}

{\bf Recall.} (this is what we used to define
the \textit{\textbf{weight of an action}}.)
$V$ $k$-vector space of finite dimension.
Then $\mathbb{G}_n \curvearrowright V$
induces a weight decomposition

$$
V = \bigoplus_{x \in \mathbb{Z}}V_\lambda
$$

where $V_\lambda=\{v \in V:tv = t^{-\lambda
v}\forall  t \in \mathbb{G}_m\}$

Conversely, given such a decomposition, we
define

$$
t\cdot v:= \sum_{\lambda \in
\mathbb{Z}}t^\lambda v_\lambda,\qquad
\text{for } v = \sum_{\lambda}v_\lambda.
$$

\textbf{New stuff.}

\begin{definition}
\label{definition-montero-1}
Let $V$ be a finite-dimensional vector space
(e.g. $V= H^{0}(X,L)$). A
\textit{\textbf{$\mathbb{Z}$-filtration}} of
$V$ is

$$
\{F^\lambda V\}_{\lambda \in
\mathbb{Z}}\subset V\qquad \text{sub v.s.}
$$

such that
\begin{enumerate}
\item \textbf{(Decreasing.)} $F^{\lambda +1}V
\subseteq F^{\lambda}V$.
\item \textbf{(Stabilizes.)} $F^\lambda V=0
\forall \lambda \gg 0$ and $F^{\lambda}V=V
\forall  \lambda \ll 0$.
\end{enumerate}
The important thing is that we have an
associated algebra: the
\textit{\textbf{associated Rees algebra}} is
the finitely generated and torsion-free
$k[t]$-module

$$
\operatorname{Rees}(F^*
V):=\bigoplus_{\lambda \in
\mathbb{Z}}(F^\lambda V) E^\lambda \subseteq
V[t,t^{-1}] \overset{\operatorname{def}}{=} V
\otimes_k k[t,t^{-1}]
$$

where $t \cdot (vt^{-\lambda})=v t ^{-
\lambda +1}= vt^{-(\lambda-1)}$.
\end{definition}

\begin{remark}
\label{remark-montero-2}
In practice you \textit{are} looking at the
vector space of sections.
\end{remark}

\medskip\noindent{\bf Rees correspondence.}

It is a correspondence of vector bundles on
$\mathbb{A}^1$ that are compatible with the
$\mathbb{C}^*$ action, i.e. toric vector
bundles in $\mathbb{A}^1$. Right so this is a
particular case of the ``Klyasko's
classification" of toric vector bundles. So
remember that the point of toric geometry is
that you have everything encoded in
combinatorics, these vector bundles are
studied combinatorically.

$$
\left\{ \substack{\text{$Z$-filtrations of}
\\ \text{fin. dim. $k$-v.s. $V$} } \right\}
\overset{1-1}{\leftrightsquigarrow}\left\{
\substack{\mathbb{G}_m \text{-linearized}  \\
\text{v.b. $V \to \mathbb{A}^1$} } \right\}
$$

$R=\operatorname{Rees}F^* \mapsto  V=
\mathbb{V}(\tilde{R}) \to
\mathbb{A}^1=\operatorname{Spec}k[t]$

$\mathbb{G}_m$-linear since
$\mathbb{Z}$-grading is compatible with the
one in $k[t]$. Conversely, given  $V \to
\mathbb{A}^1$, a $\mathbb{G}_m$-linearized
vector bundle.

So the point of ``linearizing a vector
bundle", kind of by definition, is that you
have an induced action, i.e. we have
$\mathbb{G}_m \curvearrowright
H^{0}(\mathbb{A}^1,V)=\bigoplus_{\lambda \in
\mathbb{Z}}H^{0}(\mathbb{A}^1,V)_\lambda$
given by $t \cdot
\sigma(x)=\sigma(t^{-1}\cdot x)^{\lambda \in
\mathbb{Z}}$.

\medskip\noindent{\bf My own understanding.}
You start with a projective scheme and a line
bundle on it $X,L$. You produce a thing
called \textit{\textbf{test configuration}}
$(\mathfrak{X},\mathcal{L})$ which admits a
$\mathbb{G}_m$-action, and that's awesome.
Then you construct a  $\mathbb{Z}$-filtration
on each $R_m=H^{0}(X,L^{\otimes rm})$. This
filtration behaves well under the projection:

$$
H^{0}(\mathfrak{X},\mathcal{L}^{\otimes
mr})=H^{0}(\mathbb{A}^1, \mathcal{V})
$$

and lets you go back to you original object
via the valuation map

$$
\operatorname{ev}_1:H^{0}(\mathfrak{X},
\mathcal{
L}^{\otimes rm}) \to H^{0}(X, L^{ \otimes r m
}).
$$

A thing I don't understand good (perhaps
because it is very algebraic is the Rees
algebra that plays some important algebraic
role in all this.

\begin{theorem}
\label{theorem-montero-1}
There is a correspondene between test
configurations $(\mathfrak{X},\mathcal{ L})$
of $(X,L)$ and graded
$\mathbb{Z}$-filtrations of $R(X,L^{\otimes
r})$ for some $r>0$.
\end{theorem}

\begin{proof}
This is what we have been discussing. For
every test configuration we can cook up a
filtration in the algebra $R(X,L^{ \otimes
r}$ i.e. a filtration \textit{on each piece
of degree $n$}, that is we have

$$
(\mathfrak{X},\mathcal{L}) \rightsquigarrow
F^*_{\mathfrak{X},\mathcal{L}}R(X,L^{\otimes
r }
$$

And conversely, the whole point of all this
is that you recover the variety by taking the
$\operatorname{Proj}$ of this thing (of the
test configuration?). More precisely, we have
that
$\operatorname{Rees}(F^*_{\mathfrak{X},
\mathcal{L}}R(X,L^{\otimes
r }))$ given

$$
\mathfrak{X}:=\operatorname{Proj}_{m \in
\mathbb{N}}\left(\bigoplus_{m \in \mathbb{N}}
\bigoplus_{\lambda \in
\mathbb{Z}}F^\lambda_{\mathfrak{X},
\mathcal{L}}H^{0}(X,L^{\otimes
m r })E^\lambda\right) \to \mathbb{A}^1
$$

\end{proof}

\medskip\noindent{\bf Note.}
You don't need to understand the test
configuration: it's enough with understanding
the birational geometry of the variety (the
valuations on $X$).

\begin{corollary}
\label{corollary-montero-1}
let $(\mathfrak{X},\mathcal{L})$ be a test
configuration. $\mathfrak{g}(X,L)$. Then  if
$X$ is reduced and irreducible then so is
$\mathfrak{X}$.
\end{corollary}

\begin{proof}
This is beacuse the Rees algebra preserves
the commutative-algebraic properties of
$R[t,t^{-1}]$. That's fun.
\end{proof}

\begin{remark}
\label{remark-montero-3}
Similarly, one can prove that $X$ norml and
$\mathfrak{X}_0$ reduced implies
$\mathfrak{X}$ normal.
\end{remark}

\medskip\noindent{\bf Parenthesis:
K-stability and MMP.}

Let $(X,B)$ a pair, where $X$ is a normal
projective variety and $B=\sum a_i D_i$ a
$\mathbb{Q}$-Weil divisor on $X$ such that $
a_i \in [0,1]$ for all $i$ such that
$(K_X+B)$ is $\mathbb{Q}$-Cartier (so it's an
actual line bundle, if you like).

Now if $ f:Y \to X$ is a log resolution of
$(X,B)$ with exceptional divisors
$E_1,\ldots,E_k$, then what happens with the
pullback of that line bundle? Well there will
be an error term that basically measures how
singular is $B$:

$$
K_Y+f^{-1}_*B \sim_\mathbb{Q} f^*
(K_X+B)+\sum_{i=1}^k d_i E_i
$$

for some rational $b_i$. They are the things
that measure how singular is $B$ and are
called \textit{\textbf{discrepancies} }.

\begin{definition}
\label{definition-montero-2}
$(X,B)$ is \textit{\textbf{klt lc}} if

$$
\begin{cases}
d_i>-1\qquad &\forall \text{such $f$}  \\
d_i \geq -1\qquad &\forall \text{ such $f$}
\end{cases}
$$

and we say that $X$ is \textit{\textbf{klt}}
(resp. \textit{\textbf{lc}}) if $(X,0)$ is
\textit{\textbf{klt}} (resp
\textit{\textbf{lc}}).
\end{definition}

\begin{remark}
\label{remark-montero-4}
If $\dim X=2$, then klt=quotient
singularities.
\end{remark}

\begin{theorem}[Odaka '12, '13]
\label{theorem-montero-2}
Let $X$ be a normal projective variety, $L
\in \operatorname{Pic}(X)_\mathbb{Q}$ ample.
Then
\begin{enumerate}
\item  If $K_X \sim_\mathbb{Q} 0$ then $X$
klt (resp. lc) $\iff$ $(X,L)$ is $K$-stable
(resp. K-semistable).
\item $L=K_X$, $X$ lc $\iff (X,L)$ K-stable
$\iff (X,L)$ K-semistable.
\item $L=-K_X$, $(X,L)$ K-semistable
$\implies  X$ klt.
\end{enumerate}
\end{theorem}

Key ideas behind $1.$ (using the
Donaldson-Futaki invariant that was
constructed in some past lecture I didn't
attend): we saw that if
$(\overline{\mathfrak{X}},\overline{
\mathcal{L}})
\xrightarrow{\overline{\pi}} \mathbb{P}^1$ is
a ``compactified" test configuration of
$(X,L)$ then

$$
\operatorname{DF}(\mathfrak{X},\mathcal{L}) =
\frac{\overline{\mathcal{L}}^n\cdot
K_{\overline{\mathfrak{X}}/\mathbb{P}^1}}{V}+
\frac{\overline{S}\overline{\mathcal{L}}^{n+
1}}{(n+1)V}
$$

where $V=L^n$,
$\overline{S}=\frac{n}{V}(-K_X\cdot L^{n-1}$.

\begin{remark}
\label{remark-montero-5}
You can compute the DF invariant using this
$Y$: i.e. if we consider

$$
\xymatrix{&Y\ar[dl]_{f}\ar[dr]^{g}&\\
\overline{\mathfrak{X}}\ar@{-->}[rr]&&X\times
\mathbb{P}^1}
$$

then

$$
\operatorname{DF}(\mathfrak{X},\mathcal{L})=
\frac{\overline{\mathcal{L}}^nf_*(K_{Y/X
\times \mathbb{P}^1}+g^*
\operatorname{pr}_1^*K_X}{V}+\frac{
\overline{S}\overline{\mathcal{L}}^{n+1}}{(n+
1)V}
$$

(this was proved)
\end{remark}

\medskip\noindent{\bf Key observation.}
\begin{enumerate}
\item If $X$ is lc (resp. klt) then  $K_{Y/X
\times \mathbb{P}^1}$ is effective (resp.
effective $\neq 0$).
\item If $K_X \sim_\mathbb{Q} 0$ then
$$
\operatorname{DF}(\mathfrak{X},\mathcal{L})=
\frac{\overline{\mathcal{L}}^n\cdot
f_*(K_{y/X \times \mathbb{P}^1}}{V} \geq 0
\qquad \text{(resp. $>0$}
$$

if $X$ is lc (resp. klt), i.e. $(X,L)$ is
$K$-semistable (resp. K-stable).

\item For the other implication, the idea is
that if $X$ is not lc, then $\exists
(\mathfrak{X},\mathcal{L})$ test
configuration with
$\operatorname{DF}(\mathfrak{X},
\mathcal{L})<0$.

Here's a glimplse of what's going on here:
Odaka-Xu '12 shows the existence of lc
modules $(Y,\Delta X) \to X$. So take $E:=
K_{Y/X}+ \Delta X$. So we define
$\mathcal{I}=\mathcal{I}_E\subseteq
\mathcal{O}_{X \times \mathbb{A}^1}$, the
``flag ideal". You obtain the following test
configuration that breaks the K-stability.

$$
\mathfrak{X}:=
\operatorname{Bl}_\mathcal{I}(X \times
\mathbb{A}^1) \to \mathbb{A}^1
$$

and a suitable $\mathcal{L}$ such that
$\operatorname{DF}(\mathfrak{X},
\mathcal{L})<0$.
\end{enumerate}

\medskip\noindent{\bf Lecture 4.}

\medskip\noindent{\bf Test configurations and
valuations.}

{\bf Last time}  we discussed the
correspondence of test configurations and
some filtrations on the set of sections, this
is done via something called the ``Rees"
construction.

{\bf Today}  we will be more geometric.

Let $K/k$ be a field extension with
$\operatorname{tr.deg}_k(K)<+\infty$, e.g.
$K=k(X)$.

A  \textit{\textbf{(real) valuation}} on $K$
is a function $v:K^\times \longrightarrow
\mathbb{R}$ such that
\begin{enumerate}
\item $v(fg)=v(f)+v(g)\qquad \forall f,g \in
K^\times$.
\item $v(f+g) \geq
\operatorname{min}\{v(f),v(g)\}, \qquad
\forall f,g \in K^\times$.
\item $v|_{k^\times}=0$.
\end{enumerate}
We define $v(0):=+\infty$, and if $K=k(X)$,
$X$ valued var, then we write  $v \in
\operatorname{Val}_X$ for short.

\medskip\noindent{\bf Recall.}
For $v: K^\times \longrightarrow \mathbb{R}$
valuation, we define
\begin{enumerate}
\item $\mathcal{O}_v=\{f \in K, v(f) \geq
0\}$, local ring with $\mathfrak{m}_v:=\{f
\in K: v(f)>0\}$.
\item $K(v):= \mathcal{O}_v/\mathfrak{m}_v$
and
$\operatorname{tr.deg}(v)=
\operatorname{tr.deg}K(v)$.
\item $\Gamma_v:=v(K^\times) \subseteq
\mathbb{R}$, ``value group" and
$\operatorname{rat.rk}(v):=\dim_\mathbb{Q}(
\Gamma_v
\otimes_{\mathbb{Z}}\mathbb{Q})$.
\item (Abhyankar '56) If $v \in
\operatorname{Val}_X$,
$\operatorname{tr.deg}(v)+
\operatorname{rat.rk}(v)\leq
\dim(X)$.
\end{enumerate}

\begin{example}
\label{example-montero-1}
\begin{enumerate}
\item Let $x \in X$ be a smooth (closed)
point. We define for $f \in
\mathcal{O}_{X,x}$,

$$
\operatorname{ord}_x(f):=
\operatorname{max}\{ d \in \mathbb{N}, f \in
\mathfrak{m}^d_x\}
$$

and we can extend it to

$$
\operatorname{ord}_x:k(X)^\times
\longrightarrow \mathbb{Z}
$$

This is a \textit{\textbf{discrete
valuation}}.

\item $X=\mathbb{A}^2_{(x,y)}$, $K=k(x,y)$
and let us fix  $\alpha,\beta \in
\mathbb{R}^2\setminus \{(0,0)\}$. Given $f =
\sum_{m,n \in \mathbb{N}}\lambda_{m,n}x^my^n
\in \mathcal{O}(\mathbb{A}^2)$ we define

$$
v(f):=\operatorname{min}\{ \alpha m + \beta
n: \lambda_{m,n}\neq  0\}
$$
and we obtain $v:K^\times \to \mathbb{R}$.

\item \textbf{(Divisorial valuation (most
important example))} First recall that a
\textit{\textbf{divisor}} $E$ over a variety
$X$ is a proper birational map

$$
\mu:Y \longrightarrow X
$$

with $Y$ normal, and $E \subseteq Y$ a prime
divisor. (``$\mu$ extracts the divisor $E$")
{(dani: what is a \textit{prime} divisor?)}.
In particular, $\mathcal{O}_{Y,E}$ is a DVR
with associated discrete valuation
$$
\begin{aligned}
\operatorname{ord}_E: K^\times
&\longrightarrow \mathbb{Z} \\
f &\longmapsto \operatorname{ord}_E(\mu^*f)
\end{aligned}
$$
where $K=k(X) \overset{\mu^*}{\cong}k(Y)$.
{(dani: so probably a valuation is the degree
of the polynomial right?)}

A valuation $v \in \operatorname{Val}_X$ of
the form $v = c \cdot \operatorname{ord}_E$,
$c \in \mathbb{R}^{>0}$ is called a
\textit{\textbf{divisorial valuation}}.

\end{enumerate}
\end{example}

\begin{definition}
\label{definition-montero-3}
Let $v \in \operatorname{Val}_X$. The
\textit{\textbf{center of $v$}} is the
(schematic) point $\xi=c_x(v) \in X$ such
that $v \geq 0$ on $\mathcal{O}_{X,\xi}$ and
$v>0$ on  $\mathfrak{m}_\xi$. It exists
(resp. is unique) if $X$ is proper (resp. if
$X$ separated).
\end{definition}

\begin{example}
\label{example-montero-2}
$x \in X$ a smooth point again,
$\operatorname{ord}_x$ is divisorial.
\end{example}

These notes are not complete but here's an
important theorem:

\begin{theorem}[Zerisky]
\label{theorem-montero-3}
Let $v \in \operatorname{Val}_X$ be a (real)
valuation. Them if  $n=\dim X$, $v$ is
diviorial iff $\operatorname{rat.rk}(v)=1$
and $\operatorname{tr.deg}v=n-1$.
\end{theorem}

\medskip\noindent{\bf Back to
configurations.}

Recall that $\mathfrak{X},\mathcal{L}$ is a
deformation of $(X,L)$ and notice that since
$\mathfrak{X}
\overset{\operatorname{bir}\cong}{
\dashrightarrow} X \times \mathbb{A}^1$ we
have

$$
k(\mathfrak{X}) \cong k(X)(t)
$$

That is, rational functions on deformation
space look like rational functions on base
with an extra variable.

\begin{theorem}[see. eg. Jonsson-Mustata ;12]
\label{theorem-montero-4}
Let $k \subseteq K' \subseteq K$ field
extension and $w: K^\times \longrightarrow
\mathbb{R}$ valuation and let  $v:=w
|_{(K')^\times}$. Then

$$
\operatorname{tr.deg}(w)+
\operatorname{rat.rk}(w)
\leq
\operatorname{tr.deg}(v)+
\operatorname{rat.rk}(v)+
\operatorname{tr.deg}_{K'}(K).
$$

\end{theorem}

\medskip\noindent{\bf Consequence.}
Let $w$ be a valuation of $k(X)(t)$. If $w$
is divisorial, then its restriction
$v:=r(w)=w|_{k(X)^*}$ is divisorial or
trivial.

Which is very nice because, as we will see,
we divisorial valuations are nice. They will
tell us things about stability.

\medskip\noindent{\bf Recall.}
Let $\mathbb{G}_m \curvearrowright
(\mathfrak{X},\mathcal{L})\xrightarrow{\pi}
\mathbb{A}^1$
be a \textbf{normal} test configuration of
$(X,L)$. Since $\mathfrak{X}\setminus
\mathfrak{X}_0
\overset{\mathbb{G}_m\text{-equiv.} }{\cong}X
\times (\mathbb{A}^1 \setminus\{0\}$, we can
consider the normalization of the graph of
$\mathfrak{X} \dashrightarrow X \times
\mathbb{A}^1$ to obtain

$$
\xymatrix{&Y\ar[dl]_{f}\ar[dr]^{g}&\\
\mathfrak{X}\ar@{-->}[rr]&&X\times
\mathbb{A}^1}
$$

where $f^*\mathcal{L} \sim_\mathbb{Q}g^*
(L_{\mathbb{A}^1}+D$ for some
$\mathbb{Q}$-Cartier divisor $D$ such that
$\operatorname{supp}(D) \subseteq Y_0$.
{(dani: looks like this ``test configuration"
is something like a blow-up: it has an
exceptional fiber $\mathfrak{X}_0$, which is
drawn as weird entity with singularities,
while the rest of the fibers look smooth.)}

\medskip\noindent{\bf dani recalls.}
that the strict transform of the blow up of
the cuspidal curve $Y:y^2=x^3$ inside
$\mathbb{A}^2$ is the closure of
$\varphi^{-1}(Y-O)$ where $\varphi:X \to
\mathbb{A}^2$ is the blow-up map.

\medskip\noindent{\bf Key observation.}
Every $F \subseteq \mathfrak{X}_0$
irreducible component of $\mathfrak{X}_0$
(different from the strict transorfm of $X
\times \{0\}$) induces a divisorial valuation
$\operatorname{ord}_F: k (X)(t)^\times
\longrightarrow \mathbb{Z}$.

We denote by $v_F:= r
(\operatorname{ord}_F)\overset{
\operatorname{def}}{=}
\operatorname{ord}_F |_{k(X)^\times}$ and
\textbf{we observe} that $v_F$ is a
divisorial valuation on $k(X)^*$, i.e.,
$v_F=c\cdot \operatorname{ord}_E$ for some $c
\in \mathbb{N}^{\geq 1}$ and some $E
\subseteq Y \xrightarrow{\mu}X$ divisorial
over $X$. (For a proof check the paper we
have cited several times BHJ'17).

The following proposition allows us to forget
about filtrations and work with valuations
instead (which are more geometric things).

\begin{proposition}
\label{proposition-montero-1}
Let $m \in \mathbb{N}^{\geq 1}$ such that
$\mathcal{L}^{\otimes m} \in
\operatorname{Pic}(\mathfrak{X})$. Then, for
every $\lambda \in \mathbb{Z}$ we have

$$
F^\lambda_{\mathfrak{X},\mathcal{L}}H^{0}(X,
L^{\otimes
m})=\bigcap_{\substack{F \subseteq
\mathfrak{X}_0 \\ \text{irr. comp.} }}\{s \in
H^{0}(X,L^{\otimes m}):v_F(s)+m
\operatorname{ord}_F(D) \geq \lambda
\operatorname{ord}_F(t)\}
$$

\end{proposition}

*Some computations regarding
$F^\lambda_{\mathfrak{X},\mathcal{L}}$*

\begin{definition}
\label{definition-montero-4}
Given $v: k(X)^* \longrightarrow\mathbb{R}$
divisorial valuation, we put

$$
F^\lambda_v:=\{s \in H^{0}(X,L),s(v) \geq
\lambda\}
$$

\end{definition}

\begin{definition}[K. Fujta '16]
\label{definition-montero-5}
We say that $v=c\cdot
\operatorname{ord}_E\cdot k(X)^*
\longrightarrow \mathbb{Z}$ is a
\textit{\textbf{dreamy valuation}} (and $E$
is a \textit{\textbf{dreamy divisor}})
if\ldots{} (e.g. if $Y$ is log-Fano
[BCHM'10]---they proved these are Mori
dream.)
\end{definition}

\begin{theorem}
\label{theorem-montero-5}
Let $(X,L)$ such that $X$ is Fano with klt
singularities (e.g. quotient singularities)
and $L=-K_X$. There is a bijection between
\begin{enumerate}
\item Normal test configuration
$(\mathfrak{X},\mathcal{L})$ with
$\mathcal{L}=-K_{\mathfrak{X}/\mathbb{A}^1}$
and $\mathfrak{X}_0$ reduced and irreducible.
\item $v:k(X)^* \longrightarrow \mathbb{Z}$
dreamy valuations.
\end{enumerate}
\end{theorem}

\medskip\noindent{\bf !}
Li and Xu (2014): In order to check
K-stability it is enough to consider
``special test configurations", i.e. with
$\mathcal{L}=-K_{\mathfrak{X}/\mathbb{A}^1}$
and $\mathfrak{X}_0$ (reduced, irreducible)
klt Fano variety.

\begin{proof}[Proof of the theorem (not Li
and Xu's)]
Sketched.
\end{proof}

\begin{remark}
\label{remark-montero-6}
While the variety $\mathfrak{X}_0$ seems to
be very particular, it turns out that these
conditions are enough to compute K-stability.
\end{remark}

\medskip\noindent{\bf Lecture 6.}

\medskip\noindent{\bf Abbam-Zhuang
inequality.}

\textbf{Recall} If we have $X$ Fano klt,
$\dim X=n$ (this will be the same for the
whole talk). Then, by the valuative
criterion, we have that

$$
X \text{ is }\begin{cases}
\text{K-stable} \qquad & \\
\text{K-semistable} \qquad &
\end{cases}\iff \delta(X)=\begin{cases}
>1\qquad & \\
\geq 1\qquad &
\end{cases}
$$

where

$$
\delta(X)=\operatorname{inf}_{\substack{E
\subseteq Y \xrightarrow{\mu}X \\
\text{divisor over $X$}
}}\frac{A_X(E)}{\delta_X(E)}
$$

Actually (cf Chesltsow-Rubinstein-Zhang '19),
if $X$ is a smooth Fano variety:
\begin{itemize}
\item
$\operatorname{min}\{1,\delta(X)\}=
\operatorname{sup}\{\lambda
\in [0,1], \exists \omega \in
c_1(X)\text{K\"ahler st}
\operatorname{Ric}(\omega) > \lambda
\omega\}$.
\item $\operatorname{Vol}(X)=(-K_X)^n$
($\leq_{\text{K-ss} }(n+1)^n$) $\implies
\operatorname{Vol}(X)^{1/n}\leq
\frac{(n+1)}{\delta(X)}$.
\item (Blum-Jonsson '21.) Assume $X$ is
\textit{toric} Fano klt.

$$
X \overset{\text{theory of} }{
\leftrightsquigarrow}_{\text{toric varieties}
}\text{some combinatorial data $v_i$}.
$$

So for a divisor $D=b_1D_1+\ldots+b_dD_d$ on
$X$,

$$
P_D:=\{u \in \mathbb{R}^n: \left\langle{}
u,v_i \right\rangle{}\geq -b_i \quad \forall
i\}
$$

Then the theorem is that for $D=-K_X$, i.e.
$b_i=1 \forall i$,

$$
\delta(X)=\operatorname{min}_{i=1,\ldots,d}
\frac{1}{\left\langle{}
\bar{u},v_i \right\rangle{}+1},\qquad \bar{u}
\in P_{-K_X}.
$$

In particular, if $X$ is smooth toric Fano,

$$
X \text{ KE} \iff \bar{u} =0\qquad
\text{(Wang-Zhu '04}
$$

\begin{example}
\label{example-montero-3}
For the following toric varieties we get
$\delta(\mathbb{P}^n)=1$,
$\delta(\operatorname{Bl}_p\mathbb{P}^2=
\frac{6}{7}$,
$\delta(\operatorname{Bl}_{p,q}\mathbb{P}^2=
\frac{21}{25}$.
\end{example}
\end{itemize}

Let us consider $(X,\Delta)$ log Fano pair,
$\Delta=\sum b_i\Delta_i$, $\mathbb{Q}$-Wil
divisor, $b_i \in [0,1)$ and $-(K_X+\Delta)$
ample $\mathbb{Q}$-Cartier
$\mathbb{Q}$-divisor. We define for $E
\subseteq Y \xrightarrow{\mu}X$,

$$
\delta(X,\Delta;E):=\frac{A_{(X,
\Delta)}(E)}{S_{(X,\Delta)(E)}}
$$

where
$S_{(X,\Delta)}(E)=\frac{1}{(-K_X-\Delta)^n
\int_0^{+\infty}\operatorname{Vol}(-\mu
^*(K_X+\Delta)-tE)dt}$.

We say that $(X,\Delta)$ is
\textit{\textbf{K-stable}} (resp.
\textit{\textbf{K-semistable}}) if

$$
\delta(X,\Delta)=\operatorname{inf}_{E
\subseteq Y
\xrightarrow{\mu}X}\delta(X,\Delta;E)
$$

is $>1$ (resp. $\geq 1$).

\begin{remark}
\label{remark-montero-7}
Even if you don't care about pairs, you may
get some information when you have a nice
cover. More explicitly, a theorem by Liu-Zhu
'22 says if $\pi:(X,D)\to (Y,B)$ is a finite
surjective morphism of log-Fano pairs such
that \ldots{}

As an example, if $\pi:X \to Y$ is a double
cover with $X,Y$ smooth, then $\pi^*
(K_Y+\frac{1}{2}B=K_Y+0$. Thus $(X,0)$ K-ps
iff $(Y,\frac{1}{2}B)$ is K-ps.
\end{remark}

\begin{remark}
\label{remark-montero-8}
For a normal projective variety $X$ and $L
\in \operatorname{Pic}(X)$ there is a nice
way to compute $S(V_\bullet,E)$ using certain
filtration.
\end{remark}

\medskip\noindent{\bf Abbam-Zhuang inequality
'22.}

\textbf{Objective:} prove the K-semistability
of cubic surfaces.

Also see e.g. ``The Calabi problem for Fano
3-folds" \S{}1.7.

There is a technical condition: if
$(X,\Delta)$ is log-Gano and $E \subseteq Y
\xrightarrow{\mu}X$, a
\textit{\textbf{divisor of plt type}}\ldots{}

The result by Abban-Zhuang is an
``adjunction" inequality for
$\delta(X,\Delta)$ and $E$.

It's kind of like this:

$$
\xymatrix{E\subseteq Y\ar[d]^{\mu}\\ X}
$$

you consider the thing that $E$ is sent to (a
point, for example) (that is called the
\textit{\textbf{center}} of $E$). You
consider a divisor contained in the center of
$E$, and then a thing called $F$ over a
divisor $Z'$ on $E$. Right the thing is that
with these data we can compute the invariant
$\delta_{Z'}(E,\Delta_E;W^E_{\bullet,
\bullet}$.
The caveat is that for this expression we
cannot use the $S$-invariant defined above
(that was defined using some filtration on
the sections of a projective variety and a
line bundle $L$). Now we need to define a
more subtle $S$-like invariant that depends
on two parameters.

\medskip\noindent{\bf Disgression:
Newton-Okunkov body of a divisor.}

This is a way to produce a convex body out of
a divisor. We did the example of
$\mathbb{P}^2 \supseteq Y_1=\{x=0\}\supseteq
Y_2=\{x=y=0\}$, with the divisor $D=$ line
(e.g.
$\mathcal{O}_{\mathbb{P}^2}=\mathcal{O}(1)$.
Indeed, the construction uses a flag
$Y_\bullet$ of subvarieties of codimension 1
each wrt the next, and gives the vertices of
the convex body using the order of vanishing
of the sections of the multiples of the
divisor. We concluded that
$\Delta_{Y_\bullet}(D)$ in our example is a
right-angle triangle (the convex hull of
$0,e_1,e_2$ up to rescaling).

\begin{remark}[Zikas]
\label{remark-montero-9}
This definition depends on the basis of the
section vector spaces! You could take a basis
with vanishing order equal to zero\ldots{}
\end{remark}

The definition ends up being

$$
\Delta_{Y_{\bullet}}(D):=\bigsqcup_{m \geq
1}\frac{1}{m}\nu(H^{0}(X,mD)).
$$

\begin{theorem}[Lazrsfeld-Mustata '09]
\label{theorem-montero-6}

$$
\operatorname{Vol}_{\mathbb{R}^n}(\Delta_{Y^
\bullet}(D))=\frac{1}{n!}
\operatorname{Vol}_X(D)
$$

\end{theorem}

\begin{example}
\label{example-montero-4}
For a smooth projective surface $S$ with
$Y_\bullet:S \supseteq C \supseteq \{
p\}$\ldots{}
\end{example}

\medskip\noindent{\bf Back to Abbam-Zhuang.}

$X$ smooth projective surface, $E \subseteq Y
\xrightarrow{\mu}X$ smooth divisor (cuurve).
We will estimate $\delta_p(X)$ for $p \in X$
a point. The point is that for dimensional
reasons the case of surfaces is more simple.
We get a formula for the double-indexed
invariant $S_{\bullet,\bullet}^E$.

The moral is that we compute this number by
integrating slices along the Newton-Okunkov
body. And of course this goes back to the
``adjuntion"-like formula we discussed
before. Now we chall compute de $\delta$
invariant for a special case:

\begin{example}
\label{example-montero-5}
$X$ smooth cubic on $\mathbb{P}^3$. Set $p
\in X$ and let us prove that  $\delta_p(X)
\geq  1$. Then $X$ is K-semistable.

Vamos. Consider $E \in |-K_X|$, which is
particular divisor i.e. $E \subseteq Y =X$, a
smooth elliptic curve such that $p \in E $
and $E |_{E}=3p$. In this case we get
\begin{itemize}
\item $A_X(E)=1$ (I think I might have missed
the definition of this $A$\ldots{}), ($E$
smooth). \ldots{} some formulas \ldots{} and
we compute the integral and conclude that

$$
S(W_{\bullet,\bullet}^E;p)= \int \ldots =1.
$$

On the other hand,

$$
S_X(E)=\text{this one is simpler to compute
(no Newton-Okuknov body integration)}
=\frac{1}{3}
$$

By Abbam-Zhuang,

$$
\delta_p(X) \geq \operatorname{min}\left\{
\frac{1}{1/3},1 \right\} =1.
$$

\end{itemize}
\end{example}

\medskip\noindent{\bf Editorial note.} The
notes cover Lectures 3, 4 and 6 only (the
numbering is Daniel's). In the original, the
action symbol typed as ``mathbb y'' is not
a valid symbol; it is replaced by
$\curvearrowright$, which I take to be the
intended group action. A commented-out
fragment of the original (a diagram and the
start of a remark attributed to Fujita '16
and
BHJ'17, Section 4, with an unfinished
formula)
is omitted.

\section{Jets, arcs and singularities}
\label{section-badur-jets-arcs}
%% speaker: Nero Badur, KU Leuven
%% date: February 20, 2025

\noindent
Nero Badur, KU Leuven.
Series and place not recorded in the notes.
February 20, 2025.

\medskip\noindent
No abstract available.

\medskip\noindent
Daniel's note: This is a talk that I couldn't
follow very far.

\medskip\noindent{\bf Introduction.}

\begin{definition}
\label{definition-badur-1}
An \textit{\textbf{arc}} on a
$\mathbb{C}$-algebraic variety $X$ is a
morhpism of $\mathbb{C}$-schemes
$\operatorname{Spec}\mathbb{C}[[t]] \to X$.
Equivalently, if $X = \operatorname{Spec} R$,
a morphism of $C$-algebras $R \to
\mathbb{C}[[t]]$.

The \textit{\textbf{arc space}} is
$\mathcal{L}_\infty(X) =
\operatorname{Hom}_{\mathbb{C}-
\operatorname{sch}}(\operatorname{Spec}
\mathbb{C}[[p ]],X)$.
\end{definition}

\medskip\noindent{\bf Affine case.}
$$
\begin{aligned}
\mathcal{L}_{\infty}(\mathbb{A}^n)&=
\operatorname{Hom}_{\mathbb{C}-
\operatorname{alg}}(\mathbb{C}[x_1,\ldots,
x_n],\mathbb{C}[
[ t ] ]) \cong \mathbb{A}^\infty=
\operatorname{Spec} \mathbb{C}[x_{ij}:0 \leq
i \leq n
\end{aligned}
$$

\begin{definition}
\label{definition-badur-2}
An \textit{\textbf{$m$-jet}} is just
truncating this, i.e., a morphism
$\operatorname{Spec}\mathbb{C}[t]/(t^{m+1})
\to X$. And \textit{\textbf{jet-space}} is

$$
\mathcal{L}_m(X)=
\operatorname{Hom}_{\mathbb{C}-
\operatorname{sch}}(\operatorname{Spec}(
\mathbb{C}[t]/(t^{m+1},X)
$$

\end{definition}

\textbf{Truncation gives morphisms.}

$$
\xymatrix{\mathcal{L}_\infty(X)\ar[r]
\ar@/^1pc/[rrr]^{\pi_m}&\cdots\ar[r]&
\mathcal{L}_m(X)\ar[r]&\mathcal{L}_{m-1}(X)
\ar[r]&\cdots\ar[r]&\mathcal{L}_0(X)=X}
$$

where $\pi_0$, which goes all the way to
$\mathcal{L}_0$ ``takes to the center".
\medskip\noindent{\bf Example.}
If $X$ is smooth (or $X= \mathbb{A}^n$, then
the  $\pi_{m,m-1}$ are locallu trivial
$\mathbb{A}^n$-fibrations.

{We are interested in singular subvarieties
of a smooth manifold.}

\begin{definition}
\label{definition-badur-3}
Let $f \in
\mathbb{C}[x_1,\ldots,x_n]\setminus
\mathbb{C}$. The \textit{\textbf{$m$-contact
locus of $f$}} is

$$
\mathfrak{X}_{m}^\infty(f):=\{\gamma \in
\mathcal{L}_\infty(\mathbb{A}^n):
\operatorname{ord}_\gamma f=m\}
$$

so
\end{definition}

\medskip\noindent{\bf Editorial note.} The
notes end here: the last definition is
unfinished in the original, and in the
original source the whole section was
commented out.

\section{Automorphisms of quartic surfaces
and Cremona transformations}
\label{section-paiva-quartic-cremona}
%% speaker: Daniela Paiva, IMPA
%% place: PhD thesis defence
%% date: February 25, 2025

\noindent
Daniela Paiva, IMPA.
PhD thesis defence.
February 25, 2025.

\medskip\noindent
No abstract available.

\medskip\noindent{\bf Introduction.}

The ambient space is $\mathbb{C}P^{n}$.
Recall definitions rojective algebraic
varieties (polynomials are irreducible),
hypersurfaces, automorphisms and birational
maps.

$$
\operatorname{Aut}(\mathbb{P}^{n+1},X):=\{
\varphi\in
\operatorname{Aut}(\mathbb{P}^{n+1}):
\varphi(X)=X\}
$$

\medskip\noindent{\bf dani.}
Cu\'ales automorfismos de la subvariedad son
restricciones de automorfismos del espacio
ambiente. De eso se trata la tesis.

\begin{theorem}[Matsumura-Monsky 1964]
\label{theorem-paiva-1}
restriction map is surjective as long as
$(n,d) \neq (1,4), (\text{ops!,ops!})$
\end{theorem}

Cremona group

$$
\operatorname{Bir}(\mathbb{P}^{n+1})
\supset\operatorname{Aut}(\mathbb{P}^{n+1})
$$

\medskip\noindent{\bf Gizatullin problem.}
Is $\Psi: \operatorname{Bir}(\mathbb{P}^3,S)
\longrightarrow \operatorname{Aut}(S)$
surjective?

\medskip\noindent{\bf Oguiso's question
(2012).}
Does any finite order automorphism $f \in
\operatorname{Aut}(S)$ belong to
$\operatorname{img}\Psi$?

\medskip\noindent{\bf Main result (P.-Quedo
2023, Aaraujo, P., Zikas 2024).}
\begin{enumerate}
\item[(i)] If $r(S)> 57$ or $r(S)=52$, then
$\operatorname{Bir}(\mathbb{P}^3,S)=\{1\}$;
\item[(ii)] If $r(S) \leq 57$, $r(S) \neq 52$
and $S$ is $\operatorname{Aut}$-general, then
$\Psi$ is surjective.
\end{enumerate}

\medskip\noindent{\bf K3 surfaces.}

Quartic surfaces are K3 surfaces!

\begin{definition}
\label{definition-paiva-1}
A \textit{\textbf{K3 surface}} is a smooth
projective algebraic surface with

$$
H^{0}(S,\Omega^{2}_S) = C_{\omega_S}\qquad
\text{and} \qquad h^1(S,\mathcal{O}_S)=0
$$

\end{definition}

They have lattices. The Picard lattice of a
K3 is of rank 22, though different K3s have
different lattices.

Another interesting property of K3 is that
automorphisms of the surface can be studied
via automorphisms of the lattice via global
Torelli theorem:

\begin{theorem}[Global torelli]
\label{theorem-paiva-2}

$$
\operatorname{Aut}(S)\cong
\{\phi:H^{2}(S,\mathbb{Z})\to
H^{2}(S,\mathbb{Z}):\text{preserving Hodge
structure and }\operatorname{Amp}(S) \}
$$

(The ample cone.) ``$S$ is
$\operatorname{Aut}$-general".

$$
\operatorname{Aut}(S)
\cong\{mk(\phi,\varepsilon) \in
O(\operatorname{Pic}(S) ) \times :{\text{some
conditions not so few} }\}
$$

\end{theorem}

\begin{example}[Case $\rho(S)=2$]
\label{example-paiva-1}
Setup: $S \subset \mathbb{P}^3$, smooth
quartic surface with $\rho(S)=2$, $H$ the
hypeplane section\ldots{} there is an
expresion for the intersection form (a $2
\times 2$ matrix).
\end{example}

\begin{proposition}[Galluzzi-Lobardo-Peters
2010]
\label{proposition-paiva-1}
$S \subset \mathbb{P}^3$ Aut-general with
$\rho(S)=2$,

$$
\operatorname{Aut}(S)=\begin{cases}
\{1\}\qquad & \text{ if $\exists D$ with $D^2
\in \{-2,0\}$} \\
\mathbb{Z}_2\qquad &\\
\mathbb{Z} & \text{ if $\not \exists D$ with
$D^2 \in \{-2,0\}$} \\
\mathbb{Z}_2 * \mathbb{Z}_2
\end{cases}
$$

\end{proposition}

\begin{proposition}
\label{proposition-paiva-2}
$S \subset \mathbb{P}^3$ with $\rho(S)=2$. We
have:
\begin{enumerate}
\item  If $f \in \operatorname{Aut}(S)$ has
finite order, then $f$ is an involution;
\item There is a bijection between
involutions of $S$ and ample class $A$ with
$A^2=2.$
\end{enumerate}
\end{proposition}

\begin{proposition}
[Galluzzi-Lobardo-Peters+P.-Quedo
2025]
\label{proposition-paiva-3}
$S \subset \mathbb{P}^3$ Aut-general with
$\rho(S)=2$,

$$
\operatorname{Aut}(S)=\begin{cases}
\{1\}\qquad & \text{ if $\exists D$ with $D^2
\in \{-2,0\}$ and $\not\exists  A$ ample with
$A^2=2$} \\
\mathbb{Z}_2\qquad &\text{ if $\exists D$
with $D^2 \in \{-2,0\}$ and $\exists  A$
ample with $A^2=2$}\\
\mathbb{Z} & \text{ if $\exists D$ with $D^2
\in \{-2,0\}$ and $\not\exists  A$ ample with
$A^2=2$}\\
\mathbb{Z}_2 * \mathbb{Z}_2 & \text{ if
$\exists D$ with $D^2 \in \{-2,0\}$ and $ A$
ample with $A^2=2$}
\end{cases}
$$

\end{proposition}

\medskip\noindent{\bf Sarkisov program.}

Start with a projective variety. The output
of MMP is a simpler variety with ?.

\begin{theorem}[Sakisov Program]
\label{theorem-paiva-3}
The  \textit{\textbf{Sarkisov links}}
$\varphi_i$ are

$$
\xymatrix{X_1=Y_0\ar@{-->}[r]^{\varphi_1}
\ar@/^2pc/@{-->}[rrrr]^{\varphi}&Y_1
\ar@{-->}[r]^{\varphi_2}&Y_2\ar@{-->}[r]&
\cdots\ar@{-->}[r]&Y_n=X_2}
$$

and $Y_i$ are Mori-dream (I think).
\end{theorem}

\begin{definition}[Calabi-Yau pair]
\label{definition-paiva-2}
\end{definition}

\begin{definition}[Volume preserving]
\label{definition-paiva-3}
\end{definition}

Volume preserving Sarkisov program.
Classification of Sarkisov links.
hey there

\begin{proposition}
\label{proposition-paiva-4}
$X$ is weak Fano and $X \to \mathbb{P}^3$
initiaties a Srkisov link iff $|4H-C|$ is
ample and base pont free on $S$.
\end{proposition}

\begin{proposition}[Araujo-Corti-Massarenti
2023]
\label{proposition-paiva-5}
$X \to \mathbb{P}^3$ yields an antiflip
$\implies  \exists \Gamma \subset S$ such
that $\forall \Gamma\cdot (4H-C)=-1$
\end{proposition}

\begin{proposition}
\label{proposition-paiva-6}
$X \to \mathbb{P}^3$ initiates a Sarkisov
link $\implies X $ is weak Fano iff $(g,d)
\in$ ? \ldots{}
\end{proposition}

\textbf{Result:}
\begin{theorem}[P.-Quedo 2023,
Araujo-P.-Sikas 2024]
\label{theorem-paiva-4}
?
\end{theorem}
Explanation: found the automorphism group of
$S$, found generators of
$\operatorname{Aut}(S)$ and construct Cremona
transformations realizing them.

Consider the example when the lattice has
rank 17. Then
$\operatorname{Aut}(S)=\mathbb{Z}_2=\left
\langle{}
\tau \right\rangle{}$. Let $C \subset S$ be
of type $(g,d)=(14,11)$.

Here's a result by Zikas (oops, looks like I
mixed it with another diagram by
Cutrone-Masburn 2013):

$$
\xymatrix{X\ar@{-->}[r]^{\tilde{\varphi}}
\ar[d]_{\operatorname{Bl}_C}&X_1\ar[d]^{
\operatorname{Bl}_{C_1}}\\
\mathbb{P}^3\ar@{-->}[r]_{\varphi}&{
\mathbb{P}^3?}}
$$

Anyway this is a sketch of how the proof is
done.

\textbf{Conclusion:}

\medskip\noindent{\bf Oguiso's question
(2012).}
Does any finite order automorphism $f \in
\operatorname{Aut}(S)$ belong to
$\operatorname{img}\Psi$?

$$
\text{No!}
$$

We (they) have constructed counter examples,
for example take rank = ?

\medskip\noindent{\bf Gizatullin's problem
for $\rho(S) \geq 3$.}

We start by analyzing Oguiso's example for a
smooth quartic surface with
$\operatorname{Pic}((S))=\mathbb{Z}H \oplus
\mathbb{Z}L \oplus  \mathbb{Z}M$. (Oguiso
didn't use Sarkisov program.)

\textbf{Idea}: Study how $S,
\operatorname{Aut}(S)$ and the possible
Sarkisov links deform.

Other problem: K3 surfaces in other ambient
spaces! So let $S \subset X$ where $X$ is a
Fano 3-fold with $\rho(X)=1$ and index $s$,
$S$ a K3 surface. We ask again: is every
automorphism of this K3 surface induced by a
birrational map of the ambient space? That
is,

$$
\text{ is } \qquad
\Psi:\operatorname{Bir}(X,S) \longrightarrow
\operatorname{Aut}(S)\qquad
\text{surjective?}
$$

And if $(X,S)$ is a Calabi-Yau pair we can
use Sarkisov program.

\begin{proposition}
\label{proposition-paiva-7}
Assume $\operatorname{Pic}(S) = $ no s\'e
qu\'e, then
$\operatorname{Bir}=\operatorname{Aut}$.
\end{proposition}
(This is an ongoing project with Michaela ?,
and ?

\medskip\noindent{\bf Mais uma: Intertia and
decomposition groups.}

$S \subset \mathbb{P}^3$ smooth quartic
surface

$$
0 \to \operatorname{In}(S) \to
\operatorname{Dec}? \to ?
$$

$$
0 \to \operatorname{In} \to
\operatorname{Aut}? \to \operatorname{Bir}?
$$

Anyway this one is with Zokratis and
Carolina.

\medskip\noindent{\bf Editorial note.} The
diagrams drawn with tikz-cd in the original
are written as xymatrix diagrams.

\section{Higher dimensional Fano varieties}
\label{section-moraga-higher-dim-fano}
%% speaker: Joaquín Moraga, UCLA, USA
%% place: minicourse, V-ELGA (CIMPA Cabo
%% place: Frío)
%% date: August 12-16, 2025

\noindent
Joaqu\'in Moraga, UCLA, USA.
V-ELGA (CIMPA Cabo Fr\'io), minicourse.
August 12-16, 2025 (as in the notes).

\medskip\noindent
{\bf Abstract.} This will be a 6 hour
mini-course about Fano varieties. We will
start with the classic classification of del
Pezzo (smooth Fano) and smooth Fano 3-folds
(Iskhoskikh-Prokhorov). This will be an
overview of the known results and a highlight
of why understanding Fano varieties is
important for Algebraic Geometry. Then we
will introduce Kawamata log terminal
singularities and discuss some classic and
new results about this class of
singularities. We will explain why
understanding these singularities is vital,
for instance, through the classification of
Gorenstein Fano surfaces of Picard rank one.
Finally, we will explore some new results
regarding Fano varieties and klt
singularities, discussing the existence of
complements on Fano varieties and the
boundedness of Fano varieties.

\medskip\noindent{\bf Plan.}

\begin{enumerate}
\item Introduction to Fano varieties.
\item Log terminal singularities.
\item Singular Fano varieties.
\item Cluster Fano varieties.
\end{enumerate}

\medskip\noindent{\bf Lecture 1.}

$X$ smooth projective variety of dimension
$n$. $T_{X}$ its tangent bundle.
$\Omega_{X}:=T^{*}_{X}$. The
\textit{\textbf{canonical line bundle}} is
$\omega_{X}=\Lambda^{n} \Omega_{X}$.
$\omega_{X}\cong \mathcal{O}_{X}(K_{X})$
where $K_{X}$ is the canonical divisor.

\medskip\noindent{\bf Trichotomy.}

\begin{enumerate}
\item $X$ is \textit{\textbf{Fano}} if
$\omega_{X}^{-m}$ is very ample for some $m$
divisible enough.
\item $X$ is \textit{\textbf{Calabi-Yau}} if
$\omega_{X}^{m} \cong  \mathcal{O}_{X}$.
\item $X$ is \textit{\textbf{canonically
polarized}} if $\omega_{X}^{m}$ ($m\geq 0$)
is very ample for $m$ divisible enough.
\end{enumerate}

\begin{example}
\label{example-moraga-1}
$X_{d}\subset \mathbb{P}^{n}$ is a smooth
hypersurface of degree $d$. Then
$$
\begin{aligned}
\omega_{X_{d}}&=(\omega_{\mathbb{P}^{n}}+
\mathcal{O}(X_{d})|_{X_{d}}\\
&((-n-1)H+dH)|_{X_{d}}\\
&=(J-1-n)H|_{X_{d}}
\end{aligned}
$$
\begin{itemize}
\item $d<n+1$,  $X_{d}$ is Fano.
\item $d=n+1$, $X_{d}$ is CY.
\item $d>n+1$,  $X_{d}$ is canonically
polarized.
\end{itemize}
\end{example}

So the idea is start from $\mathbb{P}^{2}$
and blow-up up to 8 points.

\begin{theorem}
\label{theorem-moraga-1}
Any smooth Fano surface (del Pezzo) is either
isomorphic to $\mathbb{P}^{1} \times
\mathbb{P}^{1}$ or the blow-up of
$\mathbb{P}^{2}$ at $k\leq 8$ points in
general position.
\end{theorem}

\begin{definition}
\label{definition-moraga-1}
The \textit{\textbf{volume}} of a smooth Fano
surface  $X$ is $(-K_{X})^{2}$.
\end{definition}

\begin{tabular}{ccc}
smooth fano surface&volume&automorphism
group\\\hline
$\mathbb{P}^{2}$ &9&
$\mathbb{P}\operatorname{GL}(3,
\mathbb{C})$\\\hline
$\operatorname{Bl}_{p}\mathbb{P}^{2},
\mathbb{P}^{1} \times \mathbb{P}^{1}$
&8&connected group of rank $\geq
2$/$\mathbb{P}\operatorname{GL}(2,\mathbb{C})
\times
\mathbb{P}\operatorname{GL}(2,\mathbb{C})
\ltimes
\mathbb{Z}_{2}$ \\\hline
$\operatorname{Bl}_{p,q}\mathbb{P}^{2}$
&7&connected group of rank 2\\\hline
$\operatorname{Bl}_{p,q,1}$
&6&$\mathbb{C}_{m}${?}$ S_{2}\times S_{3}$
\end{tabular}

\medskip\noindent{\bf Section.}

\begin{definition}
\label{definition-moraga-2}
A class of varieties $\mathcal{C}$ is
\textit{\textbf{bounded}} if there exists a
projective morphism between finite type
schemes $\mathcal{X}\to T$ such that for
every $X\in \mathcal{C}$ we can find $t\in T$
such that $X_{t}\cong \mathcal{X}_{t}$.
\end{definition}

\begin{theorem}[Koll\'ar, Miyaoka-Mori, 92]
\label{theorem-moraga-2}
Fix $n \in \mathbb{Z}_{\geq 1}$. The class of
$n$-dimensional smooth Fano varieties is
bounded. $|-mK_{X}|$ is very ample, it is
controlled in terms of $n$.
\end{theorem}

\begin{quotation}
So it works in dimension 1, it works in
dimension 2, so it should work in dimension
3, right? This is actually not true. So let
me introduce the following definition:
\end{quotation}

\begin{definition}
\label{definition-moraga-3}
A variety $X$ is \textit{\textbf{rational}}
if $ X$ is birrational to $\mathbb{P}^{n}$
for some $n$.
\end{definition}

Griffiths and Clemens got to the right
answer, though their proof was wrong.

\begin{quotation}
This is actually valid in mathematics.
\end{quotation}

\begin{theorem}[Griffiths-Clemens, 72]
\label{theorem-moraga-3}
A smooth cubic 3-fold is not rational
\end{theorem}

\begin{definition}
\label{definition-moraga-4}
A normal projective variety $X$ is
\textit{\textbf{toric}} if there is a dense
open set of $X$ isomorphic to
$\mathbb{C}^{n}_{\mathbb{T}_{m}}$ and the
action of $\mathbb{G}^{n}_{m}$ on itself
extends to $X$.
\end{definition}

{*Explanation of the polytope*}

\begin{theorem}[Cox]
\label{theorem-moraga-4}
There is a bijection between

$$
\left\{ \substack{n\text{-dimensional}
\\\text{smooth projective} \\\text{Fano toric
varieties}  } \right\}/\cong   \qquad
\longleftrightarrow\qquad  \left\{
\substack{n\text{-dimensional}  \\
\text{reflexive smooth} \\\text{lattice
polytopes} }
\right\}/\operatorname{GL}(m,\mathbb{Z})
$$

where the quotient accounts for moving around
the polytope.
\end{theorem}

\medskip\noindent{\bf Conjecture
(Folklore?).} Any $n$-dimensional smooth
reflexive lattice polytope is inside
$[-1,1]^{n}$ up to translation and
$\operatorname{GL}(n,\mathbb{Z})$.

\begin{tabular}{cc}
$\mathbb{P}^{2}$ &right-angle triangle\\
$\mathbb{P}^{1} \times
\mathbb{P}^{1}$&square\\
$\operatorname{Bl}_{p}\mathbb{P}^{2}$ &right
angle triangle with truncated right angle\\
$\operatorname{Bl}_{p,q,r}\mathbb{P}^{2}$&
the third and last corner truncated

\end{tabular}
\begin{remark}
\label{remark-moraga-1}
In 1982 Batyrev classified smooth toric Fano
3-folds. There's 16. In 00's Krauser and
Skarke classified smooth toric Fano 4-foulds,
there's 124. In 15's Kraused and Nill
classified 5-folds. In 2007 \O bro wrote and
algorithm to classify smooth toric Fano
$n$-folds.

The number of smooth toric Fano $n$-folds
grows at least as $5^{n}$ assymptotically.
\end{remark}

\begin{quotation}
What is the bad part of smooth toric Fanos?
The following theorem:
\end{quotation}

\begin{theorem}[Cox]
\label{theorem-moraga-5}
A smooth toric Fano variety is rigid.
\end{theorem}

\begin{quotation}
On one side, smooth toric Fanos are nice
because they are combinatorial. On the other
side, not so much because they do not have
moduli.
\end{quotation}

\begin{tabular}{cccccc}
Smooth Fano
surface&Volume&Aut&Rational&Toric&dimension
of moduli\\\hline\hline
$\mathbb{P}^{2}$&9&
$\mathbb{P}\operatorname{GL}(3,\mathbb{C})$&
Yes&Yes&0\\\hline
$\mathbb{P}^{1} \times \mathbb{P}^{1}$,
$\operatorname{Bl}_{p}\mathbb{P}^{2}$&8&$\geq
\mathbb{G}^{2}_{m}$&Yes&Yes&0\\\hline
$\operatorname{Bl}_{p,q}\mathbb{P}^{2}$&7&$
\geq
\mathbb{G}^{2}_{m}$&Yes&Yes&0\\\hline
$\operatorname{Bl}_{p}\mathbb{P}^{2}$&6&$\geq
\mathbb{G}^{2}_{m}$&Yes&Yes&0\\\hline
$\operatorname{Bl}_{p}\mathbb{P}^{2}$&6&
finite&Yes&Yes&0\\\hline
\end{tabular}

\begin{theorem}[Ikovsky-Prokov, Mori-Mukai,
Shokurov, 90's]
\label{theorem-moraga-6}
(I-S) There are 105 families of smooth Fano
3-folds.

The general point of precisely 88 families
corresponds to a rational Fano variety.

16 "families" correspond to toric Fano
varieties.
\end{theorem}

\medskip\noindent{\bf Question.} What happens
in dimension 4? We are far from a
classification. \textit{Can we classify
smooth Fano 4-folds?}

No.

\begin{theorem}[Casagrande,24]
\label{theorem-moraga-7}
A smooth Fano 4-fold $X$ with $\rho(X)>12$ is
a product of surfaces. In particular
$\rho(X)\leq 18$.
\end{theorem}
\begin{center}
\begin{tabular}{|c|c|c|c|}
\hline
dimension&\# number of
families&rational&toric\\\hline
1&1&1&1\\\hline
2&10&10&5\\\hline
3&105&88&16\\\hline
4&??&??&124\\\hline
\end{tabular}\end{center}

\begin{quotation}
The aim of this minicourse is to introduce a
new notion of algebraic varieties that lies
between rational and toric. They are called
\textit{\textbf{cluster type varieties}}, and
study clusted Fano varieties.
\end{quotation}

\medskip\noindent{\bf Lecture 2.}

\begin{definition}
\label{definition-moraga-5}
A \textit{\textbf{log pair}} is a couple
$(X,\Delta)$ where $X$ is a normal
quasi-projective variety and $\Delta \geq 0$
is a $\mathbb{Q}$-divisor for which
$K_{X}+\Delta$ is $\mathbb{Q}$-Cartier.
\end{definition}

\begin{example}
\label{example-moraga-2}
$E$ elliptic, $i:E\to E$.
$\mathbb{Z}_{12}\hookrightarrow E$.
{\ldots{}} The conclusion is that studying
the pair $(\mathbb{P}^{1},D)$ is equivalent
to studying $\mathbb{P}^{1}$.
In-equivariantly.
\end{example}

\begin{example}
\label{example-moraga-3}
$S\subset X$, $(K_{X}+S)|_{S}\sim K_{S}$.
Understanding the log pair $(X,S)$ is useful
to understand $S$.
\end{example}

\begin{definition}
\label{definition-moraga-6}
Let $(X,\Delta )$ be a log pair and $\pi:Y\to
X$ a projective birational morphism from a
normal variety $E\subset Y$ prime divisor.
For

$$
\pi^{*}(K_{X}+\Delta )=K_{Y}+\Delta_{Y}
$$

where $\Delta_{Y}$ may not be effective, the
\textit{\textbf{log discrepancy}} of
$(X,\Delta )$ at $E$ is

$$
\alpha_{E}(X,\Delta
):=1-\operatorname{coeff}(\Delta_{Y}).
$$

\end{definition}

\begin{remark}
\label{remark-moraga-2}
The number $\alpha_{E}(X,\Delta )$ gets more
negative as the multiplicity of $(X,\Delta )$
along the tangent directions corresponding to
$E$ "gets higher".
\end{remark}

\begin{example}
\label{example-moraga-4}
\begin{enumerate}
\item $X=\mathbb{A}^{n}$, $\Delta =0$,
$\pi:Y\supset E\cong \mathbb{P}^{1}\to
\mathbb{A}^{n}$ the blow-up at
$(0,\ldots,0)$.
$pi^{*}(K_{\mathbb{A}^{n}}=K_{\mathbb{P}^{n}
+(1-n)}E$, $\alpha_{E}(\mathbb{A}^{n} =n$.
\item $ X=\mathbb{A}^{n}$, $\Delta
=\lambda_{1}H_{1}+\ldots+\lambda_{n}H$,
$\pi^{*}(K_{\mathbb{A}^{n}
+\Delta}=K_{Y}+(1-n)E+\sum_{i=1}^{n}
\lambda_{i}E+\sum_{i=1}^{n}
\lambda_{i}\tilde{H}_{i}$,
$\alpha_{E}(\mathbb{A}^{n} +\Delta
)=n-\sum_{i}\lambda_{i}$.
\end{enumerate}
\end{example}

\begin{definition}
\label{definition-moraga-7}
A log pair $(X,\Delta)$ has
\textit{\textbf{log terminal singularities}}
if all its log discrepancies are $>0$.

A log pair $(X,\Delta )$ is
\textit{\textbf{log canonical}} if all its
log discrepancies are $\geq 0$.

A log pair $(X,\Delta )$ is
\textit{\textbf{terminal}} if all its log
discrepancies from exceptional divisors are
$>1$.

A log pair $(X,\Delta)$ is
\textit{\textbf{canonical}} if all its log
discrepancies are $\geq 1$.
\end{definition}

\begin{exercise}
\label{exercise-moraga-1}
$(X,\Delta )$ is log terminal iff
$\lambda_{i}<1$ for all $i$.

$(X,\Delta )$ is log canonical iff
$\lambda_{i}\leq 1$ for all $i$.
\end{exercise}

\begin{remark}[Historical note]
\label{remark-moraga-3}
$X$ smooth pair, $K_{X}$ is big.
$X\dashrightarrow X'$. $K_{X'}$ big and nef.
{Terminal singularities are the singularities
appearing in the terminal (minimal) model.}

$$
\xymatrix{X\ar@{-->}[r]\ar@{-->}[dr]&X'
\ar[d]\\
&X''}
$$

{Canonical singularities are the
singularities appearing in the canonical
model.}
\end{remark}

\begin{example}
\label{example-moraga-5}
$G\hookrightarrow (\mathbb{A}^{n},0)$,
$X=\mathbb{A}^{n} /G$. $ X$ has log terminal
singularity.

$$
\xymatrix{\tilde{E}\subset\tilde{Y}
\ar[r]^{/G}\ar[d]&Y\supseteq
E\ar[d]\\ \mathbb{A}^{n}\ar[r]^{/G}&X}
$$

$\alpha_{E}(X)=\alpha\overset{\sim}{E}(
\mathbb{A}^{n}
)/r>0$. And $r$ is the ramification index at
$E$.
\end{example}

\begin{proposition}
\label{proposition-moraga-1}
Let $X$ be a smooth projective variety with
$\rho(X)=1$, $H$ an ample divisor on $X$,
$v\in C_{X}:=\operatorname{Spec}\left(
\bigoplus_{n\geq 0}
H^{0}(X,\mathcal{O}_{X}(mH))  \right) $,
\begin{itemize}
\item $C_{X}$ is log terminal iff $X$ is
Fano.
\item $C_{X}$ is log canonical iff $X$ is CY.
\item $C_{X}$ is not lc iff $X$ is
canonically polarized.
\end{itemize}
\end{proposition}

\begin{exercise}
\label{exercise-moraga-2}
$H\subseteq \mathbb{A}^{n}$ if
$\operatorname{mult}_{0}H>n+1$, then $(H,0)$
is not log canonical.
\end{exercise}

\medskip\noindent{\bf Dimension 1.}
{is not very intersting:} locally
analytically this is the only example:
$(\mathbb{A}^{1},\lambda \{0\} )$, log
terminal iff $\lambda<1$, log canonical iff
$\lambda\leq 1$.

\medskip\noindent{\bf Dimension 2.}

$(X,x)$ is terminal iff smooth.

$(X,x)$ is canonical iff ADE sing.

$(X,x)$ is log terminal iff quotient sing.

$(X,x)$ is strictly lc iff elliptic sing.

\begin{theorem}[Reid, 80's]
\label{theorem-moraga-8}
A terminal 3-fold singularity $(X,x)$ is a
hyperquotient singularity.
\end{theorem}

\begin{definition}
\label{definition-moraga-8}
Let $(X,x)$ be a the germ of a minimal
singularity over $\mathbb{C}$. $\varepsilon
>0$ small and
$\operatorname{Link}(X;x)=S_{\varepsilon
}^{(x)} \cap X$, $S_{\varepsilon }$ is sphere
of radius $\varepsilon $ around $x$ in
$\mathbb{A}^{N}$.

The \textit{\textbf{local fundamental group}}
of  $(X;x)$ is
$\pi_{1}(\operatorname{Link}(X;x)) $.
\end{definition}

\begin{theorem}[Koll\'ar, 2010]
\label{theorem-moraga-9}
There is a sequence $(X_{r};0)$ of terminal
4-fold singularities with
$\operatorname{emdim}(X_ {r})\to \infty$.
$\pi_{1}(X_{r};0) =\{1\} $ and
$\operatorname{Cl}(X_{r},0)=0$
\end{theorem}

{*Some table classifing dimensions 2,3,4,
Terminal-smooth, hyperquotient,?, Canonical,
Ade-?,?,?, log termi-quotient,?,?, log
cas-Elliptic singularity finetely many
families,?,?}

\begin{definition}
\label{definition-moraga-9}
A singularity $(X;x)$ is of
\textit{\textbf{log terminal type}} if there
exists $\Delta \geq 0$ such that
$(X,\Delta;x)$ is log terminal.
\end{definition}

\begin{remark}
\label{remark-moraga-4}
The $\Delta$ can be used as a correction term
for the lack of $\mathbb{Q}$-Gorensteiness.
\end{remark}

\begin{example}
\label{example-moraga-6}
$\sigma=\left\langle{}
e_{1},e_{2},e_{3},2e_{1}+2e_{2}-e_{3}
\right\rangle{} $. $X(\sigma)$ is not
$\mathbb{Q}$-Gorenstein.
$(X(\sigma),\frac{2}{3}D_{4}$ is log
terminal.
\end{example}

\begin{theorem}
\label{theorem-moraga-10}
A toric singularity is log terminal type.
\end{theorem}

\begin{theorem}
\label{theorem-moraga-11}
A singularity $(X;x)$ is toric if and only if
it is a torus quotient singularity.
\end{theorem}

\begin{theorem}[Biaun-Greb-Langlass-M, 22]
\label{theorem-moraga-12}
Log terminal type singularities are preserved
by reductive quotient
\end{theorem}

\begin{theorem}[Jordan, 1870's]
\label{theorem-moraga-13}
There is a constant $c(n)$ satisfying the
following. Let $G\leq
\operatorname{GL}(n,\mathbb{C})$ be a finite
subgroup. Then $G$ admits a normal abelian
subgroup if index $\leq c(n)$.
\end{theorem}

\begin{theorem}[Collins, 2010]
\label{theorem-moraga-14}
$c(n)=(n+1)!$ provided $n\geq 71$
\end{theorem}

\begin{corollary}
\label{corollary-moraga-1}
Let $(X;x)$ be an $n$-dimensional quotient
singularity. Then
$\pi_{1}^{\operatorname{loc}}(X;X)$ admits a
normal abelian group index $\leq c(n)$.
\end{corollary}

\begin{theorem}[Braun-Filipazzi-M-Svaldi, 20]
\label{theorem-moraga-15}
There is a constant $K(n)$ only depending on
$n$ and satisfying the following. Let $(X;x)$
be a log terminal singularity of dimension
$n$. Then $\pi_{1}^{\operatorname{loc}}(X;x)$
admits a normal finite abelian subgroup of
rank $\leq n$ and $\operatorname{index}\leq
K(n)$.
\end{theorem}

\begin{definition}
\label{definition-moraga-10}
Let $(X,\Delta)$ be a log pair. The
\textit{\textbf{regularity}} of $(X,\Delta)$
is

The \textit{\textbf{absolute regularity}} of
$(X,\Delta)$ is

$$
\hat{\operatorname{reg}}(X,\Delta):=
\operatorname{max}\left\{
\operatorname{reg}(X,\Delta +B|(X,\Delta +B),
B\geq 0, \text{ is lc}  \right\}
$$

\end{definition}

\medskip\noindent{\bf Editorial note.} In the
original notes the formula defining the
regularity of $(X,\Delta)$ is commented out
and the definition is left unfinished; it is
missing here.

\begin{theorem}[M (speaker),21]
\label{theorem-moraga-16}
Let $(X;x)$ be a $n$-dimensional log terminal
singularity of absolute regularity $r\in
\{-1,\ldots,n-1\} $. Then there exists an
analytic embedding $(\mathbb{D}^{*} )^{r+1}
\hookrightarrow X^{\operatorname{sm}}$ such
that the image of the induced homomorphism

$$
\pi_{1}(\left(\mathbb{D}^{*}
\right)^{r+1})\to
\pi_{1}^{\operatorname{loc}}(X;x)
$$

is finite, normal and of cokernel of order
$\leq K(n)$.
\end{theorem}

\medskip\noindent{\bf Summary.}

\begin{enumerate}
\item Definition of log terminal.
\item Log terminal sing cannot be classified.
\item Reductive quotient singularities are
log terminal type.
\item Reductive quotient sings are log
terminal type.
\item  The fundamental groups of log
terminals sing is controlled by {heav} toric
geometry.
\end{enumerate}

\medskip\noindent{\bf Lecture 3.}

\medskip\noindent{\bf Singular Fano
varieties.}

\medskip\noindent{\bf Why singular Fano
varieties?}

\begin{enumerate}
\item Aloows to understand symmetries of
smooth Fano varieties.
\item Singular Fano varieties naturally
appear when compactigyinf moduli spaces og
smooth Fano varieties.
\item Smooth CY manifolds often admit
degeneratoions to recudible varities where
the central fiber has singular Fanos as
irreducible components.
\end{enumerate}

\begin{definition}
\label{definition-moraga-11}
A \textit{\textbf{Fano variety}} is a normal
projective variety $X$ with log terminal
singularities and $-K_{X}$ ample. A
\textit{\textbf{log Fano}} is a pair
$(X,\Delta)$ with $-(X_{X}+\Delta)$ ample and
$(X,\Delta)$ log terminal. A variety $X$ is
\textit{\textbf{Fano type}} if $(X,\Delta)$
is log terminal and $(-K_{X},\Delta)$ is nef
and big for some suitable $\Delta \geq 0$. A
pair $(X,\Delta)$ is \textit{\textbf{log CY}}
if $(X,\Delta)$ is lc and $K_{X}+\Delta$.
\end{definition}

\begin{exercise}
\label{exercise-moraga-3}
$\mathbb{P}^{1} \times \mathbb{P}^{1}$ smooth
Fano, blow-up all torus inv points $X\to
\mathbb{P}^{1} \times \mathbb{P}^{1}$, then
$X$ is not Fano.
\end{exercise}

\begin{remark}
\label{remark-moraga-5}
Any projective normal toric variety is Fano
type.
\end{remark}

\begin{theorem}[Prokhorov, X 14]
\label{theorem-moraga-17}
Let $(X,x)$ be a log terminal sing. Then
there exists a projective birational morphism
$\pi:Y\to  X$ extracting a unique normal
divisor $E$ with $\pi(E)=x$, $X\setminus
\{x\} \cong Y\setminus E$, and $E$ is a Fano
type.
\end{theorem}

$$
\xymatrix{\left\{ \substack{\text{Fano
type}\\ \text{varieties of dim } n}
\right\}\ar@/^1pc/[r]^{\text{cones}}&\left\{
\substack{\text{log terminal}\\ \text{type
sing of dim } n+1}
\right\}\ar@/^1pc/[l]^{\text{plt blow-ups}}}
$$

and now another answer to the initial
question of this lecture:

\begin{enumerate}
\item[4.] Singular Fanos allow us to
understand log terminal singularities (of one
dimension more).
\end{enumerate}

\medskip\noindent{\bf What we know
classificationwise.}

\begin{itemize}
\item Canonical Fano surfacs of $\rho(x)=1$
are classified. 28 families.
\item Terminal Fano 3-folds, not classified
but a lot of explicit work.
\item Canonical toric Fano varities are
classified up to dimension 7 (at least six
months ago) (Nill,\ldots{}) These grow double
exponentially with dimension.
\end{itemize}

\medskip\noindent{\bf Regularity of a pair
$(X,\Delta)$.}

\begin{quotation}
Let's do some qualitative mathematics. So we
aim to describe the properties of these
varieties.
\end{quotation}

Let $(X,\Delta)$ be a log pair. Recall the
definitions of regularity and absolute
regularity.

\begin{remark}
\label{remark-moraga-6}
The \textit{\textbf{absolute regularity}} of
a toric singularity $= \dim -1$.
\end{remark}

\begin{definition}
\label{definition-moraga-12}
The \textit{\textbf{absolute coregularity}}
of  $(X,\Delta ;x)$ is

$$
\hat{\operatorname{coreg}}(X,\Delta,x)=\dim
X-\hat{\operatorname{reg}}(X,\Delta;x)
$$

\end{definition}

By definition,
$\hat{\operatorname{coreg}}(X,\Delta ;x)\in
\{0,\ldots,\dim X\} $.

\begin{definition}
\label{definition-moraga-13}
Let $X$ be a Fano type variety. The
\textit{\textbf{absolute regularity}} of $X$
is

$$
\hat{\operatorname{reg}}(X):=\max
\{\operatorname{reg}(X,B):(X,B)\text{ is CY}
\}
$$

$$
\hat{\operatorname{coreg}}(X):=\dim
X-\hat{\operatorname{reg}}(X)-1
$$

\end{definition}

\begin{theorem}[M,21]
\label{theorem-moraga-18}
Let $X$ be a Fano type variety of dimension
$n$ and absolute regularity $r$. Then there
is a local analytic embedding

$$
(\mathbb{D}^{*} )^{r+1} \hookrightarrow
X^{\operatorname{sm}}
$$

such that the image of the homomorphism on
fundamental groups $\mathbb{Z}^{r+1} \to
\pi_{1}(X^{\operatorname{em}}) $ is finite,
normal and of index $\leq k(n)$.
\end{theorem}

\begin{corollary}
\label{corollary-moraga-2}
Fano varieties with regularity $-1$ and
dimension $n$ have size of the fundamental
group of the smooth loci bounded by $k(n)$,
that is,  $|\pi_{1}(X^{*m}) |\leq k(n)$.
\end{corollary}

\begin{remark}
\label{remark-moraga-7}
The expectation is that $k(n)$ is
$\mathcal{O}((in)!)$.
\end{remark}

\begin{definition}
\label{definition-moraga-14}
An \textit{\textbf{exceptional Fano variety}}
is a Fano variety $X$ for which
$\hat{\operatorname{coreg}}(X)=\dim X$. Iff
$\alpha(X)>1$ (for the
birrationally-inclined).
\end{definition}

\medskip\noindent{\bf Boundedness of Fano
varieties.}

\medskip\noindent{\bf Question.} Are all the
Fano varieties of the same dimension bounded?

\medskip\noindent{\bf Answer.}
No. $\mathbb{P}(1,1,n)$, $n\to \infty$.

\begin{definition}
\label{definition-moraga-15}
A Fano variety $X$ is
\textit{\textbf{$\varepsilon $-log terminal}}
if all its log discrepancies are
$>\varepsilon $.

$$
0\text{-log termial} \iff\text{log terminal}
$$

$$
1\text{-log terminal} \iff\text{terminal}
$$

\end{definition}

\begin{theorem}[Birkari, 2016, "singularities
of linear systems and boundedness", hardcore
and technical, there are surveys]
\label{theorem-moraga-19}
Fix $n$ a positive integer and $\varepsilon
>0$. Then the class of $n$-dimensional
$\varepsilon $-log terminal varieties is
bounded.

$\varepsilon -H$ in toric geometry
corresponds to controlling angles in the
polytope away from $\varepsilon $.
\end{theorem}

\begin{theorem}[Birkari, 2016, used by the
one above, "Anti-pluricanonical linear
systems", hardcore and technical, there are
surveys]
\label{theorem-moraga-20}
Fix $n$ a positive integer. Then the class of
$n$-dimensional exceptional Fano varieties is
bounded.
\end{theorem}

\begin{quotation}
We shall see come recent developments about
this.
\end{quotation}

\begin{definition}
\label{definition-moraga-16}
Let $X$ be a normal projective variety. A
\textit{\textbf{$\mathbb{Q}$ement}} is a
divisor $0\leq B \sim_{\mathbb{Q}}-K_{X}  $
for which $(X,B)$ is log CY.

A $\mathbb{Q}$-complement $B$ is said to be a
\textit{\textbf{$N$-complement}} if
$N(K_{x}+B)\sim 0$.
\end{definition}

\begin{theorem}[Birkari, 2016, the real
breakthrough, after this we posted the one on
the top]
\label{theorem-moraga-21}
Fix $n \in \mathbb{Z}_{>0}$, there exists
$N(n)$ satisfying the following. Any
$n$-dimension Fano variety admits a
$N(n)$-complement.
\end{theorem}

\begin{theorem}[Totaro, 22]
\label{theorem-moraga-22}
$N(n)$ grows at least doubly exponentially.
There is a Fano 4-fold (exceptional) $X$ with
$|-mK_{X}|=\varnothing $ for all $m\leq
4.677.233$.
\end{theorem}

\begin{quotation}
To prove this essentially you need to give
examples.
\end{quotation}

\begin{remark}
\label{remark-moraga-8}
All toric varieties admit 1-complements of
coreg=0.
\end{remark}

\begin{theorem}[Figueroa-Filipazzi-M-Peng,22]
\label{theorem-moraga-23}
A Fano variety of absolute corregularity zero
amits either a complement or a 2-complement
of corregularity zero. In particular,
$|-2K_{X}| \neq \varnothing $.
\end{theorem}

\begin{quotation}
This is a generalization that the toric
boundary exists for toric variety.
\end{quotation}

$$
\xymatrix{&\{n\text{-dimensional Fano
variety}\}&\\ \left\{
\substack{n\text{-dimensional rational}\\
\text{Fano varieties}}
\right\}\ar@{^{(}->}[ur]\ar@{^{(}->}[rr]^{?
\text{(not
sure)}}&&\left\{ \substack{n\text{-dim
coregularity}\\ \text{zero Fano varieties}}
\right\}\\ &\left\{ \substack{n\text{-dim
toric}\\ \text{Fano varieties}}
\right\}\ar@{^{(}->}[ul]\ar@{^{(}->}[ur]&}
$$

\begin{theorem}[Kaloghiro, 10]
\label{theorem-moraga-24}
There exists a terminal Fano 3-fold which
birationally superrigid that has absolut
coreg =0.
\end{theorem}

\medskip\noindent{\bf Editorial note.} The
notes cover Lectures 1, 2 and 3 only; the
last item of the plan (cluster Fano
varieties) is not in the notes. The tables
were retyped as simple tabulars: two cells
containing line breaks were joined on one
line. A commented-out table in the first
lecture is omitted. The diagrams drawn with
tikz-cd in the original are written as
xymatrix diagrams.

\section{Algebraic tori in the complement of
quartic surfaces}
\label{section-alves-algebraic-tori-quartic}
%% speaker: Eduardo Alves da Silva, Institut
%% speaker: de Mathématiques d'Orsay
%% place: Seminário das Sextas, PUC-Rio
%% date: January 22

\noindent
Eduardo Alves da Silva, Institut de
Math\'ematiques d'Orsay.
Semin\'ario das Sextas, PUC-Rio.
January 22 (year not recorded in the notes).

\medskip\noindent
No abstract available.

\begin{enumerate}
\item Recall from Eduardo's last talk that

\begin{definition}
\label{definition-alves-tori-1}
A \textit{\textbf{log Calabi-Yau}} pair is a
lc pair  $(X,D)$ consisting of a normal
projective variety $X$ and a reduced Weil
divisor $D$ such that
$K_X+D\sim_{\mathbb{Z}}0$.
\end{definition}

Now we have given a slightly different
definition.

\item  It was possible to classify Fano
threefolds using these methods. (Recent.)
Why? The anticanonical bundle is ample so
there is a power that gives a section\ldots{}
\item
\begin{theorem}[Gross-Hacking-Keel, 2017]
\label{theorem-alves-tori-1}
Log CY pair of index 1 and of corregularity 0
is log rational.
\end{theorem}
\medskip\noindent{\bf Conjecture (Shokirov).}
Every threefold rational log CY pair $(X,B )$
and corregularity 0 is log rational.

\item  Perhaps this has to do with a
classification of toric varieties:
\begin{theorem}[Brown, McKerman, Svaldi,
Zong, 2018]
\label{theorem-alves-tori-2}
$(X,B)$ log CY pair implies $c (X,B ) \geq
0$.
\end{theorem}
No I think it's this one:
\medskip\noindent{\bf Idea.}
If $c(X,B)<1$, then $X,\left\lfloor B
\right\rfloor$ is toric.

\item Toric $\implies $ cluster type
$\implies $ log rational. Counter examples
for both of the reverse arrows are known.
\item
\begin{theorem}[-,Figueroa, Moraga, 2024]
\label{theorem-alves-tori-3}
$(\mathbb{P}^3,B)$ log CY pair, $i=1$
corregularity 0. Assume Bb is non-normal.
$(\mathbb{P}^3, B )$ is cluster type unless
the nodal locus of  $B$ is centered in a
plane.
\end{theorem}
\end{enumerate}

\medskip\noindent{\bf Editorial note.} In the
original the speaker's surname is typed
``Alvez''; I use ``Alves'', as in the other
section on this speaker's talk. The item
beginning ``No I think it's this one'' is
Daniel's own comment.

\end{document}
